%! TEX root = ../main.tex

\graphicspath{{../Paper_IV/Paper_IV_AOS_latex/}{../Paper_III/Paper_III_latex_arxiv/}}

\chapter[Wicksell's problem: the Bayesian nonparametric approach]{\huge Wicksell's problem: the Bayesian nonparametric approach}
\chaptermark{Bayesian approach to Wicksell's problem}

\section[Semiparametric BvM via Isotonized Posterior]{Semiparametric Bernstein-von Mises phenomenon via Isotonized Posterior}
\sectionmark{Semiparametric BvM via Isotonized Posterior}

\subsection{Introduction}

Before diving into the results obtained using Bayesian nonparametric modelling approaches, we recall to the reader the main formulas that are useful to understand fully the following results. We start by formally recalling the main mathematical objects of interest in the material science setting and in the astronomical setting, respectively.

The setting of the opaque material having globular microstructure is formalized by assuming that the embedded particle squared radii are i.i.d.\ realizations from a random variable $X^* \sim F^*_0$ and our observations consist of cross sections of the medium (i.e.\ "planes cutting the medium"). Note that the spheres with bigger radii are more likely to be sampled. Thus the distribution of the squared radii of the spheres actually cut is not $F^*_0$ but a biased version of it; we write for the squared radius of the sphere actually cut: $X \sim F_0$. Since we assume that the samples are gathered uniformly, by the Pythagorean Theorem, the observable squared circle radius $Z$ is related to $X$ via: $ Z = (1-U^2) X = Y X,$ where $U \sim U[0,1]$ so $Y \sim \operatorname{Be}(1,1/2)$, with density $y \mapsto \frac{1}{2} (1-y)^{-1/2} \ind_{(0,1)}(y)$. So, the observed squared radii $Z$ are a random fraction of the true squared radii depending on where the plane cuts through the sphere. Consequently, since both $X$ and $Y$ are nonnegative, a version of the density $g_0$ of $Z$ can, for $z\geq 0$, be written as:
\begin{align}\label{eq: g_F_0 expression}
	g_0 (z) = \frac{d}{dz} \int_{(0,\infty)} \pb(Y \leq z/x ) \, dF_0(x) =\int_{z}^{\infty} \frac{dF_0(x)}{2\sqrt{x^2 - xz}}
\end{align}
The problem in the astronomical setting can be modelled as follows: let $\mathbf{Z}= (Z_{p_1}, Z_{p_2},Z_{p_3})$ be a spherically symmetric random vector (the three-dimensional position of a star in a galaxy) of which only $(Z_{p_1}, Z_{p_2})$ can be observed. Under the assumption that the joint of $\mathbf{Z}$ is a function of the distance to the center of the galaxy, astronomers are interested in estimating the distribution of the squared distance of a star to the center of the globular cluster or galaxy $X = Z^2_{p_1}+ Z^2_{p_2}+Z^2_{p_3} \sim F_0$ using $Z = Z^2_{p_1}+ Z^2_{p_2}$. From the assumed spherical symmetry, $\mathbf{Z}
\sim O \mathbf{Z}$ for all orthogonal matrices $O \in \mathrm{O}(3)$. Therefore $\mathbf{U} := \mathbf{Z} / \|\mathbf{Z}\|$ (note $X = \|\mathbf{Z}\|^2$) is uniform on the unit sphere $\mathbb{S}^2$. It follows that the random variable $Z_{p_3}$ (the third coordinate) is distributed as: $\sqrt{X} \cdot U_3$ where $U_3 \sim U[-1,1]$ independent of $X$. Note also that $Y \equiv 1-U_3^2 \sim \operatorname{Be}(1,1/2)$ as above. Because:
\begin{align*}
	\pb (Z \leq z) = \pb ( Z^2_{p_1}+ Z^2_{p_2}+Z^2_{p_3} \leq z + Z^2_{p_3}) = \pb (X \leq z + X \cdot U_3^2) = \pb (YX \leq z),
\end{align*}
it follows that both for the application in materials science and the one in astronomy, the density of the observations $Z \sim G_0$ is the one given on the right-hand side of \eqref{eq: g_F_0 expression}.

By using Abel-integral equation's theory, we can invert \eqref{eq: g_F_0 expression} and obtain (c.f. \cite{19}):
\begin{align}\label{eq: solution range equation}
  &F_0(x) = 1 -\frac{2}{\pi} \sqrt{x} \, V_0(x) - \frac{2}{\pi} \int_{x}^{\infty} \frac{V_0(s)}{2 \sqrt{s}} \, ds,  
\end{align}
where:
\begin{align}\label{eq: V_0 chapter3}
    V_0(x) \vcentcolon=\int_x^{\infty} \frac{g_0(z)}{\sqrt{z-x}} \: d z.
\end{align}
(The proof of \eqref{eq: solution range equation} is provided in Lemma \ref{lemma: solution range equation} in the Appendix). In Corollary \ref{thm: BvM for F}, we formally show that the functional $x \mapsto - \frac{2}{\pi} \int_{x}^{\infty} V_0(s) / (2 \sqrt{s}) \, ds$ can be estimated at the rate $\sqrt{n}$, whereas the functional $V_0(x)$ can be estimated (nonparametrically) only at the slower rate $\sqrt{n}/\sqrt{\log{n}}$. This can be intuitively seen because, by plugging \eqref{eq: V_0 chapter3} into the second term of \eqref{eq: solution range equation} and using Fubini's theorem to exchange order of integration:
\begin{align}\label{eq: second equation range}
\int_{x}^{\infty} \frac{V_0(s)}{2 \sqrt{s}} \, ds =\frac{\pi}{2}-\int_0^{\infty} \sin ^{-1} \sqrt{1 \wedge \frac{x}{z}} \,  g_0(z)\, dz 
\end{align}
and $\mathrm{Var}[\sin ^{-1} \sqrt{1 \wedge\left(x / Z\right)}] < \infty$, for $Z \sim G_0$, in contrast with $\mathrm{Var}[(Z-x)^{-1/2}] = \infty$. It thus follows that the estimation of $V_0$ and $F_0$ are essentially equivalent. Furthermore, as shown in \cite{1}, the relation
$ 
F_0^*(x) = 1 - \frac{V_0(x)}{V_0(0)}
$ 
holds, which implies that also in the materials science setting, the primary target of estimation is $V_0$. Accordingly, our focus in this chapter will first be on estimating $V_0(x)$, and we will subsequently show that these results naturally extend to the estimation of $F_0(x)$. The latter will be the main object of interest in the astronomical application discussed in the simulation study.

We aim to develop an easy-to-compute and statistically efficient nonparametric Bayesian method for estimating \( V_0(x) \) at a given point \( x \geq 0 \). We begin by placing a Dirichlet Process prior directly on the distribution function of the observables: \( G \sim \operatorname{DP}(\alpha) \). Given an observed sample \( Z_1, \dots, Z_n \) from cdf \( G_0 \) (associated with density $g_0$), the posterior distribution is (c.f.\ \cite{59}, Theorem 4.6 in \cite{4}):
\begin{align}\label{eq: dirichlet posterior chapter3}
	G \mid Z_1,\ldots,Z_n \sim \text{DP}(\alpha + n \mathbb{G}_n), \quad \quad  \mathbb{G}_n := \frac{1}{n} \sum_{i=1}^n \delta_{Z_i}.
\end{align}
This choice offers computational simplicity due to the conjugacy of the posterior. Leveraging the inverse relations inherent to the problem, we can estimate the functional of interest using \eqref{eq: solution range equation} by
\begin{align}\label{eq: naive bayes chapter3}
	V_{\scriptscriptstyle{G}}(x) := \int_x^\infty \frac{dG(z)}{\sqrt{z - x}}.
\end{align}
This initial step yields the \textit{Naive Bayes Posterior} (NBP). We show below that this satisfies a BvM phenomenon with asymptotic variance \( g_0(x) \). However, the realisations \( V_{\scriptscriptstyle{G}} \) of this plug-in posterior are not necessarily monotone functions. To remedy this, we project them into the \( \mathbb{L}_2 \) subspace of decreasing and right-continuous functions. It turns out that this not only corrects the non-monotonicity, but also improves the estimator. We show that the resulting \textit{Isotonized Inverse Posterior} (IIP) $\hat{V}_{\scriptscriptstyle{G}}(x)$ satisfies a BvM theorem with improved variance that depends on the local smoothness of $V_0$. The smoothness at $x$ is measured through a parameter $\gamma > \frac{1}{2}$ that can be loosely understood as the degree of H\"{o}lder continuity at $x$. Using \eqref{eq: solution range equation}, we can define the IIP to target the estimation of $F_0(x)$, which we will denote by $\hat{F}_{\scriptscriptstyle{G}}(x)$. The formal definitions of $\hat{V}_{\scriptscriptstyle{G}}(x)$ and $\hat{F}_{\scriptscriptstyle{G}}(x)$, being quite technical, are given in \eqref{eq: IIP} and \eqref{eq: IIP F0} below, respectively. 

Thus the approach in this chapter has three advantages:
\begin{itemize}
	\item[(i)] \textbf{Computational speed}: The Dirichlet posterior \eqref{eq: dirichlet posterior chapter3} is explicit and its projection can be effectively computed by the Pool-Adjacent-Violators Algorithm (PAVA, c.f.\ \cite{61}).  
	\item[(ii)]  \textbf{Efficiency and adaptation}: The IIP achieves the asymptotic variance of the efficient estimators introduced in \cite{27} and \cite{10} and it is adaptive to the level of local smoothness of $V_0$.
	\item[(iii)] \textbf{Uncertainty quantification}: The Bayesian framework inherently provides uncertainty quantification without the need to estimate \( \gamma \) or other parameters. We show below that this is asymptotically correct in the frequentist sense.
\end{itemize}

The Naive Bayes Posterior (NBP) is analogous to the naive estimator introduced in \cite{1} and \cite{15} which is constructed as follows: given a sample $Z_1, \ldots, Z_n$ from the density $g_0$ and for $\mathbb{G}_n$ the empirical distribution function based on this sample,
\begin{align}\label{eq: naive estimator of V chapter3}
    V_n(x) \vcentcolon=\int_x^{\infty} \frac{d\mathbb{G}_n(z)}{\sqrt{z-x}}.
\end{align}
In \cite{27}, the authors showed that the projection of the naive estimator onto the set of monotone functions, the \textit{Isotonized Inverse Estimator} (IIE) $\hat{V}_n$, is asymptotically efficient and the attained asymptotic variance depends on the local smoothness of $V_0$ in the same way as for the IIP.

Projecting the naive posterior distribution departs from a classical Bayesian framework which might assign a Dirichlet Process (DP) prior to \( F \) and employ a hierarchical model to sample from the density of the observations. A Bernstein-von Mises (BvM) result for this approach is currently unavailable. Moreover, the available algorithms for posterior sampling rely on Gibbs sampling procedures, which are much slower than the method we propose here. 

From a practical perspective, the Bernstein-von Mises (BvM) result established in this chapter assures practitioners that the Bayesian methodology based on DP prior provides all the efficiency benefits of the Isotonic Inverse Estimator (IIE) of \cite{27}, but without requiring the estimation of difficult-to-access nuisance quantities such as the density of the observations $g_0(x)$ or the local H\"{o}lder smoothness $\gamma > 1/2$. In applications, where one aims to recover the distribution of squared distances of stars from their projected two-dimensional positions or the distribution of the embedded particles' squared radii, this is especially valuable: the underlying smoothness of these distributions (and thus of $V_0$) may vary across radii, and the BvM result holds for all local smoothness levels $\gamma > 1/2$. This significantly improves over existing bootstrap-based uncertainty guarantees of \cite{19}, which only hold under the stronger assumption that $V_0$ is differentiable at the estimation point $x$ (which is even stronger than $\gamma = 1$). 

\subsection{Main results}\label{sec: main results chapter3}

We start by proving a BvM result for the \textit{Naive Bayes Posterior} (NBP), analogous to the result obtained for the naive estimator (c.f.\ Theorem 2 in \cite{1}). For our first result, we need our version of $g_0$ to be well-defined and bounded, i.e.:
\begin{align*}
    \| g_0 \|_{\infty} := \sup_{z \geq 0} \bigg| \int_{z}^{\infty} \frac{dV_0(s)}{\pi \sqrt{s - z}} \bigg| < \infty. \numberthis \label{eq: condition finite g at x chapter3}
\end{align*}
Assumption \eqref{eq: condition finite g at x chapter3} asserts that projected 2D star positions do not become arbitrarily concentrated. From a mathematical perspective, without this assumption the inverse problem becomes highly unstable: small errors in the data can cause arbitrarily large errors in estimating $V_0$. For instance, if $g_0$ is unbounded at $x$, this leads to instability in estimating $V_0(x)$, reflected in an infinite asymptotic variance. Thus, boundedness of $g_0$ ensures stability of the inverse mapping. Mathematically, if $F_0$ is H\"{o}lder continuous with degree $\gamma > 1/2$, then $g_0$ is bounded (cf.\ Lemma A.3 in \cite{27}).

In equation \eqref{eq: claim naive} and \eqref{eq: claim} below we use the notion of conditional weak convergence in distribution (in probability). For its precise definition we refer the reader to \eqref{eq: formal definition conditional convergence chapter3}-\eqref{eq: formal definition conditional convergence 1 chapter3} in Appendix A.
\begin{theorem}[BvM for NBP]\label{thm: BvM naive}
	Let $x \geq 0$ and suppose that \eqref{eq: condition finite g at x chapter3} holds true. Let $G \sim \operatorname{DP}(\alpha)$, with $\int_{x}^{\infty} (s-x)^{-1/2} \, d \alpha(s) < \infty$. Then, in probability under $G_0$: 
	\begin{align}\label{eq: claim naive}
	\frac{\sqrt{n}}{\sqrt{\log n}}\big( V_{\scriptscriptstyle{G}}(x)-V_n(x)\big) \mid Z_1, \ldots, Z_n \rightsquigarrow N\big(0, g_0(x)\big).
	\end{align}
\end{theorem}
\begin{proof}
	For a measure $P$ and a function $f$, we denote by $Pf$ the integral $\int f \, dP$, this holding also for empirical measures like $\mathbb{G}_n = \frac{1}{n} \sum_{i=1}^n \delta_{Z_i}$. 
	
	Since $G \mid Z_1,\ldots,Z_n \sim \text{DP}(\alpha + n \mathbb{G}_n)$ by proposition G.10 in \cite{4}, for \( Q \sim \text{DP}(\alpha) \), \( \mathbb{B}_n \sim \text{DP}(n \mathbb{G}_n) \), \( V_n \sim \text{Be}(|\alpha|, n) \) and $f(z) = 1/\sqrt{(z-x)_+} $, $x \geq 0 $, conditionally on $Z_1,\ldots,Z_n$, we have the following representation in distribution:
	\begin{align}\label{eq: first representation in distribution}
	G f = V_n Q f + (1 - V_n) \mathbb{B}_n f,
	\end{align}
	where \( Q  \), \( \mathbb{B}_n \), \( V_n  \) are independent. First note that by assumption $Qf$ is a.s.\ well-defined (c.f.\ Remark 4.4 in \cite{4}). By Proposition G.2 in \cite{4} for \( \varepsilon_i \overset{\text{i.i.d.}}{\sim} \text{Exp}(1) \) and \( W_{ni} := \varepsilon_i / \bar{\varepsilon}_n \), we have \( (W_{n1}, \ldots, W_{nn}) \sim \text{Dir}(n; 1/n, \dots, 1/n) \). Since \( \mathbb{B}_n \sim \text{DP}(n \mathbb{G}_n) \), for $\delta_{Z_i}$ Dirac measure at $Z_i$, we have the following representation in distribution of $\mathbb{B}_n$ (c.f.\ Example 3.7.9 in \cite{5}):
	\[
	\mathbb{B}_n = \frac{1}{n} \sum_{i=1}^n W_{ni} \delta_{Z_i}.
	\]
	Because of this representation (c.f.\ section 3.7.2. in \cite{5}), $\mathbb{B}_n$ is also known as \textit{Bayesian bootstrap}. Conclude that: 
	\[
		\frac{\sqrt{n}}{\sqrt{\log n}} (\mathbb{B}_n - \mathbb{G}_n) f = \frac{1}{\sqrt{n \log n}} \sum_{i=1}^n \left( W_{ni} - 1 \right) f(Z_i).
	\]
	 Because $|\bar{\varepsilon}_n - 1| = O_p(1/\sqrt{n})$, we study: $(n \log n)^{-1/2} \sum_{i=1}^n \left( \varepsilon_i - 1 \right) f(Z_i)$. Now note that by Lemma \ref{lemma: weak conv epsilons scaled} in the Appendix:
	\begin{align}\label{eq: weak conv for NBP}
		\frac{1}{\sqrt{n \log{n}}} \sum_{i=1}^n \frac{(\varepsilon_i -1)}{\sqrt{Z_i - x}} \ind_{\{ Z_i >x \}} \mid Z_1, \ldots, Z_n \rightsquigarrow N\left(0, g_0(x)\right) \, \, \text { under } \, G_0 \text {. }
	\end{align}
	Because $Qf$ is a.s.\ well-defined, we obtain the claim by Slutsky's lemma upon noticing $ \frac{\sqrt{n}}{\sqrt{\log{n}}} V_n Q f \stackrel{G_0}{\rightarrow} 0$.
\end{proof}
We consider a projection-type posterior by isotonizing draws from the obtained naive posterior \eqref{eq: naive bayes chapter3} as follows. First define, for $G \sim \text{DP}(\alpha + n \mathbb{G}_n)$ and $\mathbb{G}_n$, $G_0$ as above:
\begin{align*}
&U_{\scriptscriptstyle{G}}(x) := \int_{0}^x V_{\scriptscriptstyle{G}}(y) \, dy = 2 \int_{0}^\infty \sqrt{z} \, dG(z) - 2 \int_x^\infty \sqrt{z - x} \, dG(z), \numberthis \label{eq: U_G}\\
&U_n(x) := \int_{0}^x V_n(y) \, dy = 2 \int_{0}^\infty \sqrt{z} \, d\mathbb{G}_n(z) - 2 \int_x^\infty \sqrt{z - x} \, d\mathbb{G}_n(z),\numberthis \label{eq: U_n chapter3} \\
&U_0(x) := \int_{0}^x V_0(y) \, dy = 2 \int_{0}^\infty \sqrt{z} \, dG_0(z) - 2 \int_x^\infty \sqrt{z - x} \, dG_0(z). \numberthis \label{eq: U_0}
\end{align*}

Furthermore, for a continuous function $k : [0,\infty) \to [0,\infty)$, the Least Concave Majorant (LCM) and its right-hand side derivative, are given for $x \in [0,\infty)$ by:
\begin{align*}
    &k^*(x) := \inf {\{ f(x) \, : \, f \geq k \, \, \text{on} \, \, [0,\infty), \, f \, \text{concave} \, \, \text{on} \, \, [0,\infty) \}},\\
    & (k^*)^{\prime}_+ (x):= \lim_{h \to 0^+} \frac{k^*(x + h) - k^*(x)}{h}.
\end{align*}
We denote by $\hat{V}_n$ (IIE) the right-hand side derivative of $U^*_n$, the LCM of $U_n$. Similarly, $\hat{V}_{\scriptscriptstyle{G}}$ is the right-hand side derivative of $U^*_{\scriptscriptstyle{G}}$, the LCM of $U_{\scriptscriptstyle{G}}$, i.e., the isotonization of a draw from the naive Bayes posterior. Thus, for $G$ as in \eqref{eq: dirichlet posterior chapter3}:  
\begin{align*}
	&\hat{V}_n(x) := (U^*_n)^{\prime}_+ (x), \numberthis \label{eq: IIE}\\
	&\hat{V}_{\scriptscriptstyle{G}}(x) :=  (U^*_{\scriptscriptstyle{G}})^{\prime}_+(x). \numberthis \label{eq: IIP}
\end{align*} 
By Lemma 2 in \cite{1}, $\hat{V}_n$ and $\hat{V}_{\scriptscriptstyle{G}}$ can be viewed as (equivalent to) $\mathbb{L}_2$-projections of the naive estimators $V_n$ and $V_{\scriptscriptstyle{G}}$ into the space of nonnegative, decreasing and right-continuous functions on $[0,\infty)$. We can finally define the Naive posterior and the IIP to target the estimation of $F_0$ as follows:
\begin{align*}
    &F_{\scriptscriptstyle{G}}(x) := 1 -\frac{2}{\pi} \sqrt{x} \, V_{\scriptscriptstyle{G}}(x) - \frac{2}{\pi} \int_{x}^{\infty} \frac{V_{\scriptscriptstyle{G}}(s)}{2 \sqrt{s}} \, ds, \numberthis \label{eq: naive post F0} \\
	&\hat{F}_{\scriptscriptstyle{G}}(x) := 1 -\frac{2}{\pi} \sqrt{x} \, \hat{V}_{\scriptscriptstyle{G}}(x) - \frac{2}{\pi} \int_{x}^{\infty} \frac{\hat{V}_{\scriptscriptstyle{G}}(s)}{2 \sqrt{s}} \, ds,  \numberthis \label{eq: IIP F0}
\end{align*}
and analogously we define $F_n(x)$, with $V_n$ instead of $V_{\scriptscriptstyle{G}}$ in \eqref{eq: naive post F0} and $\hat{F}_n(x)$, with $\hat{V}_n$ instead of $\hat{V}_{\scriptscriptstyle{G}}$ in \eqref{eq: IIP F0}.

We now state two additional conditions needed for Theorem \ref{thm: BvM}. First, a natural assumption is a moment constraint for $Z$, (which can be translated into a moment condition on $X$ c.f.\ \cite{27}): 
\begin{align*}
    \quad \int_{0}^{\infty} z \, g_0(z) \, dz < \infty. \numberthis \label{eq: finite first moment chapter3}
\end{align*}
This is a mild assumption, which is practically always satisfied, given that materials studied in materials science are of bounded size and globular clusters typically do not extend indefinitely. Additionally, assume $\exists$ $K >0$ and $\gamma > \frac{1}{2}$ such that
\begin{align*}
    &  \int_{0}^{1} ( V_0(x) - V_0(x+u \delta) ) \, du \sim \operatorname{sgn}(\delta)|\delta|^{\gamma} K, \quad \quad \text{as} \: \delta \rightarrow 0, \numberthis \label{eq: condition roughness} 
\end{align*}
where $\sim$ means that the ratio of the left-hand side and of the right-hand side goes to 1 as $\delta \rightarrow 0$. The constant $\gamma$ indicates the level of smoothness at $x$ of $V_0$. This is analogous to equation (5) from \cite{27}. Note that it is possible to see what this condition means in terms of the underlying $F_0$, using \eqref{eq: solution range equation}, we obtain that $\exists$ $K >0$ and $\gamma > \frac{1}{2}$ such that as $\delta \rightarrow 0$:
\begin{align}\label{eq: condition roughness F0}
	\int_{0}^{1} ( \sqrt{x} \, V_0(x) - \sqrt{x+ u \delta} \, V_0(x+u \delta) ) \, du \sim  \operatorname{sgn}(\delta)|\delta|^{\gamma} K - \frac{V_0(x)}{4 \sqrt{x}} \operatorname{sgn}(\delta)|\delta|, 
\end{align}
(upon using a Taylor expansion for $\int_0^1 (\sqrt{x} - \sqrt{x+ u \delta}) \, du$). Because $\frac{2}{\pi} \int_{x}^{\infty} \frac{V_0(s)}{2 \sqrt{s}} \, ds$ is Lipschitz continuous, the local smoothness of $F_0$ at $x$ is $\min(\gamma, 1)$. Therefore for the most relevant levels of smoothness, i.e.\ $\gamma \in (\frac{1}{2},1]$, the smoothness of $V_0$ carries over to $F_0$. Because $ 
F_0^*(x) = 1 - \frac{V_0(x)}{V_0(0)}
$, equation (5) from \cite{27} is equivalent to \eqref{eq: condition roughness} (note that in that paper $F$ denotes the current $F_0^*$). Thus, for a better mathematical understanding of the role of $\gamma$, we refer the interested reader to the discussion given in \cite{27}, in particular Example 2.1, 2.2, 2.3 and Remark 2.1. 

From a physical point of view, the role of $\gamma$ is also interpretable. In a globular cluster, different regions (e.g., dense core vs. sparse halo) may correspond to different levels of smoothness of $F_0$, and the semiparametric theory reflects this. The local adaptivity of the Bayesian estimator is encoded in the limiting variance of the Bernstein-von Mises mirrors this phenomenon. If there is a sharp increase in the concentration of stars with size close to $x$ (i.e., $\gamma$ close to $1/2$), then the 2D observations are much noisier. This makes the estimation more difficult and results in a higher asymptotic variance. 

\begin{theorem}[BvM for IIP]\label{thm: BvM}
	Let $x > 0$, suppose that \eqref{eq: condition finite g at x chapter3}, \eqref{eq: finite first moment chapter3}, and \eqref{eq: condition roughness} hold true. Moreover let $G \sim \operatorname{DP}(\alpha)$, where the base measure $\alpha$ has bounded density. Then, in probability under $G_0$: 
	\begin{align}\label{eq: claim}
		\frac{\sqrt{n}}{\sqrt{\log n}}\left(\hat{V}_{\scriptscriptstyle{G}}(x)-\hat{V}_n(x)\right) \mid Z_1, \ldots, Z_n \rightsquigarrow N\left(0, \frac{g_0(x)}{2 \gamma}\right).
	\end{align}
\end{theorem}
\begin{proof}

The proof essentially consists of three parts. The first part is to rewrite the event
$$
\left\{\sqrt{\frac{n}{\log n}}\left(\hat{V}_{\scriptscriptstyle{G}}(x) - \hat{V}_n(x)\right) \leq a\right\}
$$
given $Z_1,\ldots,Z_n$ in terms of the location of the maximum (arg max) of a stochastic process $t \mapsto \tilde{Z}^{\scriptscriptstyle{G}}_n(t)-a t$. It turns out that the process $\tilde{Z}^{\scriptscriptstyle{G}}_n(t)$ can be decomposed into three essential parts. The second step consists in studying these processes. One process, which can be interpreted as the Bayesian bootstrap counterpart of the process studied in the proof of Theorem 2.1 in \cite{27} (c.f.\ eq.\ (14) in \cite{27}). This process converges conditionally in distribution to a Gaussian process $\mathbb{Z}$, which we show can be written as $\mathbb{Z}(t) = tX$, where $X \sim N(0,g_0(x)/(2\gamma))$. A second process (see \eqref{eq: deterministic part process}), which is a deterministic process, whose behavior is controlled by the local smoothness of $V_0$ at $x$ and which converges to $- |t|^{\gamma+1} K$ (with $\gamma$ and $K$ as in \eqref{eq: condition roughness}). Lastly, a third remainder process which converges uniformly to zero with respect to the topology of uniform convergence on compacta. The final step is to prove that from the conditional convergence in distribution of $t \mapsto \tilde{Z}^{\scriptscriptstyle{G}}_n(t)-a t$ to $\mathbb{Z}(t)-a t$ it follows that the arg max of ($\tilde{Z}^{\scriptscriptstyle{G}}_n(t)-a t$) converges in distribution to the arg max of $\mathbb{Z}(t)-a t$. Since the arg max functional is certainly not continuous on the space of continuous functions equipped with the uniform norm, the simplest form of the continuous mapping theorem cannot be applied. For this reason, we prove in Lemma \ref{lemma: cond argmax} specifically for this application: a conditional argmax continuous mapping theorem (extension of Theorem 3.2.2 in \cite{5}). From \eqref{eq: conditional stoch boundedness chapter3} below, it follows that the sequence of locations of the maximum of ($\tilde{Z}_n(t)-a t$) is conditionally uniformly tight. This enables us to use Lemma \ref{lemma: cond argmax}. Moreover, since the limiting process turns out almost surely to have a unique maximum, the asserted conditional convergence in distribution in \eqref{eq: claim} follows.


For sequences $\delta_n = \sqrt{\log{n}}/\sqrt{n}$ and $\delta_n^{*}=\delta_n^{1 / \gamma}$ where $\gamma >1 / 2$ satisfies condition \eqref{eq: condition roughness}, define
\begin{align}\label{eq: Z_n process}
	Z^{\scriptscriptstyle{G_0}}_n(t) & :=\delta_n^{-1}\left(\delta_n^*\right)^{-1}\left\{U_n\left(x+\delta_n^* t\right)-U_n(x)-U_0\left(x+\delta_n^* t\right)+U_0(x)\right\} \\
	& =2 \delta_n^{-1} (\delta^*_n)^{-1} \int\left\{\sqrt{(z-x)_+} -\sqrt{(z-x-\delta^*_n t)_+} \right\} \, d\left(\mathbb{G}_n-G_0\right)(z).
\end{align}
By Lemma \ref{eq: lemma equiv to claim} in Appendix A it suffices to show, for $X \sim N\left(0, \frac{g_0(x)}{2 \gamma}\right)$ and every $a \in \mathbb{R}$:
\begin{align}\label{eq: equiv to claim}
	\Pi_n\left( \sqrt{\frac{n}{\log n}}\left(\hat{V}_{\scriptscriptstyle{G}}(x)-V_0(x)\right) -Z^{\scriptscriptstyle{G_0}}_n(1) \leq a \mid Z_1, \ldots, Z_n \right) \xrightarrow{G_0} \mathrm{P}(X \leq a).
\end{align}
Define the process \( T_{\scriptscriptstyle{G}} \) for \( a > 0 \):
\[
T_{\scriptscriptstyle{G}}(a) := \inf \{ t \geq 0 : U_{\scriptscriptstyle{G}}(t) - at \text{ is maximal} \}.
\]
By the switch relation (see Lemma 3.2 in \cite{2}):
\[
T_{\scriptscriptstyle{G}}(a) \leq x \iff \hat{V}_{\scriptscriptstyle{G}}(x) \leq a.
\]
Therefore, for sufficiently large \( n \), conditionally on $Z_1, \ldots, Z_n$:
\[
\delta_n^{-1} ( \hat{V}_{\scriptscriptstyle{G}}(x) - V_0(x) ) - Z^{\scriptscriptstyle{G_0}}_n(1) \leq a 
\iff \left( \delta_n^* \right)^{-1} \left( T_{\scriptscriptstyle{G}} \left( V_0(x) + \delta_n (a + Z^{\scriptscriptstyle{G_0}}_n(1)) \right) - x \right) \leq 0.
\]
Because the location of a maximum of a function is equivariant under translations and invariant under multiplication by a positive number and addition of a constant:
\begin{align*}
	&\left( \delta_n^* \right)^{-1} \left( T_{\scriptscriptstyle{G}} \left( V_0(x) + \delta_n (a + Z^{\scriptscriptstyle{G_0}}_n(1)) \right) - x \right)\\
	& = \left( \delta_n^* \right)^{-1} \big( \inf \big\{ x + t \geq 0 : U_{\scriptscriptstyle{G}}(x + t) - V_0(x)(x+t) \\
    &\quad \quad \quad \quad \quad - \delta_n (a + Z^{\scriptscriptstyle{G_0}}_n(1))(x + t) \, \, \text{ is max} \big\} - x \big)\\
	&= \inf \big\{t \geq-(\delta^*_n)^{-1} x: U_{\scriptscriptstyle{G}}\left(x+\delta^*_n t\right)-U_{\scriptscriptstyle{G}}(x) -U_n\left(x+\delta^*_n t\right)+U_n(x)\\
	&\quad + U_0(x+\delta^*_nt) - U_0(x)-V_0(x) \delta^*_n t-\delta_n \delta^*_n (a t + Z^{\scriptscriptstyle{G_0}}_n(1)t - Z^{\scriptscriptstyle{G_0}}_n(t)) \, \text { is max}\big\} \numberthis \label{eq: useful for stoch}\\
	&= \inf \big\{t \geq-(\delta^*_n)^{-1} x: \tilde{Z}^{\scriptscriptstyle{G}}_n(t)-a t \, \text { is max}\big\}.
\end{align*}
where we define the localized process \( \tilde{Z}_n^{\scriptscriptstyle{G}} \) by,
\begin{align*}
\tilde{Z}_n^{\scriptscriptstyle{G}}(t) &:= \delta_n^{-1} \left( \delta_n^* \right)^{-1} \big\{ U_{\scriptscriptstyle{G}} \left( x + \delta_n^* t \right) - U_{\scriptscriptstyle{G}}(x) - U_n \left( x + \delta_n^* t \right) + U_n(x) \big\} \numberthis \label{eq: first components Z tilde}\\
&\quad + \delta_n^{-1} \left( \delta_n^* \right)^{-1} \big\{ U_0 \left( x + \delta_n^* t \right) - U_0(x) - V_0(x) \delta^*_n t \big\} +  Z^{\scriptscriptstyle{G_0}}_n(t) -  t Z^{\scriptscriptstyle{G_0}}_n(1) \numberthis \label{eq: second components Z tilde}
\end{align*}
for $t \in I_x$, where $I_x := [0, \infty) \text{ if } x = 0, \; \text{and} \; I_x := \mathbb{R} \text{ if } x > 0$.

We determine the asymptotic behavior of the process $\tilde{Z}^{\scriptscriptstyle{G}}_n$ by studying separately the components in \eqref{eq: first components Z tilde} and \eqref{eq: second components Z tilde}, and next we conclude by applying Slutsky's lemma. First note that $|Z^{\scriptscriptstyle{G_0}}_n(1)t - Z^{\scriptscriptstyle{G_0}}_n(t)| = o_{\scriptscriptstyle{G_0}}(1)$ uniformly in $t$ in compacta by the proof of Theorem 2.1 in \cite{27} (c.f.\ display after equation (16)). Furthermore, using assumption \eqref{eq: condition roughness} (as in equation (13) in \cite{27}), uniformly in $t$ in compacta:
\begin{align}\label{eq: deterministic part process}
	\delta_n^{-1} \left( \delta_n^* \right)^{-1} \big\{ U_0 \left( x + \delta_n^* t \right) - U_0(x) - V_0(x) \delta^*_n t \big\} =  - |t|^{\gamma+ 1} (K + o(1)).
\end{align}
This determines the behavior of \eqref{eq: second components Z tilde}. Now we study the posterior behaviour of the process in \eqref{eq: first components Z tilde}, which can be interpreted as a \textit{Bayesian bootstrap} counterpart of the process \eqref{eq: Z_n process} studied in the proof of Theorem 2.1 in \cite{27} (c.f.\ eq.\ (14) in \cite{27}):
\begin{align*}
	Z^{\scriptscriptstyle{G}}_n(t) &:= \delta_n^{-1} \left( \delta_n^* \right)^{-1} \big\{ U_{\scriptscriptstyle{G}} \left( x + \delta_n^* t \right) - U_{\scriptscriptstyle{G}}(x) - U_n \left( x + \delta_n^* t \right) + U_n(x) \big\} \\
	&= 2 \delta_n^{-1} \left( \delta_n^* \right)^{-1} \int\left\{\sqrt{(z-x)_+} -\sqrt{(z-x-\delta^*_n t)_+} \right\} \, d\left(G-\mathbb{G}_n\right)(z)
\end{align*}
As in proof of Theorem \ref{thm: BvM naive} in equation \eqref{eq: first representation in distribution}, we have by Proposition G.10 in \cite{4}, for independent \( Q \sim \text{DP}(\alpha) \), \( \mathbb{B}_n \sim \text{DP}(n \mathbb{G}_n) \), and \( V_n \sim \text{Be}(|\alpha|, n) \), the following representation in distribution:
\begin{align}\label{eq: representation process G}
	G f_t^n = V_n Q f_t^n + (1 - V_n) \mathbb{B}_n f_t^n,
\end{align}
where $f_t^n(z) := 2(\delta_n^*)^{-1} \big\{ \sqrt{(z-x)_+} -\sqrt{(z-x-\delta^*_n t)_+}  \big\}$ and thus the following representation in distribution of the process \( Z^{\scriptscriptstyle{G}}_n(t) \):
\[
Z^{\scriptscriptstyle{G}}_n(t) = \delta_n^{-1} V_n (Q - \mathbb{G}_n) f_t^n + \delta_n^{-1} (1 - V_n) (\mathbb{B}_n - \mathbb{G}_n) f_t^n.
\]
This representation is crucial for the asymptotic analysis of the process \( Z^{\scriptscriptstyle{G}}_n(t) \). The first thing to do is to identify which parts of the process \( Z^{\scriptscriptstyle{G}}_n(t) \) do not give any contribution to limiting dstribution of the process, as they converge conditionally to zero in probability. The second part is to identify the limiting distribution of the remaining parts of the process, which will be a Gaussian process. In Appendix A, we give the \hyperlink{proof of reduction}{proof} that uniformly in $t$ in compacta:
\begin{align}\label{eq: representation without remainders chapter3}
	Z^{\scriptscriptstyle{G}}_n(t) = \mathbb{Z}^*_n (t) + o_p(1),
\end{align}
where for \( \varepsilon_i \overset{\text{i.i.d.}}{\sim} \text{Exp}(1) \), the process $\mathbb{Z}^*_n$ is defined as:
\begin{align}\label{eq: mathbbZ process}
    \mathbb{Z}^*_n (t) := \frac{1}{\sqrt{n \log n}} \sum_{i=1}^n \left( \varepsilon_i - 1 \right) f_t^n(Z_i).
\end{align}
Then by Slutsky's lemma, if we determine the conditional weak limit of \( (\mathbb{Z}_n^\ast(t) \, : \, t \in I_x) \), then we determine the one of \( (Z^{\scriptscriptstyle{G}}_n(t) \, : \, t \in I_x) \). We now study the marginals of $\mathbb{Z}^*_n$. Note that by Lemma \ref{lemma: result with cov structure} for fixed $t \in I_x$, in probability under $G_0$:
\begin{align}\label{eq: weak conv for IIP}
	\mathbb{Z}^*_n (t) \mid Z_1, \ldots, Z_n \rightsquigarrow N\left(0, t^2  g_0(x)/ (2 \gamma) \right),
\end{align}
and again by Lemma \ref{lemma: result with cov structure} that:
\begin{align}\label{eq: cov structure}
\text{Cov}\left(\mathbb{Z}_n^\ast(t), \mathbb{Z}_n^\ast(s) \,\middle|\, Z_1, \dots, Z_n\right) 
= \frac{1}{2 \gamma} g_0(x) st \left( 1 - O_{\scriptscriptstyle{G_0}}\left(\frac{\log \log n}{\log n} \right) \right) + O_{\scriptscriptstyle{G_0}}\left(\frac{1}{\log n}\right),
\end{align}
By the conditional Chebyshev inequality and \eqref{eq: cov structure} we conclude $\forall \: \varepsilon >0$:
\begin{align*}
    \pb \left( | \mathbb{Z}^*_n (t) - t \mathbb{Z}^*_n (1) | > \varepsilon \mid Z_1, \ldots, Z_n \right) \leq \frac{1}{\varepsilon^2} \mathrm{Var}( \mathbb{Z}^*_n (t) - t \mathbb{Z}^*_n (1)  \mid Z_1, \ldots, Z_n ) \stackrel{G_0}{\rightarrow} 0.
\end{align*}
By \eqref{eq: representation without remainders chapter3} conclude that \eqref{eq: weak conv for IIP} holds true with $\mathbb{Z}^*_n $ replaced by $Z^{\scriptscriptstyle{G}}_n$. Therefore, the finite-dimensional distributions of $Z^{\scriptscriptstyle{G}}_n$ converge conditionally in distribution to the finite-dimensional distributions of the process: $ \mathbb{Z}(t) = tX$, where $X \sim N(0,g_0(x)/(2\gamma))$. 

Lemma \ref{lemma: cond asymp equicontinuity} gives the conditional asymptotic equicontinuity of the process $\mathbb{Z}^*_n$. The conditional asymptotic equicontinuity together with the conditional convergence in distribution of the finite dimensional distributions, implies in probability under $G_0$ as $n \rightarrow \infty$:
\begin{align}\label{eq: weak conv Z^*_n}
	\mathbb{Z}^*_n \mid Z_1,\ldots, Z_n \rightsquigarrow \mathbb{Z} \quad \text{in} \quad \ell^{\infty}(K).
\end{align}
for every compact $K \subset I_x$ and $\ell^{\infty}$ being the space of bounded functions on $K$ equipped with the supremum norm. For a review of the conditional weak convergence in $\ell^{\infty}$, see \cite{5} (c.f.\ Example 1.5.1 and Theorem 1.5.4). By \eqref{eq: representation without remainders chapter3} and \eqref{eq: weak conv for IIP} conclude that \eqref{eq: weak conv Z^*_n} holds true with $\mathbb{Z}^*_n $ replaced by $Z^{\scriptscriptstyle{G}}_n$, i.e.\ the conditional weak limit of $Z^{\scriptscriptstyle{G}}_n$ relative to the topology of uniform convergence on compacta is the same Gaussian process. Altogether, recalling also \eqref{eq: deterministic part process}, we can conclude that conditionally on $Z_1,\ldots, Z_n$ the weak limit of $(\tilde{Z}_n^{\scriptscriptstyle{G}}(t) - at \, : \, t \in K )$ in $\ell^{\infty}$ is given by $( tX - |t|^{\gamma+1} K - at \, : \, t \in K )$ for every compact $K\subset I_x$.

In order to apply Lemma \ref{lemma: cond argmax}, notice that \( \mathbb{Z}\) has a.s.\ continuous sample paths and a.s.\ unique maxima (see Lemma 2.6 in \cite{12}) and that $\hat{t}^{\scriptscriptstyle{G}}_n :=  \inf \big\{t \geq-(\delta^*_n)^{-1} x: \tilde{Z}^{\scriptscriptstyle{G}}_n(t)-a t \, \text { is max}\big\}$ is an argmax process conditionally on $Z_1, \ldots, Z_n$, satisfying the second condition on Lemma \ref{lemma: cond argmax}. It is thus enough to show that there exists a conditionally stochastically bounded sequence \((\tilde{M}_n)\) that satisfies:
\begin{align}\label{eq: conditional stoch boundedness chapter3}
\Pi_n \left( \left( \delta_n^* \right)^{-1} \left| T_{\scriptscriptstyle{G}} \left( V_0(x) + \delta_n (a + Z^{\scriptscriptstyle{G_0}}_n(1)) \right) - x \right| \leq \tilde{M}_n \mid Z_1, \ldots, Z_n \right) 
\stackrel{G_0}{\to} 1,
\end{align}
The \hyperlink{proof of stoch boundedness}{proof} of \eqref{eq: conditional stoch boundedness chapter3} is given in Appendix A. Thus by Lemma \ref{lemma: cond argmax} we can conclude:
\begin{align*}
&\Pi_n\left( \sqrt{\frac{n}{\log n}}\left(\hat{V}_{\scriptscriptstyle{G}}(x)-V_0(x)\right) -Z^{\scriptscriptstyle{G_0}}_n(1) \leq a \mid Z_1, \ldots, Z_n \right)  \\
& \quad \stackrel{G_0}{\longrightarrow} \pb \left\{ \underset{t \in I_x}{\arg \max }\left\{t X-|t|^{\gamma+ 1}  K  -a t \right\} \leq 0 \right\} = \pb\{ X \leq a\},
\end{align*}
because $t \mapsto t(X-a) - |t|^{\gamma +1}K$ attains its maximum on $(-\infty,0]$ iff $X-a\leq 0$.
\end{proof}

We conclude this section by showing the implications of the above results for the uncertainty quantification of the posterior. Under the conditions of Theorem \ref{thm: BvM} (c.f.\ \cite{27}) 
\begin{align*}
	\frac{\sqrt{n}}{\sqrt{\log n}}\left(\hat{V}_n(x) - V_0(x)\right) \rightsquigarrow N\left(0, \frac{g_0(x)}{2 \gamma}\right).
\end{align*}
Thus as consequence of \eqref{eq: solution range equation} and the analogous definition of $\hat{F}_n(x)$, we have that (see proof of Corollary \ref{thm: BvM for F} in the Appendix):
\begin{align}\label{eq: convergence IIE for F}
	\frac{\sqrt{n}}{\sqrt{\log n}}\left(\hat{F}_n(x) - F_0(x)\right) \rightsquigarrow N\left(0, \frac{2 x}{\pi^2 \gamma} g_0(x)\right).
\end{align}
Together with \eqref{eq: convergence for F} below, this will imply that the credible intervals constructed from the Isotonized Inverse Posterior (IIP) of $\hat{F}_{\scriptscriptstyle{G}}(x)$ are (asymptotic) frequentist confidence sets. The following Corollary \ref{thm: BvM for F}, whose proof is given in the Appendix, summarizes formally the above discussion. Note that because the uncertainty quantification based on the IIP is always preferable to the one based on the Naive Bayes Posterior, we only propose uncertainty quantification using the latter.

\begin{corollary}[BvM for $\hat{F}_{\scriptscriptstyle{G}}$ and $F_{\scriptscriptstyle{G}}$ and UQ for $F_0$]\label{thm: BvM for F} Let $F_{\scriptscriptstyle{G}}$ as in \eqref{eq: naive post F0} and $\hat{F}_{\scriptscriptstyle{G}}$ as in \eqref{eq: IIP F0} and analogously let $F_n(x)$, with $V_n$ instead of $V_{\scriptscriptstyle{G}}$ in \eqref{eq: naive post F0} and $\hat{F}_n(x)$, with $\hat{V}_n$ instead of $\hat{V}_{\scriptscriptstyle{G}}$ in \eqref{eq: IIP F0}. Under the same conditions of Theorem~2.1 in \cite{65}, in probability under $G_0$ as $n \rightarrow \infty$:
    \begin{align}
        &\frac{\sqrt{n}}{\sqrt{\log n}}\left(F_{\scriptscriptstyle{G}}(x)-F_n(x)\right) \mid Z_1, \ldots, Z_n \rightsquigarrow N\left(0, \frac{4 x}{\pi^2} g_0(x)\right). \numberthis \label{eq: convergence for naive F} 
    \end{align}
    Under the same conditions of Theorem~2.2 in \cite{65}, in probability under $G_0$ as $n \rightarrow \infty$:
    \begin{align}
        &\frac{\sqrt{n}}{\sqrt{\log n}}\left(\hat{F}_{\scriptscriptstyle{G}}(x)-\hat{F}_n(x)\right) \mid Z_1, \ldots, Z_n \rightsquigarrow N\left(0, \frac{2 x}{\pi^2 \gamma} g_0(x)\right).\numberthis \label{eq: convergence for F}
    \end{align}
    Further, if we define the posterior quantiles for isotonized posterior $\hat{F}_{\scriptscriptstyle{G}}$, for $\alpha \in [0,1]$, as:
    \begin{align*}
        q_{\alpha} := \inf \{ t \in \mathbb{R} \, : \, \Pi_n (\hat{F}_{\scriptscriptstyle{G}}(x) \leq t \mid Z_1, \ldots, Z_n) \geq \alpha\}
    \end{align*}	
    then under $G_0$ as $n \rightarrow \infty$:
    \begin{align}\label{eq: UQ}
        \pb \left( F_0(x) \in [q_{\alpha/2}, q_{1- \alpha/2}]\right) \rightarrow 1- \alpha.
    \end{align}
    \end{corollary}
    
    \begin{proof}
		We start by proving \eqref{eq: convergence for F}, as the proof of \eqref{eq: convergence for naive F} is analogous but simpler. First note that by the inverse relation \eqref{eq: solution range equation} inherent to Wicksell's problem:
		\begin{align*}
			\hat{F}_{\scriptscriptstyle{G}}(x)-\hat{F}_n(x) = -\frac{2}{\pi} \sqrt{x}\left\{\hat{V}_{\scriptscriptstyle{G}}\left(x\right)-\hat{V}_n\left(x\right)\right\}-\frac{2}{\pi} \int_{x}^{\infty} \frac{\hat{V}_{\scriptscriptstyle{G}}(u)-\hat{V}_n(u)}{2 \sqrt{u}} \, du .
		\end{align*}
		From Theorem~2.2 in \cite{65}, we obtain the weak conditional limit of $\frac{\sqrt{n}}{\sqrt{\log{n}}}\left\{\hat{V}_{\scriptscriptstyle{G}}\left(x\right)-\hat{V}_n\left(x\right)\right\} \mid Z_1,\ldots,Z_n$. Integrating by parts, the integral on the last display is at most
		$$
		\begin{aligned}
			\frac{\left|U^*_{\scriptscriptstyle{G}}(x)-U_n^{*}(x)\right|}{2 \sqrt{x}}+\frac{1}{4}\left|\int_{x}^{\infty} \frac{U^*_{\scriptscriptstyle{G}}(u)-U_n^{*}(u)}{u^{3 / 2}} \, d u\right| &\leq \frac{\left\|U^*_{\scriptscriptstyle{G}}-U_n^{*}\right\|_{\infty}}{2 \sqrt{x}}+\left\|U^*_{\scriptscriptstyle{G}}-U_n^{*}\right\|_{\infty} \frac{1}{2 \sqrt{x}} \\
		& =\frac{\left\|U^*_{\scriptscriptstyle{G}}-U_n^{*}\right\|_{\infty}}{\sqrt{x}} \leq \frac{\left\|U_{\scriptscriptstyle{G}}-U_n\right\|_{\infty}}{\sqrt{x}}.
		\end{aligned}
		$$
		where in the last inequality we applied a generalized version of Marshall's inequality (c.f.\ Lemma 2.2 in \cite{66}). Now note that $\sqrt{n} \left\|U_{\scriptscriptstyle{G}}-U_n\right\|_{\infty} \mid Z_1,\ldots,Z_n$ has a conditional weak limit by Theorem 12.2 in \cite{4} combined with the Continuous Mapping Theorem. We can apply Theorem 12.2 because the class of functions $\mathcal{F} := \{ z \mapsto \sqrt{(z)_+} - \sqrt{(z - s)_+} \, : \, s \geq 0 \}$ is Donsker by Lemma 2.1 in \cite{27}. Conclude thus by Slutsky's lemma upon noticing $\frac{\sqrt{n}}{\sqrt{\log{n}}} \left\|U_{\scriptscriptstyle{G}}-U_n\right\|_{\infty} \mid Z_1,\ldots,Z_n \, \to 0$ under $G_0$.    

		\eqref{eq: convergence for naive F} follows using the fact that $\sqrt{n} \left\|U_{\scriptscriptstyle{G}}-U_n\right\|_{\infty} \mid Z_1,\ldots,Z_n$ has a conditional weak limit and the rewriting:
		\begin{align*}
			F_{\scriptscriptstyle{G}}(x)-F_n(x) = -\frac{2}{\pi} \sqrt{x}\left\{V_{\scriptscriptstyle{G}}\left(x\right)-V_n\left(x\right)\right\}-\frac{2}{\pi} \int_{x}^{\infty} \frac{V_{\scriptscriptstyle{G}}(u)-V_n(u)}{2 \sqrt{u}} \, du .
		\end{align*}
		
		We now show the third claim \eqref{eq: UQ}. By \eqref{eq: convergence for F} and Lemma 21.2 in \cite{6}, in probability under $G_0$ as $n \rightarrow \infty$:
	\begin{align*}
		\widetilde{q}_{\alpha} &:= \inf \left\{ t \in \mathbb{R} \, : \, \Pi_n \left( \frac{\sqrt{n}}{\sqrt{\log n}}\left(\hat{F}_{\scriptscriptstyle{G}}(x)-\hat{F}_n(x)\right) \leq t \mid Z_1, \ldots, Z_n \right) \geq \alpha \right\} \\
		&\quad \stackrel{G_0}{\longrightarrow} \sqrt{\frac{2 x g_0(x)}{\pi^2 \gamma}} \Phi^{-1}(\alpha), 
	\end{align*}
	where $\Phi^{-1}$ is the quantile function of the standard normal distribution. Then note:
	\begin{align*}
		q_{\alpha} &= \inf_{t \in \mathbb{R}} \left\{  \Pi_n \left( \frac{\sqrt{n}}{\sqrt{\log n}}\left(\hat{F}_{\scriptscriptstyle{G}}(x)-\hat{F}_n(x)\right) \leq \frac{\sqrt{n}}{\sqrt{\log n}} \left( t - \hat{F}_n(x) \right)\mid Z_1, \ldots, Z_n \right) \geq \alpha \right\} \\
		&= \hat{F}_n(x) + \frac{\sqrt{\log n}}{\sqrt{n}} \widetilde{q}_{\alpha} = \hat{F}_n(x) +  \sqrt{\frac{\log{n}}{n} \cdot \frac{2 x g_0(x)}{\pi^2 \gamma} } \, \Phi^{-1}(\alpha) + o_{\scriptscriptstyle{G_0}}\left(\sqrt{\frac{\log n}{n}}\right).
	\end{align*}
	 Note that
    \begin{align}\label{eq: convergence IIE for F proof}
	\frac{\sqrt{n}}{\sqrt{\log n}}\left(\hat{F}_n(x) - F_0(x)\right) \rightsquigarrow N\left(0, \frac{2 x}{\pi^2 \gamma} g_0(x)\right)
    \end{align}
    follows by the same rewriting in \eqref{eq: solution range equation} once we notice, by Marshall's inequality (c.f.\ Exercise 3.1 in \cite{2}) and Lemma 2.1 in \cite{27}, that as $n \rightarrow \infty$:
		\begin{align*}
			\bigg| \frac{\sqrt{n}}{\sqrt{\log{n}}} \int_{x}^{\infty} \frac{\hat{V}_n(u)- V_0(u)}{\pi \sqrt{u}} \, du \bigg| \leq \frac{\sqrt{n}}{\sqrt{\log{n}}} \frac{\left\|U_n - U_0\right\|_{\infty}}{\sqrt{x}} \stackrel{G_0}{\longrightarrow} 0.
		\end{align*}
	Using \eqref{eq: convergence IIE for F proof} and the symmetry of the normal distribution, for $X \sim N(0, (2 x g_0(x))/ (\pi^2 \gamma))$ 
	\begin{align*}
		&\pb \left( F_0(x) \in [q_{\alpha/2}, q_{1- \alpha/2}]\right) = \pb \left( \sqrt{\frac{n}{\log{n}}} \left( \hat{F}_n(x)- F_0(x) \right) \in \left[-\widetilde{q}_{1-\alpha/2}, -\widetilde{q}_{\alpha/2} \right]  \right)\\
		&\quad \quad \quad \rightarrow \pb \left( X \in \sqrt{ \frac{2 x g_0(x)}{\pi^2 \gamma} } [\Phi^{-1}(\alpha/2), \Phi^{-1}(1-\alpha/2)]\right) = 1 - \alpha,
	\end{align*}
	where we assumed that $\alpha \in [0,1]$.
	\end{proof}

\subsection[Simulation Study and Discussion about UQ]{Simulation Study and Discussion about Uncertainty Quantification}
\sectionmark{Simulation Study and Discussion about UQ}

In this section, we present simulations to illustrate the practical behavior of the IIP in comparison to the IIE studied in \cite{27}. Furthermore we present a simulation study that compares the IIP-uncertainty quantification to the one proposed in \cite{19}.

\begin{figure}[H]
	\centering
    \begin{subfigure}[t]{0.48\linewidth}
        \centering
        \includegraphics[width=\linewidth]{200_1.pdf}
        \caption*{$n=200$}
    \end{subfigure}
    \hspace{-0.2cm}
    \begin{subfigure}[t]{0.48\linewidth}
        \centering
        \includegraphics[width=\linewidth]{2000_1.pdf}
        \caption*{$n=2000$}
    \end{subfigure}
	\\
    \centering
	\begin{subfigure}[t]{0.5\linewidth}
        \centering
        \includegraphics[width=\linewidth]{2000_zoom_1.pdf}
		\caption*{Zoom for $n=2000$}
    \end{subfigure}
	\\
	\vspace{0.2cm}
    \begin{subfigure}[t]{0.48\linewidth}
        \centering
        \includegraphics[width=\linewidth]{2000_h.pdf}
        \caption*{Histograms}
    \end{subfigure}
    \hspace{-0.3cm}
    \begin{subfigure}[t]{0.52\linewidth}
        \centering
		\vspace*{-4.65cm}
        \includegraphics[width=\linewidth]{2000_qq.pdf}
        \caption*{QQ-plot}
    \end{subfigure}
	\caption{\textit{Synthetic} setting 1: uncertainty for different sample sizes and analysis of normality at $x=1.5$, based on 300 draws from the posterior.}\label{fig: analysis normality}
\end{figure}
We consider three distinct settings. The first two are \textit{synthetic} settings where the true underlying distribution function \( F_0 \) is known. The first two simulation settings serve as valuable benchmarks for assessing the quality of our methodology. In particular, the first one verifies that the variance attained by the IIP coincides with the theoretical one that we derived in Theorem \ref{thm: BvM for F}. The second \textit{synthetic} setting compares the uncertainty quantification of our methodology to the bootstrap-based one derived in \cite{19}. The last simulation setting is a \textit{real-data} application involving the positions of stars in the Globular Cluster Messier 62. In the last case, the ground truth distribution of the stars is unknown. Also in this setting, we compare our methodology with the bootstrap of \cite{19}. 

In Figure \ref{fig: analysis normality}, we analyze the setting where the underlying cdf \( F_0 \) is exponential with rate 1.2. Samples \( Z_1, \ldots, Z_n \) used to compute the projection posterior of the IIP are generated hierarchically: \( X_1, \ldots, X_n \sim F_0 \), \( Y_1, \ldots, Y_n \sim \mathrm{Be}(1, 1/2) \), with \( Z_i = X_i \cdot Y_i \) for all \( i \). The plots use the following color scheme: green for estimates of \( F \) based on prior draws from \( G \sim \operatorname{DP}(M \alpha) \) (parameters \( M \) and \( \alpha \) specified in the plots); purple for draws from the isotonized posterior; red for the average of these draws (approximating the posterior mean);  yellow for the Isotonic Inverse Estimator (IIE) of \cite{27}; and orange for \( F_0 \).  

The first row of Figure \ref{fig: analysis normality} compares the behavior of the IIP and the IIE as the sample size goes from 200 to 2000. With \( n = 200 \), the IIE and the average of the IIP draws already behave very similarly and with \( n = 2000 \), the uncertainty significantly reduces and the two estimation procedures essentially coincide. The average IIP appears to be somewhat smoother, which will typically be seen as an advantage. The middle row focuses on \( n = 2000 \), providing a zoomed-in view and highlighting \( x = 1.5 \), the point used for subsequent analysis of the normality. The final row examines normality (\( n = 2000 \)): histograms, QQ-plots, and statistical tests strongly confirm normality of the gathered samples. Furthermore, the Bayesian-estimated standard deviation aligns closely with the theoretical value attained by the frequentist procedure (c.f.\ \cite{27}, note that \( \gamma = 1 \) because \( F_0 \) is Lipschitz). In all the plots we employ 300 posterior draws. 
 
In Figure \ref{fig: compare UQ}, we analyze six different settings where at $x=5$ the underlying cdfs $F_0$ have varying degrees of smoothness, making them locally H\"{o}lder smooth of degree $\gamma \in \{ 0.55, 0.6, 0.65, 0.7, 0.8, 1.5 \}$. In particular, for an appropriate constant $K>0$, $F_0(y) = \frac{1}{2} - K (5-y)^\gamma$, for $y \in [0,5]$ and $F_0(y) = \frac{1}{2} + K (y-5)^\gamma$, for $y \in [5,10]$. The plots use the following color scheme: pink for the Bayesian-based credible sets; red for the approximated posterior mean; light blue for the bootstrap-based confidence sets of \cite{19}; yellow for the Isotonic Inverse Estimator (IIE) of \cite{27}; and lime green for \( F_0 \). The prior base measure used for these simulation is the same used in the previous simulation study. In all cases \( n = 2000 \) and we employed 300 posterior-based and bootstrap-based draws. This comparison confirms that the IIP delivers valid uncertainty quantification for varying degrees of smoothness and that the uncertainty quantification delivered by \cite{19} behaves closely, even though the authors formally showed it works only for a differentiable underlying truth $F_0$.

In Figure \ref{fig:stars}, we apply the methodology to star position data from globular cluster Messier 62. Let \( \mathbf{Z} = (Z_{p_1}, Z_{p_2}, Z_{p_3}) \) represent a star's three-dimensional position, where only \( (Z_{p_1}, Z_{p_2}) \) are observed. Here, \( X = Z_{p_1}^2 + Z_{p_2}^2 + Z_{p_3}^2 \sim F_0 \), and \( Z = Z_{p_1}^2 + Z_{p_2}^2 \). Beyond serving as a compelling application of our methodology, this example highlights a few key conclusions. First of all, also for real data the IIE closely resembles the IIP. However, because in this case the true $\gamma$ is unknown, it is not immediately clear how the frequentist uncertainty quantification compares with the Bayesian one. From the work of \cite{19}, we know that bootstrap procedures work for Wicksell's problem, at least when the underlying truth is differentiable at $x$. We thus compared our methodology with the one proposed in \cite{19} and we observed that they behave very similarly with $1000$ or more samples. In Figure \ref{fig:stars} the reader can observe an instance with $n=1400$.
\begin{figure}[H]
	\centering
    \begin{subfigure}[t]{0.52\linewidth}
        \centering
        \includegraphics[width=\linewidth]{gamma_055.pdf}
        %\caption*{$\gamma=0.55$}
    \end{subfigure}
    \hspace{-0.8cm}
    \begin{subfigure}[t]{0.52\linewidth}
        \centering
        \includegraphics[width=\linewidth]{gamma_06.pdf}
        %\caption*{$\gamma=0.55$}
    \end{subfigure}
	\\
    \begin{subfigure}[t]{0.52\linewidth}
        \centering
        \includegraphics[width=\linewidth]{gamma_065.pdf}
        %\caption*{Histograms}
    \end{subfigure}
    \hspace{-0.8cm}
    \begin{subfigure}[t]{0.52\linewidth}
        \centering
        \includegraphics[width=\linewidth]{gamma_07.pdf}
        %\caption*{QQ-plot}
    \end{subfigure}
	\\
    \begin{subfigure}[t]{0.52\linewidth}
        \centering
        \includegraphics[width=\linewidth]{gamma_08.pdf}
        %\caption*{Histograms}
    \end{subfigure}
    \hspace{-0.8cm}
    \begin{subfigure}[t]{0.52\linewidth}
        \centering
        \includegraphics[width=\linewidth]{gamma_15.pdf}
        %\caption*{QQ-plot}
    \end{subfigure}
	\caption{\textit{Synthetic} setting 2: comparison uncertainty quantification IIP and bootstrap method from \cite{19}.}\label{fig: compare UQ}
\end{figure}
\vspace{-1cm}
\begin{table}[H]
\centering
\resizebox{\textwidth}{!}{%
\begin{tabular}{c|cccccc}
   & $\gamma = 0.55$ & $\gamma = 0.6$ & $\gamma = 0.65$ & $\gamma = 0.7$ & $\gamma = 0.8$ & $\gamma = 1.5$ \\ \hline
Bootstrap CI width at $x = 5$ &  0.116   &  0.110    &  0.104   &  0.108   &   0.100  &   0.073  \\
Bayesian CI width at $x = 5$ &   0.114 &  0.121   &   0.103  &   0.101  &   0.099  &   0.077  \\
\end{tabular}}
%\caption{Your table caption here}
\label{tab:your_label}
\end{table}


\vspace{-0.5cm} 
\begin{figure}[H]
    \centering
	\begin{subfigure}{0.55\linewidth}
        \includegraphics[width=\linewidth]{ngc2622_1400_projected_stars.pdf}
    \end{subfigure}

    \begin{subfigure}{0.51\linewidth}
        \includegraphics[width=\linewidth]{1400_stars.pdf}
    \end{subfigure}  
	\hspace{-0.6cm}
    \begin{subfigure}{0.515\linewidth}
        \includegraphics[width=\linewidth]{CI_1400.pdf}
    \end{subfigure}

    \caption{\textit{Real-data} application: $n=1400$ projected star positions. On the left-hand side, the IIE from \cite{27} and the IIP with 300 posterior samples; on the right-hand side, a comparison between the $95\%$ bootstrap confidence set for the IIE from \cite{19} and the $95\%$ credible sets for the IIP.}
    \label{fig:stars}
\end{figure}



\subsubsection*{Discussion about Uncertainty Quantification}

The main limitation of the approach in \cite{27} is the lack of uncertainty quantification, which is essential for practical applications. The asymptotic variance of the limiting distribution achieved by the IIE involves two nuisance parameters: the local smoothness $\gamma$ and the density $g_0(x)$ of the observed data at the evaluation point $x$. While $g_0(x)$ may be estimable (for most practical applications the rate of convergence would be slower than $\sqrt{n}/\sqrt{\log n}$), estimating $\gamma$ is particularly challenging, as it requires recovering the local smoothness of the target distribution function $F_0$ at a specific point $x$ using samples from $g_0$, effectively a separate inverse problem. To our knowledge, no practical methods currently exist for this task. Moreover, because of the method proposed in this chapter, estimating $\gamma$ offers little practical benefit and thus makes this research direction of limited value. This highlights the key advantage of the Bayesian approach proposed in this chapter: it bypasses the need to estimate either $\gamma$ or $g_0(x)$, as the posterior distribution inherently provides valid uncertainty quantification over a wide range of local smoothness levels, i.e.\ all $\gamma > 1/2$. For $\gamma \leq 1/2$, the inverse problem becomes severely ill-behaved due to the unboundedness of the density of the observations $g_0$. This might even lead to a slower rate of convergence both for the naive plug-in estimator and for the IIE. Therefore, for all practically relevant cases with $\gamma > 1/2$, practitioners can apply our method without needing to estimate nuisance parameters or make strong assumptions about the form of $F_0$, thanks to its adaptive nature and the automatic uncertainty quantification.


The results displayed in Figure \ref{fig: compare UQ} suggest that the bootstrap method of \cite{19} behaves very closely to our method. We can thus make the empirically grounded hypothesis that the bootstrap method proposed in \cite{19} might work under the less stringent assumptions of the current chapter. The value of this simulation study has to be taken carefully, as a formal mathematical proof is still needed. However, proving this hypothesis is nontrivial and we therefore leave this as potential furure work. 

From a practitioner's perspective, thus, the key advantage of our Bayesian methodology lies in its provable robustness across a broad class of smoothness levels: it provides valid uncertainty quantification for all $\gamma > 1/2$. In contrast, the bootstrap approach in \cite{19} is provably consistent only when $F_0$ is differentiable at $x$, a substantially stronger condition (stricter than $\gamma = 1$). This distinction is crucial, since determining whether $F_0$ is differentiable at a given point based solely on samples from $g_0$ constitutes a challenging inverse problem both theoretically and in practice. By avoiding this requirement, our method offers a more broadly applicable and practitioner-friendly solution.







\subsection{Conclusion}

We have introduced a novel Bayesian approach for nonparametric estimation in Wicksell's problem, leveraging a Dirichlet Process prior on the distribution of observables rather than unobservables. This deviation from the classical Bayesian framework in inverse problems allows for computational efficiency through posterior conjugacy. By projecting the posterior draws onto the \( \mathbb{L}_2 \) subspace of increasing, right-continuous functions, we obtained the Isotonized Inverse Bayes Posterior (IIP) and we verified that it allows for a Bernstein-von Mises theorem with minimax asymptotic variance, reflecting the H\"{o}lder continuity of the true distribution. This proves that our methodology gives asymptotically correct uncertainty quantification under the frequentist point of view and it highlights its potential for Bayesian nonparametrics in inverse problems. 

Regarding potential future work, several interesting directions emerge. First, it would be important to verify whether the bootstrap procedure of \cite{19} remains valid for all degrees of smoothness of the underlying truth, not just for $F_0$ differentiable at $x$. Second, an open question is whether a Bernstein-von Mises (BvM) result can be obtained for Wicksell's problem within the classical Bayesian framework that is, by placing a Dirichlet Process (DP) prior $F$ on \( F_0 \). Formally obtaining a nonparametric hierarchical Bayesian model for the observations $Z_1,\ldots,Z_n$. In the following section we will illustrate the obtained results with this hierarchical Bayesian approach.

We conclude by remarking that the work of the present chapter introduces a fresh perspective on Bayesian inverse problems by employing the concept of the \textit{projection posterior}, first proposed in \cite{54}. The methods proposed here are also significantly computationally faster if compared to the standard hierarchical Bayes model illustrated above, because they do not rely on computationally expensive Gibbs samplers. Finally, a valuable direction for future research is to identify other problems where this approach could be effectively utilized.


%%%%%%%%%%%%%%%%%%%%%%%%%%%%%%%%%%%%%%%%%%%%%%
%% Single Appendix:                         %%
%%%%%%%%%%%%%%%%%%%%%%%%%%%%%%%%%%%%%%%%%%%%%%

\section[Posterior Contraction via Hierarchical Modelling]{Posterior Contraction via Dirichlet-Beta Hierarchical Modelling}\label{III.sec: paper III}
\sectionmark{Posterior Contraction via Hierarchical Modelling}

\subsection{Introduction}\label{III.sec: intro}

We closed the previous section by mentioning that one of the most interesting directions for future research is the study of Wicksell's problem within the classical Bayesian framework that is, by placing a Dirichlet Process (DP) prior $F$ on \( F_0 \). In particular, obtaining a Bernstein-von Mises result in this setting, remains an open question. However, in this section we show consistency and posterior contraction results obtained for this modelling approach.  

We perform statistical inference in the setting of Wicksell's problem using the following nonparametric hierarchical Bayesian model for the observations $Z_1,\ldots,Z_n$:
\begin{itemize}
    \item[(i)] A probability distribution $F$ on $\mathbb{R}_+$ is generated from the Dirichlet Process prior $\operatorname{DP}(\alpha)$.
    \item[(ii)] An i.i.d.\ sample $X_1,\ldots,X_n$ is generated from $F$.
    \item[(iii)] An i.i.d.\ sample $Y_1,\ldots,Y_n$ is generated from a density $\psi(y) = \frac{1}{2\sqrt{1-y}}, \, y \in (0,1).$ of a beta distribution $\operatorname{Be}(1,1/2)$, independent of the other samples $X_1,\ldots,X_n$.
    \item[(iv)] The observations are $Z_i = Y_i \cdot X_i$, for $i=1,\ldots,n$, with density as in \eqref{eq: g_F_0 expression}.
\end{itemize}
In the above model, the posterior is a Mixture of Dirichlet Processes (MDP), with distribution and mean, for any measurable $A$:
\begin{align*}
    &F \mid Z_1,\ldots,Z_n \sim \int \mathrm{DP}\bigg( \alpha + \sum_{i=1}^{n} \delta_{X_i}\bigg) \: d \mathrm{P}(X_1,\ldots,X_n \mid Z_1,\ldots,Z_n), \numberthis \label{III.eq: posterior distribution} \\
    &\mathrm{E}  [F (A) \mid Z_1,\ldots,Z_n] = \numberthis \label{III.eq: posterior mean} \\
    &= \int \bigg( \frac{|\alpha|}{| \alpha| + n} \overline{\alpha}(A) + \frac{1}{| \alpha| + n} \sum_{i=1}^{n} \delta_{X_i}(A) \bigg) \: d \mathrm{P}(X_1,\ldots,X_n \mid Z_1,\ldots,Z_n),
\end{align*}
where \eqref{III.eq: posterior distribution} follows directly from Corollary 3.1 in \cite{21} and \eqref{III.eq: posterior mean} follows from \eqref{III.eq: posterior distribution}, Fubini's theorem and Proposition 4.2 in \cite{4}.

\subsection{Main results}\label{III.sec: main results}


Given two probability distributions $G$ and $G_0$ with densities $g$ and $g_0$, the Kullback-Leiber discrepancy and corresponding variance are denoted by
$$
K\left(g_0, g\right) :=\int \log \left(g_0 / g\right) d G_0, \quad K_2\left(g_0, g\right) :=\int\left(\log \left(g_0 / g\right)\right)^2 d G_0.
$$
We write $g_0 \in \mathrm{KL}(\Pi)$ if, with positive prior probability, $\forall \: \varepsilon >0$:
\begin{align*}
    K(g_0;g) < \varepsilon.
\end{align*}
We will consider three notions of recovering the truth: the bounded Lipschitz distance $d_{\scriptscriptstyle{BL}}$ on probability distributions on $\mathbb{R}$, the Hellinger metric and the Wasserstein metric. The $d_{\scriptscriptstyle{BL}}$ distance metrizes the topology of weak convergence on the space of probability measures over a separable metric space and it is formally defined as:
\begin{align}\label{III.eq: formal definition conditional convergence}
    d_{\scriptscriptstyle{ B L}}\left(G, G_0\right):=\sup _{f \in B L_1}\left| \,\int f d G -\int f d G_0\right| ,
\end{align}
The Hellinger distance on the space of densities is given by
$$
h(g_0, g):=\left(\int\left(g_0^{1 / 2}(z)-g^{1 / 2}(z)\right)^2 d z\right)^{1 / 2}
$$
We are going to use several times that:
\begin{align}\label{III.eq: bound hellinger with L1}
h^2(g_0, g) \leq\|g_0-g\|_1 \leq 2 h(g_0, g),     
\end{align} 
The third notion of recovery consists in measuring recovery of the underlying truth via the Wasserstein metric
\begin{align*}
W_k \left( G, G_0 \right) := \inf_{\pi \in \Pi^* \left(G, G_0\right)}\left(\int |x-y|^k d \pi(x, y)\right)^{1 / k},  
\end{align*}
where $\Pi^*\left(G, G_0\right)$ is the collection of all couplings $\pi$ of $G$ and $G_0$ into a bivariate measure with marginals $G$ and $G_0$ (i.e. if $(x, y) \sim \pi$, then $x \sim G$ and $y \sim$ $G_0$), and $k \geq 1$.

\subsubsection{Consistency}\label{III.sec: consistency}

We assume that the distribution $F$ is supported inside a finite interval $[0,L]$. In practice, the endpoint will be known (e.g.\ the maximal length of the studied surface) and hence we can choose a base measure of the $\operatorname{DP}(\alpha)$ prior to have support at least $[0,L]$. In short:  
\begin{align}\label{III.eq: running assumption}
    \mathrm{supp}(F_0) \subseteq [0,L] \subset \mathrm{supp}(\alpha) \quad \text{where:} \quad F \sim \mathrm{DP}(\alpha),  
\end{align}
where $L$ is not an atom of $F_0$. 
\begin{theorem}[Consistency]\label{III.thm: consistency}
Let $F \sim  \, \, \Pi  = \mathrm{DP} (\alpha)$ be the Dirichlet process prior and $g_{\scriptscriptstyle{F}}$ as in \eqref{eq: g_F_0 expression}. Suppose that $\ex_0 \big| \log{g_0} (Z) \big| \vee \ex_0 \big| \log{F_0(\max{(Z,1)},L)} \big|  < \infty$ and that \eqref{III.eq: running assumption} holds true, then the following assertions are true: 
\begin{thmlist}
\item $g_0 \in \mathrm{KL}(\Pi)$ and the posterior distribution is consistent relative to the weak topology at $g_0$, i.e.\ $\forall \, \varepsilon >0$ as $n \rightarrow \infty$:\label{III.thm: weak topology consistency}
$$  \Pi_n(F \, : \, d_{\scriptscriptstyle{BL}} (G_{\scriptscriptstyle{F}} , G_0)> \varepsilon \mid Z_1,\ldots,Z_n) \rightarrow 0, \: \: \: \: \: \: \: G^n_0 \: \: \text{a.s.}$$
\item If $\overline{\alpha}([0,n^{-\beta}]) \leq e^{-Cn}$ for some $C >0$ and $\beta >0$, then the posterior is consistent relative to the $\mathbb{L}_1$ norm at $g_0$, i.e.\ $\forall \, \varepsilon >0$ as $n \rightarrow \infty$:\label{III.thm: L1 consistency}
$$  \Pi_n (F \, : \, \|g_{\scriptscriptstyle{F}} - g_0\|_1 > \varepsilon \mid Z_1,\ldots,Z_n ) \rightarrow 0, \: \: \: \: \: \: \: G^n_0 \: \: \text{a.s.}$$
\item If the condition of b.\ holds, then the posterior is consistent relative to the weak topology at $F_0$, i.e.\ $\forall \: \varepsilon > 0$ as $n \rightarrow \infty$:\label{III.thm: weak consistency underlying}
\begin{align*}
     &\Pi_n \left(F \, : \, d_{\scriptscriptstyle{BL}} \left( F,F_0 \right) > \varepsilon \mid Z_1,\ldots,Z_n \right) \rightarrow 0, \: \: \: \: \: \: \:  G^n_0 \: \: \text{a.s.}
\end{align*}
\item If the condition of b.\ holds, then the posterior is consistent relative to the weak topology at $F^*_0$, i.e.\ $\forall \: \varepsilon > 0$ as $n \rightarrow \infty$:\label{III.thm: weak consistency underlying F^*}
\begin{align*}
     &\Pi_n \left(F \, : \, d_{\scriptscriptstyle{BL}} \left( F^*,F^*_0 \right) > \varepsilon \mid Z_1,\ldots,Z_n \right) \rightarrow 0, \: \: \: \: \: \: \:  G^n_0 \: \: \text{a.s.}
\end{align*}
\end{thmlist} 
\end{theorem}

The proof of Theorem \ref{III.thm: consistency} is deferred to Section \ref{III.sec: proof consistency results}.





\subsubsection{Contraction rates}\label{III.sec: contraction rates}

In this section we will focus on the contraction rate of the three fundamental functionals for the application in material sciences: $g_0, V_0 $ and $F^*_0$. However we will use three different notions of recovery for the three different functionals: for the first one in Hellinger distance, for the second one pointwise convergence at a fixed point $x \geq 0$ and for the third one the Wasserstein metric. The reason for this choice is that the first one will help to prove the second one, and the first and the second will help to prove the third one. Again for mathematical convenience, we are going to make use of a slightly more stringent assumption than \eqref{III.eq: running assumption}, i.e.\ for $0<\delta<L$:
\begin{align}\label{III.eq: running assumption2}
    \mathrm{supp}(F_0) \subseteq [\delta,L] \subset \mathrm{supp}(\alpha) \quad \text{where:} \quad F \sim \mathrm{DP}(\alpha).    
\end{align}

\begin{theorem}[Contraction rates]\label{III.thm: contraction rates}
    Let $g_{\scriptscriptstyle{F}}$ as in \eqref{eq: g_F_0 expression} and $F \sim  \, \, \Pi  = \mathrm{DP} (\alpha)$ be the Dirichlet process prior with center measure $\alpha$ that has a Lebesgue density bounded away from 0 and $\infty$ on its bounded support that satisfies \eqref{III.eq: running assumption2}. Further, let $F_0$ have density $f_0$ bounded away from zero and infinity,  $\| 1/ g_0 \|_1 < \infty$ and $Z_1,\ldots,Z_n \stackrel{i.i.d.}{\sim} g_0$. Then:
    \begin{thmlist}
        \item there exists a constant $M$ such that, in $G_0$-probability: \label{III.thm: contraction rate hellinger}
        $$
            \Pi_n\Big(F: h\big(g_{\scriptscriptstyle{F}}, g_0 \big) \geq M n^{-1 / 4} (\log n)^{1/4} \mid Z_1, \ldots, Z_n\Big) \rightarrow 0,
        $$
        \item for every $x \geq 0$, there exists a constant $M$ such that, in $G_0$-probability:\label{III.thm: contraction rate V}
        \begin{align*}
            \Pi_n\left(F : |V_{\scriptscriptstyle{F}}(x) - V_0(x)| \geq M n^{-1 / 4} (\log n)^{3/4} \mid Z_1, \ldots, Z_n\right) \rightarrow 0.
        \end{align*}
        \item for every $k \geq 1$, there exists a constant $M$ such that, in $G_0$-probability:\label{III.thm: contraction rate F}
        \begin{align*}
            \Pi_n\left(F : W_k\left(F^*, F^*_0\right) \geq M  n^{-1 /(2 k + 4)}(\log n)^{\frac{4k +3}{2k+4}} \mid Z_1, \ldots, Z_n\right) \rightarrow 0.
        \end{align*}
    \end{thmlist}
     
\end{theorem}

The proof of Theorem \ref{III.thm: contraction rates} is deferred to Section \ref{III.sec: proof contraction rates}. The contraction rate $n^{-1/4}$ for the Hellinger distance is equal to the $\mathbb{L}_2$-minimax rate for estimating a density in a Sobolev space of order $1/2$ (see \cite{45}). Because the densities $g_0$ are in $H^{1/2}[0,L]$ if $F_0$ possesses a bounded density, this suggests that the rate  $n^{-1/4}$ is close to optimal. The rate with respect to the Wasserstein distance remains unknown. The pointwise contraction rate for the functional \( V_0 \) follows from the contraction rate in Hellinger distance, combined with a proof technique based on a local perturbation. This perturbation, which is believed to correspond to the efficient influence function for the functional of interest, is essential for handling arguments involving \( 1/g_{\scriptscriptstyle{F}} (z)\), which is unbounded as $z \to L$. However, this localization comes at the cost of contracting at the Hellinger distance rate. Based on the findings of \cite{51} and the results of our simulation study, we believe that this rate is suboptimal and that a Bernstein-von Mises phenomenon may hold.

\subsection{Simulation study}\label{III.sec: simulation}
 
In this section, we present simulations to illustrate the practical behavior of the presented methodology. We compare it also with the Inverse Isotonic Estimator (IIE) studied in \cite{27}. We consider two distinct settings. The first is a \textit{synthetic} setting where the true underlying distribution function \( F_0 \) is known. The first simulation setting serves as a valuable benchmark for assessing the quality of our methodology. The second is a \textit{real-data} application involving the positions of stars in the Globular Cluster Messier 62. In the second case, the ground truth distribution of the stars is unknown. 

\vspace{-0.4cm}
\begin{figure}[H]
    \centering
    \includegraphics[width=0.6\linewidth]{img_sim/comp.pdf}
    \setlength{\belowcaptionskip}{-10pt}
    \vspace{-0.1cm}
    \caption{Comparison for $n=200$ isotonic estimator and bayesian Dirichlet-based estimator.}
    \label{III.fig: example}
\end{figure}
\vspace{-0.1cm}


\begin{figure}[H]
    \centering
    \begin{subfigure}{0.55\textwidth}
        \includegraphics[width=\linewidth]{img_sim/GCcomparison3(450 samples)foto.pdf}
    \end{subfigure}
    \\
    \begin{subfigure}{0.51\textwidth}
        \includegraphics[width=\linewidth]{img_sim/gc200.pdf}
    \end{subfigure}
    \hspace{-0.6cm}
    \begin{subfigure}{0.51\textwidth}
        \includegraphics[width=\linewidth]{img_sim/gc400.pdf}
    \end{subfigure}    
    \caption{Comparison for $n=112$ and $n=450$ projected star positions between isotonic estimator and bayesian Dirichlet-based estimator.}
    \label{III.fig:stars}
\end{figure}


\begin{figure}[H]
    \centering
    \hspace{0.12cm}
    \begin{subfigure}{0.51\textwidth}
        \includegraphics[width=0.9\linewidth]{img_sim/200.pdf}
    \end{subfigure}
    \hspace{-0.7cm}
    \begin{subfigure}{0.51\textwidth}
        \includegraphics[width=0.93\linewidth]{img_sim/2000.pdf}
    \end{subfigure}
    \\
    \hspace{0.05cm}
    \begin{subfigure}{0.51\textwidth}
        \includegraphics[width=0.89\linewidth]{img_sim/hist_200.pdf}
    \end{subfigure}
    \hspace{-0.6cm}
    \begin{subfigure}{0.51\textwidth}
        \includegraphics[width=0.89\linewidth]{img_sim/hist_1000.pdf}
    \end{subfigure}
    \\
    \hspace{-0.4cm}
    \begin{subfigure}{0.52\textwidth}
        \includegraphics[width=\linewidth]{img_sim/qq_200.pdf}
        \caption*{$n=200$}
    \end{subfigure}
    \hspace{-0.7cm}
    \begin{subfigure}{0.52\textwidth}
        \includegraphics[width=\linewidth]{img_sim/qq_1000.pdf}
        \caption*{$n=1000$}
    \end{subfigure}
    \caption{Analysis normality.}
    \label{III.fig: analysis normality}
\end{figure}

In Figure \ref{III.fig: example} and \ref{III.fig: analysis normality}, we analyze the setting where the underlying cdf \( F_0 \) is exponential with rate 1.2. Samples \( Z_1, \ldots, Z_n \) used to compute the posterior are generated hierarchically: \( X_1, \ldots, X_n \sim F_0 \), \( Y_1, \ldots, Y_n \sim \mathrm{Be}(1, 1/2) \), with \( Z_i = X_i \cdot Y_i \) for all \( i \). The plots use the following color scheme: green for prior draws \( F \sim \operatorname{DP}(M \alpha) \) (parameters \( M \) and \( \alpha \) specified in the plots); purple for draws from the posterior; red for the average of these draws (approximating the posterior mean); yellow for the Isotonic Inverse Estimator (IIE) of \cite{27}; and orange for \( F_0 \).  

In Figure \ref{III.fig:stars}, we apply the methodology to star position data from globular cluster Messier 4. Let \( \mathbf{Z} = (Z_{p_1}, Z_{p_2}, Z_{p_3}) \) represent a star's three-dimensional position, where only \( (Z_{p_1}, Z_{p_2}) \) are observed. Here, \( X = Z_{p_1}^2 + Z_{p_2}^2 + Z_{p_3}^2 \sim F_0 \), and \( Z = Z_{p_1}^2 + Z_{p_2}^2 \). Beyond serving as a compelling application of our methodology, this example highlights that also in real data applications the IIE of \cite{27} behaves very closely to the present methodology. 

The first column of Figure \ref{III.fig: analysis normality} shows that with \( n = 200 \). The average of the posterior appears to be somewhat smoother, which will typically be seen as an advantage. By increasing the sample size to \( n = 1000 \) the uncertainty significantly reduces. The two bottom rows examine normality: histograms, QQ-plots, and statistical tests strongly confirm normality of the gathered samples. Furthermore, the Bayesian-estimated standard deviation aligns closely with the theoretical value attained by the frequentist procedure, suggesting that a Bernstein-von Mises phenomenon is satisfied. In all the plots we employ 300 posterior draws.

\subsection{Conclusion}\label{III.sec: conclusion}

In this section we studied Wicksell's problem within the classical nonparametric Bayesian framework, placing a Dirichlet Process prior directly on $F_0$. We established consistency of the posterior in several metrics (Theorem~\ref{III.thm: consistency}) and derived contraction rates for the key functionals $g_0$, $V_0$ and $F^*_0$ (Theorem~\ref{III.thm: contraction rates}).

The contraction rate $n^{-1/4}$ for the Hellinger distance equals the $\mathbb{L}_2$-minimax rate for estimating a density in a Sobolev space of order $1/2$ (see \cite{45}). Since $g_0 \in H^{1/2}[0,\infty)$ whenever $F_0$ possesses a bounded density, this suggests that $n^{-1/4}$ is close to optimal. The pointwise contraction rate for $V_0$ follows from the Hellinger contraction rate, combined with a Laplace transform rewriting of the posterior as in equation (2.13) of \cite{37}, a local perturbation and a LAN expansion as in \cite{27}. This perturbation, believed to correspond to the efficient influence function for the functional of interest, is essential for controlling terms involving $1/g_{\scriptscriptstyle{F}}(z)$, which is unbounded as $z \to L$. However, the local nature of the perturbation in Wicksell's problem comes at the cost of contracting at the Hellinger rate $h$ rather than at the faster rate $h^2 = O_p(\sqrt{\log n}/\sqrt{n})$ that arises in the proof style of \cite{37}. Despite significant efforts, overcoming this limitation remains an open problem.

The question of whether a Bernstein-von Mises theorem holds in this setting is also open. The simulation study presents encouraging evidence: the posterior concentrates around the truth at the expected frequentist rate, and normality of the centered and rescaled posterior draws is strongly supported by histograms, QQ-plots and formal tests. Moreover, the Bayesian standard deviation closely matches the theoretical frequentist value, suggesting that efficiency is attained. Establishing this rigorously, however, remains a challenging open problem.

From a practical standpoint, the hierarchical Bayesian approach presented here is considerably more demanding to implement than the projection-based method of the previous section. The latter requires no Markov Chain Monte Carlo: isotonized posterior draws are obtained analytically via the Greatest Convex Majorant, making computation fast and straightforward. The present approach, by contrast, requires MCMC sampling from the Mixture of Dirichlet Processes posterior \eqref{III.eq: posterior distribution}, which is computationally slow and harder to tune. Unless the full posterior distribution over $F$ is of interest, the projection approach is therefore preferable in practice.

\newpage

\section{Appendix}\label{sec: appendix}
\textit{Measure-theoretical framework}: in this chapter, the sample space is $(\mathbb{R}^+, \mathscr{B}(\mathbb{R}^+))$, and $\mathfrak{M}$ (with Borel $\sigma$-field $\mathscr{M}$ for the weak topology) is the collection of all Borel probability measures on $(\mathbb{R}^+, \mathscr{B}(\mathbb{R}^+))$. The Dirichlet process prior is a probability measure on $\mathfrak{M}$. A prior on the set of probability measures $\mathfrak{M}$ is a probability measure on a sigma-field $\mathscr{M}$ of subsets of $\mathfrak{M}$. Alternatively, it can be viewed as map from some probability space $(\Omega, \mathscr{U}, \mathrm{Pr})$ into $(\mathfrak{M}, \mathscr{M})$. We can think of the hierarchy of the spaces as follows: $(\Omega, \mathscr{U}, \mathrm{Pr}) \rightarrow (\mathfrak{M}, \mathscr{M}) \rightarrow (\mathbb{R}^+, \mathscr{B}(\mathbb{R}^+))$. We put a Dirichlet process prior over $G$ and write $G \sim \Pi = \mathrm{DP}( \alpha)$ for some base measure $\alpha$ and prior precision $|\alpha|$. This means that for any measurable set $B \in \mathscr{M}$ we have $\Pi(B) = \mathrm{Pr} ( G \in B)$. Throughout the chapter, we use the notations $\pb$ and $\ex$, which should be understood as the outer probability measure and outer expectation, respectively, whenever issues of measurability in the considered processes arise (see \cite{5} for a rigorous treatment of these topics). 

Furthermore, we use the notion of conditional convergence in distribution. This happens if the bounded Lipschitz distance between the conditional law of Borel measurable $X_n$ given $\mathcal{B}_n$ and a Borel probability measure $L$ tends to zero, where the convergence can be in outer probability. The sequence $X_n$ tends to $L$ conditionally given $\mathcal{B}_n$ in outer probability if
\begin{align}\label{eq: formal definition conditional convergence chapter3}
d_{\scriptscriptstyle{ B L}}\left(\mathcal{L}\left(X_n \mid \mathcal{B}_n\right), L\right) :=\sup _{f \in B L_1}\left| \, \mathrm{E}\left(f\left(X_n\right) \mid \mathcal{B}_n\right)-\int f d L\right| \xrightarrow{\mathrm{P}} 0,
\end{align}
where $ B L_1$ is the space of bounded Lipschitz functions with Lipschitz constant 1.

In particular if $Z_1,\ldots,Z_n \overset{\text{i.i.d.}}{\sim} \pb_{\scriptscriptstyle{Z}}$, then $X_n \mid Z_1,\ldots,Z_n \rightsquigarrow N(0,\sigma^2)$ means:
\begin{align}\label{eq: formal definition conditional convergence 1 chapter3}
	d_{\scriptscriptstyle{ B L}}\left( \pb (X_n \in \cdot \mid Z_1,\ldots,Z_n), N(0,\sigma^2) \right) \xrightarrow{\pb_{\scriptscriptstyle{Z}}} 0.
\end{align}

\subsection{Proofs of complemental statements for Semiparametric BvM}

\begin{lemma}\label{lemma: solution range equation}
    \begin{align}\label{supp:eq: solution range equation}
        &F_0(x) = 1 -\frac{2}{\pi} \sqrt{x} \, V_0(x) - \frac{2}{\pi} \int_{x}^{\infty} \frac{V_0(s)}{2 \sqrt{s}} \, ds,  
        \end{align}
    where:
    \begin{align}\label{supp:eq: V_0}
        V_0(x) \vcentcolon=\int_x^{\infty} \frac{g_0(z)}{\sqrt{z-x}} \: d z.
    \end{align}
\end{lemma}
\begin{proof}
    We will prove \eqref{supp:eq: solution range equation} using the materials science application setting, which is the same that Wicksell \cite{14} considered. In this way we will use his result. However the same formula can be obtained with similar reasonings in the astronomical application setting (c.f.\ eq.\ (1.3) and proof of Proposition 2.2 in \cite{19}).
    
    Recall that the embedded particle squared radii are i.i.d.\ realizations from a random variable $X^* \sim F^*_0$. However, when sampling, two opposing effects shape the model. First, a circular cross-section cannot exceed the radius of its parent sphere. Second, larger spheres are more likely to be intersected---sampling is proportional to radius. As a result, the observed distribution is not $F^*_0$ but its size-biased version:
    \begin{align}\label{supp:F as funct of F*}
        F_0(x)=\frac{1}{m_F} \int_0^x \sqrt{v} \, d F_0^*(v),
    \end{align}
    where $0<m_F=\int_0^{\infty} \sqrt{y} \, d F^*_0(y)<\infty$ is the expected radius under $F^*$.

    By plugging this formula into \eqref{eq: g_F_0 expression} we obtain the equivalent formula for $g_0$ as in \cite{27}: 
    $$
    g_0(z)=\frac{1}{2 m_F} \int_z^{\infty} \frac{d F^*_0(x)}{\sqrt{x-z}}.
    $$
    Wicksell \cite{14} inverted this equation, by recognising an Abel-type integral, and found an expression for $F^*_0$ in terms of $g_0$:
    \begin{align}\label{supp:eq: inv}
    F_0^*(x)=1-\frac{\int_x^{\infty}(z-x)^{-1 / 2} g_0(z) \: dz}{\int_0^{\infty} z^{-1 / 2}g_0(z) \: dz}, \quad x>0
    \end{align}
    and $F_0^*(x) = 0$ for all $x \leq 0$. Using \eqref{supp:F as funct of F*}, by defining $c_F = \int_0^{\infty} z^{-1 / 2} g_0(z) \, d z$, and integrating by parts we obtain:
    \begin{align*}
    &F_0(x) \cdot m_F=\left.\sqrt{v} F_0^*(v)\right|_0^x-\int_0^x F_0^*(v) \frac{1}{2 \sqrt{v}} \, d v=\sqrt{x} F_0^*(x)-\int_0^x F_0^*(v) \frac{1}{2 \sqrt{v}} \, d v  \\
    &=\sqrt{x}\left(1-\frac{\int_x^{\infty}(z-x)^{-1 / 2} g_0(z) d z}{\int_0^{\infty} z^{-1 / 2} g_0(z) d z}\right) -\int_0^x\left(1-\frac{\int_v^{\infty}(z-v)^{-1 / 2} g_0(z) d z}{\int_0^{\infty} z^{-1 / 2} g_0(z) d z}\right) \frac{1}{2 \sqrt{v}} \, d v\\
    &=-\sqrt{x} \frac{1}{c_F} \int_x^{\infty}(z-x)^{-1 / 2} g_0(z) \, d z + \frac{1}{c_F} \int_0^x \int_v^{\infty}(z-v)^{-1 / 2} g_0(z) \, d z \, \frac{1}{2 \sqrt{v}} \, d v,
    \end{align*}
    since:
    $\frac{2 m_F}{\pi}=(\int_0^{\infty} \frac{g_0(z)}{\sqrt{z}} d z)^{-1} \Rightarrow \frac{2 m_F}{\pi}=\frac{1}{c_F}$, therefore: $m_F c_F=\frac{\pi}{2}$, and $V_0(x) = \int_x^{\infty}(z-x)^{-1 / 2} g_0(z) \, d z$ we obtain:
    \begin{align*}
        F_0(x) = - \frac{2}{\pi} \sqrt{x} V_0(x) + \frac{2}{\pi} \int_0^x \frac{V_0(v)}{2\sqrt{v}} \, dv
    \end{align*} 
    because the density $g_0$ integrates to 1 (c.f.\ \cite{1} p.\ 1519), it follows by integration by parts:
    \begin{align*}
        1 = -\frac{2}{\pi} \int_0^{\infty} \sqrt{u} \, dV_0(u) = \frac{2}{\pi} \int_0^{\infty} \frac{V_0(u)}{2\sqrt{u}} \, du = \frac{2}{\pi} \int_0^x \frac{V_0(v)}{2\sqrt{v}} \, dv + \frac{2}{\pi} \int_x^{\infty} \frac{V_0(v)}{2\sqrt{v}} \, dv,
    \end{align*}	
    thus the claim in \eqref{supp:eq: solution range equation} follows.
\end{proof}

\begin{lemma}\label{eq: lemma equiv to claim}
	\eqref{eq: equiv to claim} implies \eqref{eq: claim}. 
\end{lemma}
	\begin{proof}
	By Theorem 1.13.1 (i) and (ii) in \cite{5} the claim \eqref{eq: claim} is obtained as soon as we show (c.f.\ \eqref{eq: formal definition conditional convergence chapter3}), for $X \sim N\left(0, \frac{g_0(x)}{2 \gamma}\right)$:
	\begin{align*}
		\sup _{f \in B L_1}\left| \, \mathrm{E}\left(f\bigg(\frac{\sqrt{n}}{\sqrt{\log n}}\big(\hat{V}_{\scriptscriptstyle{G}}(x)-\hat{V}_n(x)\big)\bigg) \mid Z_1,\ldots,Z_n \right)- \ex \left(f(X) \right) \right| \xrightarrow{G_0} 0.	
	\end{align*}
	By Lemma 2.1 given in the Online Supplement \cite{64}, for any bounded 1-Lipschitz function $f : \mathbb{R} \mapsto \mathbb{R}$ and $Z^{\scriptscriptstyle{G_0}}_n$ as in \eqref{eq: Z_n process}:
	\begin{align*}
		&\bigg|\mathrm{E} \bigg[f\bigg(\frac{\sqrt{n}}{\sqrt{\log n}}\big(\hat{V}_{\scriptscriptstyle{G}}(x)-\hat{V}_n(x)\big)\bigg) - f\bigg(\frac{\sqrt{n}}{\sqrt{\log n}}\big(\hat{V}_{\scriptscriptstyle{G}}(x)-V_0(x)\big)-Z^{\scriptscriptstyle{G_0}}_n(1)\bigg)\mid  Z_1, \ldots, Z_n  \bigg] \bigg| \\
		&   \quad \quad \leq \bigg| \frac{\sqrt{n}}{\sqrt{\log n}}\left(\hat{V}_n(x)-V_0(x)\right)-Z^{\scriptscriptstyle{G_0}}_n(1) \bigg| \xrightarrow{G_0} 0.
	\end{align*} 
	This means that the claim is obtained as soon as we show:
	\begin{align*}
		\sup _{f \in B L_1}\left| \, \mathrm{E}\left(f\bigg(\frac{\sqrt{n}}{\sqrt{\log n}}\big(\hat{V}_{\scriptscriptstyle{G}}(x)-V_0(x)\big)-Z^{\scriptscriptstyle{G_0}}_n(1)\bigg) \mid Z_1,\ldots,Z_n \right)- \ex \left(f(X) \right) \right| \xrightarrow{G_0} 0.	
	\end{align*}
	By Theorem 1.13.1 (iii) in \cite{5} the last display is implied by \eqref{eq: equiv to claim}.
	\end{proof}
	
	\begin{proof}[\hypertarget{proof of reduction}{\textbf{Proof of}} \eqref{eq: representation without remainders chapter3}]
		First note that by assumption $Qf$ is a.s.\ well-defined (c.f.\ Remark 4.4 in \cite{4}). Moreover note that \( \ex V_n = |\alpha|/(|\alpha| + n) \) and $|f_n^t(Z)| \leq 2 \sqrt{|t|/ \delta^*_n}$ and \( \forall \, \varepsilon > 0 \), by independence, that:
	\begin{align*}
	\pb \left( \delta_n^{-1} \left| V_n (Q - \mathbb{G}_n) f_t^n \right| > \varepsilon \mid Z_1, \ldots, Z_n \right) &\leq \delta_n^{-1} \varepsilon^{-1} \ex V_n \, \ex \left[  (Q + \mathbb{G}_n) | f_t^n | \mid Z_1, \ldots, Z_n \right] \\
	&\hspace{-0.2cm}\leq 4 \delta_n^{-1} \varepsilon^{-1} \frac{|\alpha|}{|\alpha| + n} \frac{\sqrt{|t|}}{\sqrt{\delta^*_n}}= O\left(\delta^{-1}_n n^{-1} (\delta^{*}_n)^{-1/2} \right).
	\end{align*}
	As $\gamma > \frac{1}{2}$, note that $\delta^{-1}_n n^{-1} (\delta^{*}_n)^{-1/2} = n^{-1} (n / \log{n})^{\frac{1}{2} + \frac{1}{4 \gamma}} \ll  n^{-1} (\sqrt{n}/ \sqrt{\log{n}})^2  \rightarrow 0$. Similarly: $\pb \left( \delta_n^{-1} \left| V_n (\mathbb{B}_n - \mathbb{G}_n) f_t^n \right| > \varepsilon \mid Z_1, \dots, Z_n \right) \stackrel{G_0}{\rightarrow} 0.$



	By Proposition G.2 in \cite{4} for \( \varepsilon_i \overset{\text{i.i.d.}}{\sim} \text{Exp}(1) \) and \( W_{ni} := \varepsilon_i / \bar{\varepsilon}_n \), we have \( (W_{n1}, \ldots, W_{nn}) \sim \text{Dir}(n; 1/n, \dots, 1/n) \). Since \( \mathbb{B}_n \sim \text{DP}(n \mathbb{G}_n) \), for $\delta_{Z_i}$ Dirac measure at $Z_i$, we have the following representation in distribution of $\mathbb{B}_n$ (c.f.\ Example 3.7.9 in \cite{5}):
	\[
	\mathbb{B}_n = \frac{1}{n} \sum_{i=1}^n W_{ni} \delta_{Z_i}.
	\]
	Because of this representation (c.f.\ section 3.7.2. in \cite{5}), $\mathbb{B}_n$ is also known as \textit{Bayesian bootstrap}. Conclude that: 
	\[
		\frac{\sqrt{n}}{\sqrt{\log n}} (\mathbb{B}_n - \mathbb{G}_n) f^n_t = \frac{1}{\sqrt{n \log n}} \sum_{i=1}^n \left( W_{ni} - 1 \right) f^n_t(Z_i).
	\]
	Because $|\bar{\varepsilon}_n - 1| = O_p(1/\sqrt{n})$, for almost every sequence $Z_1,\ldots, Z_n$ conclude that, uniformly in $t$ in compacta: 
	\begin{align}\label{eq: process representation}
		\frac{1}{\sqrt{n \log n}} \sum_{i=1}^n \left( W_{ni} - 1 \right) f_t^n(Z_i)=\frac{1}{\sqrt{n \log n}} \sum_{i=1}^n \left( \varepsilon_i - 1 \right) f_t^n(Z_i) + o_p(1),
	\end{align}
	conclude by recalling the definition of $\mathbb{Z}^*_n$ in \eqref{eq: mathbbZ process}.
	\end{proof}

    \begin{lemma}\label{supp:lemma: asymptotic linearity iso estimator}
		For fixed $x \geq 0$, $\gamma > \frac{1}{2}$ which satisfies \eqref{eq: condition roughness}, let $\hat{V}_n$ be the Inverse Isotonic Estimator (IIE) (see also Theorem 2 in \cite{27}). Then for $n \rightarrow \infty$:
		\begin{align*}
			\frac{\sqrt{n}}{\sqrt{\log{n}}} \left( \hat{V}_n(x) - V_0(x) \right)= Z^{\scriptscriptstyle{G_0}}_n(1) +o_{\scaleto{G_{0}}{4.8pt}}(1).  
		\end{align*}
		where $ Z^{\scriptscriptstyle{G_0}}_n(1)$ is given, for sequences $\delta_n = \sqrt{\log{n}}/\sqrt{n}$ and $\delta_n^{*}=\delta_n^{1 / \gamma}$ where $\gamma >1 / 2$ satisfies condition \eqref{eq: condition roughness}, 
    \begin{align}\label{supp:eq: Z_n process}
        Z^{\scriptscriptstyle{G_0}}_n(t) & :=\delta_n^{-1}\left(\delta_n^*\right)^{-1}\left\{U_n\left(x+\delta_n^* t\right)-U_n(x)-U_0\left(x+\delta_n^* t\right)+U_0(x)\right\} \\
        & =2 \delta_n^{-1} (\delta^*_n)^{-1} \int\left\{\sqrt{(z-x)_+} -\sqrt{(z-x-\delta^*_n t)_+} \right\} \, d\left(\mathbb{G}_n-G_0\right)(z).
    \end{align}
	\end{lemma}
	\begin{proof}
	    Fix $x \geq 0$ such that the conditions stated above are satisfied. Define, for $U_n$ as above and $a>0$, the process $T_n$ as follows:
		$$
		T_n(a):=\inf \left\{t \geq 0: U_n(t)-a t \text { is maximal }\right\}.
		$$
		The relation between $T_n$ and $\hat{V}_n$ is given by the switch relation (Lemma 3.2, \cite{2}):
		\begin{align}\label{supp:eq: switch relation}
		T_n(a) \leq x \quad \Leftrightarrow \quad \hat{V}_n(x) \leq a .
		\end{align}
		From this relation, it follows that for all $n$ sufficiently large:
		\begin{align*}
			\delta_n^{-1}\left(\hat{V}_n(x)-V_0(x)\right) \leq a +& Z^{\scriptscriptstyle{G_0}}_n(1) \:  \Leftrightarrow \: (\delta^*_n)^{-1}\left(T_n\left(a_0+\delta_n \left( a + Z^{\scriptscriptstyle{G_0}}_n(1) \right) \right)-x\right) \leq 0. 
		\end{align*}
		We aim to study the probability of the event on the left-hand side. We write $a_0$ for $V_0(x)$, which is fixed. Using the same steps as in the proof of Theorem 2 in \cite{27} we obtain:
		\begin{align*}
		& (\delta^*_n)^{-1}\left(T_n\left(a_0+\delta_n \left( a + Z^{\scriptscriptstyle{G_0}}_n(1) \right) \right)-x\right) \\
        &\quad = \inf \big\{t \geq-(\delta^*_n)^{-1} x: \tilde{Z}_n(t)- \left( a + Z^{\scriptscriptstyle{G_0}}_n(1) \right) t \text { is max}\big\} ,
		\end{align*}
		where:
		\begin{align*}
		\tilde{Z}_n(t)-a t :=& \: \delta_n^{-1} ( \delta^*_n)^{-1} \left\{U_n\left(x+\delta^*_n t\right)-U_n(x)- U_0\left(x+\delta^*_n t\right)+U_0(x)\right\} + \numberthis \label{supp:eq: rewrite Z tilde 1}\\ 
		&+ \delta_n^{-1} ( \delta^*_n)^{-1}  \left\{ U_0\left(x+\delta^*_n t\right)-U_0(x)-V_0(x)\delta^*_n t \right\} - at. \numberthis \label{supp:eq: rewrite Z tilde}
		\end{align*}
		Again by the proof of Theorem 2 in \cite{27} we have:
		\begin{align}
			\tilde{Z}_n(t) = Z^{\scriptscriptstyle{G_0}}_n(t) - |t|^{\gamma + 1}  (K + o(1)),
		\end{align}
		moreover: $Z^{\scriptscriptstyle{G_0}}_n(t) - t Z^{\scriptscriptstyle{G_0}}_n(1) \stackrel{ G_0}{\longrightarrow} 0$. From these facts, we conclude by applying Theorem 2.7 from \cite{12} to the process $\tilde{Z}_n$:
		\begin{align}\label{supp:eq: convergence argmax}
			(\delta^*_n)^{-1}\left(T_n\left(a_0+\delta_n \left( a + Z^{\scriptscriptstyle{G_0}}_n(1) \right)\right)-x\right) & \rightsquigarrow \underset{t \in I_x}{\arg \max }\left\{-|t|^{\gamma + 1} K  -a t \right\}. 
			\end{align}
			Since: 
			\begin{align*}
				\begin{cases}
					&\underset{t \in I_x}{\arg \max }\left\{-|t|^{\gamma + 1}  K  -a t \right\} < 0 \quad \text{if} \: a>0, \\
					&\underset{t \in I_x}{\arg \max }\left\{-|t|^{\gamma + 1}  K  -a t \right\} >0 \quad \text{if} \: a<0.
				\end{cases}
			\end{align*}
			We conclude then:
			\begin{align*}
			\mathrm{P} &\left\{\delta_n^{-1}\left(\hat{V}_n(x)- V_0(x)\right) -Z^{\scriptscriptstyle{G_0}}_n(1) \leq a \right\}  \\
			&\quad \quad \quad  \stackrel{\eqref{supp:eq: convergence argmax}}{\longrightarrow} \mathrm{P} \left\{\underset{t \in I_x}{\arg \max }\left\{-|t|^{\gamma + 1}  K  -a t \right\} \leq 0 \right\} = \begin{cases}
				&1 \quad \text{if} \: a>0, \\
				&0 \quad \text{if} \: a<0,
			\end{cases}
			\end{align*}
			which gives the claim.
	\end{proof}
	
    \begin{lemma}\label{supp:lemma: slow varying second moment}
			If 
			$g_0(x) = \int_x^\infty (4s(s-x))^{-1/2} \, dF_0(s) < \infty$, 
			then as $t \downarrow 0$:
			\begin{align}\label{supp:eq: first statement}
				\ex \left\{ \frac{1}{Z-x} \ind_{ \left\{ Z \geq x + t \right\}} \right\} \sim \log{\left(\frac{1}{t}\right)} \, g_0 (x),
			\end{align}
			where $\sim$ means that the ratio of the left-hand side and of the right-hand side goes to 1 as $t \downarrow 0$. Moreover let $Z \indep \Lambda \sim \mathrm{Exp}(1)$, then it is also true:
			\begin{align}\label{supp:eq: second statement}
				\ex \left\{ \frac{(\Lambda - 1)^2}{Z-x} \ind_{ \left\{ Z \geq x + t (\Lambda - 1)^2  \right\}} \right\} \sim \log{\left(\frac{1}{t}\right)} \, g_0 (x).
			\end{align}    
			\end{lemma}
			\begin{proof}
			We start with \eqref{supp:eq: first statement}. Consider independent $X \sim F_0$ and $V \sim \text{Be}(1,1/2)$, so that $Z \sim VX$. Consider $t \downarrow 0$, we study:
			\begin{align*}
			& \ex\left\{\frac{1}{Z-x} \ind_{\{Z>x+t\}}\right\}=\ex\left\{\frac{1}{V X - x} \ind_{\{V X>x+t\}}\right\}= \iint \frac{ 1_{\{v s>x+t\}} }{v s-x}\frac{{1}_{v \in [0,1]} }{2 \sqrt{1-v}} \, d v \, d F_0(s) \\
			& =\ex \ind_{\{X>x+t\}} \int_{\frac{x+t}{X}}^1 \frac{1}{v X -x} \frac{1}{2 \sqrt{1-v}} d v \stackrel{(\star)}{=} \ex \left\{ \int_{\frac{t}{X-x}}^1 \frac{1}{z} \frac{1}{2 \sqrt{1-z}} d z \frac{\ind_{\{X-x > t\}}}{\sqrt{X} \sqrt{X-x}} \right\} 
			\end{align*}
			where in $(\star)$ we changed variable by using $vX-x = z(X-x)$, which means $1-v = \frac{X-x}{X}(1-z)$. Define now $X_x := X-x$ and $\mathcal{J}(t) := \int_{t}^{1} \frac{1}{z} \frac{1}{\sqrt{1-z}} dz$. We will therefore show:
			\begin{align*}
				\ex\left\{ \mathcal{J}\left(\frac{t}{X_x}\right) \frac{1}{2\sqrt{X_x} \sqrt{X_x + x}} \ind_{\{X_x > t \}} \right\} \sim \log{ \left( \frac{1}{t} \right) }\ex\left\{\frac{1}{2\sqrt{X_x} \sqrt{X_x + x}} \ind_{\{X_x > 0 \}} \right\}
			\end{align*}
			In order to show the above, we first prove as $t \downarrow 0$:
			\begin{align*}
				\mathcal{J}(t) = \int_{t}^{1} \frac{1}{z} \frac{1}{\sqrt{1-z}} dz \sim \log{\frac{1}{t}}
			\end{align*}
			and $\forall \: 0 < t < 1$, it holds $\mathcal{J}(t) \lesssim \log{\left( \frac{1}{t}\right)} + 1$. 
			
			For any $b \in [t,1]$, write: $\mathcal{J}(t) = \int_{t}^{b} \frac{1}{z} \frac{1}{\sqrt{1 - z}} dz + \int_{b}^{1} \frac{1}{z} \frac{1}{\sqrt{1 - z}} dz$. For $r >1$:
			\begin{align*}
				\int_{b}^{1} \frac{1}{z} \frac{1}{\sqrt{1 - z}} dz &\leq \left(\log \frac{1}{b}\right)^r \int_{b}^{1} \frac{1}{z (\log \frac{1}{z})^r} \frac{1}{\sqrt{1 - z}} dz \\
				&\leq \left(\log \frac{1}{b} \right)^r \int_{0}^{1} \frac{1}{z (\log \frac{1}{z})^r} \frac{1}{\sqrt{1 - z}} dz \lesssim \left(\log \frac{1}{b}\right)^r.
			\end{align*}
			Furthermore:
			\begin{align*}
				\left| \int_t^b \frac{1}{z} \frac{1}{\sqrt{1 - z}} dz - \int_t^b \frac{1}{z} dz \right| 
				\leq \int_t^b \frac{1}{z} dz \left( \frac{1}{\sqrt{1 - b}} - 1 \right) \lesssim \left( \log \frac{b}{t} \right) \frac{b}{(1 - b)^{3/2}}.
			\end{align*}
			Together the above gives:
			\begin{align*}
				&\left| \frac{\int_t^1 \frac{1}{z} \frac{1}{\sqrt{1 - z}} dz}{\log \frac{b}{t}} - 1 \right| 
				= \left| \frac{\int_t^b \frac{1}{z} \frac{1}{\sqrt{1 - z}} dz + \int_b^1 \frac{1}{z} \frac{1}{\sqrt{1 - z}} dz - \log \frac{b}{t}}{\log \frac{b}{t}} \right| \\
				&= \left| \frac{\int_t^b \frac{1}{z \sqrt{1 - z}} dz - \int_t^b \frac{1}{z} dz + \int_b^1 \frac{1}{z \sqrt{1 - z}} dz}{\log \frac{b}{t}} \right| \lesssim \frac{b}{(1-b)^{3/2}} + \left| \frac{1}{\log \frac{b}{t}} \right| \left| \log \frac{b}{t} \right|^r.
			\end{align*}
			Take  $b = 1/ \sqrt{\log \frac{1}{t}}$ to see that
			\begin{align*}
				\left| \frac{\mathcal{J}(t)}{\log \frac{1}{t} - \frac{1}{2} \log \log \frac{1}{t}} - 1 \right| \lesssim \frac{1}{\sqrt{\log \frac{1}{t}}} \left( 1 + o(1) \right) + \frac{\left( \log \log \frac{1}{t} \right)^r}{\log \frac{1}{t} - \frac{1}{2} \log \log \frac{1}{t}} \longrightarrow 0
			\end{align*}
			This gives $\mathcal{J}(t) \sim \log \frac{1}{t}$. Also $\mathcal{J}(t) \lesssim \mathcal{J}(\varepsilon) \lesssim 1$, for $t \geq \varepsilon, \, \, \forall \,\varepsilon > 0$ and hence $\mathcal{J}(t) \lesssim \log \frac{1}{t} + 1$.
			
			Now in order to prove the claim we need to be able to apply the Dominated Convergence Theorem to the following:
			\begin{align*}
				 \frac{1}{\log{\frac{1}{t}} + 1}\ex\left\{ \mathcal{J}\left(\frac{t}{X_x}\right) \frac{1}{2\sqrt{X_x} \sqrt{X_x + x}} \ind_{\{X_x > t \}} \right\} 
			\end{align*}
			To this end note that since $\log{\frac{X_x}{t}} \leq \log{X_x \ind_{X_x >1}}+ \log{\frac{1}{t}}$:
			\begin{align*}
				\frac{\mathcal{J}\left( \frac{t}{X_x}\right)\ind_{X_x > t}}{\log{\frac{1}{t}} + 1}  = \frac{\mathcal{J}\left( \frac{t}{X_x}\right) \ind_{X_x > t}}{\log{\frac{X_x}{t}} + 1} \frac{\log{\frac{X_x}{t}} + 1}{\log{\frac{1}{t}} + 1}  &\leq \frac{\log{\frac{X_x}{t}} + 1}{\log{\frac{1}{t}} + 1} \ind_{X_x > t} \leq \log{X_x \ind_{X_x > 1}} + \ind_{X_x > 0}
			\end{align*}
			note that $\ex\left\{ \ind_{X_x > 0} + \left( \log{X_x} \ind_{X_x > 1} \right) \right\} < \infty$, which means we can apply the Dominated Convergence Theorem upon noticing:
			\begin{align*}
				\frac{\mathcal{J}\left( \frac{t}{X_x}\right)}{\log{\frac{1}{t}} + 1} \ind_{X_x > t} \stackrel{t \rightarrow 0}{\longrightarrow} \ind_{X_x > 0} \: \: \text{a.s.}
			\end{align*}
			We proceed now with \eqref{supp:eq: second statement}. Recall $Z \indep \Lambda \sim \mathrm{Exp}(1)$, $Z = XV$ and $X \indep V \sim \mathrm{Be}(1, 1/2)$. Then for $t \downarrow 0$ we study:
			\begin{align*}
				\ex \left[ \ind_{\{Z-x > (\Lambda - 1)^2 t\}} \frac{(\Lambda - 1)^2}{Z-x} \right] = \ex \left[ \ind_{\{X > x + (\Lambda - 1)^2 t\}} \int_{\frac{x + (\Lambda - 1)^2 t}{X}}^1  \frac{(\Lambda - 1)^2}{vX - x} \frac{1}{2 \sqrt{1 - v}} dv \right]
			\end{align*}
			using $vX - x = z(X-x)$ and $1-v = \frac{X-x}{X}(1-z)$ the above display equals:
			\begin{align*}
				&\ex \left[ \ind_{\{X > x + (\Lambda - 1)^2 t\}} \int_{\frac{(\Lambda - 1)^2 t}{X-x}}^1  \frac{(\Lambda - 1)^2}{z(X - x)} \frac{1}{2\sqrt{1-z}} \sqrt{\frac{X-x}{X}} \frac{X-x}{X} dz \right] \\
				&= \ex \left[ \ind_{\{X > x + (\Lambda - 1)^2 t\}} \mathcal{J}\left( \frac{(\Lambda - 1)^2 t}{X-x} \right) \frac{(\Lambda - 1)^2}{2\sqrt{X} \sqrt{X-x}} \right]
			\end{align*}
			where again $\mathcal{J}(t) = \int_{t}^{1} \frac{1}{z} \frac{1}{\sqrt{1-z}} dz$, which satisfies $\mathcal{J}(t) \sim \log \frac{1}{t}$ and $\mathcal{J}(t) \lesssim \log \frac{1}{t} + 1$ for $t \in (0,1)$. Thus:
			\begin{align*}
				&\frac{\mathcal{J}(\frac{(\Lambda-1)^2 t}{X-x})}{\log \frac{1}{t}} \frac{(\Lambda-1)^2}{2\sqrt{X}\sqrt{X-x}} \ind_{\{X > x + (\Lambda-1)^2 t\}} \\
				&\sim \frac{\log\left( \frac{X-x}{(\Lambda-1)^2 t} \right)}{\log(\frac{1}{t})} \frac{(\Lambda-1)^2}{2\sqrt{X}\sqrt{X-x}} \ind_{\{X > x + (\Lambda-1)^2 t\}} \overset{t \downarrow 0}{\longrightarrow} \frac{(\Lambda-1)^2}{2\sqrt{X}\sqrt{(X-x)_+}} \: \text{a.s.}
			\end{align*}
			Moreover:
			\begin{align*}
				&\left| \frac{\mathcal{J}(\frac{(\Lambda-1)^2 t}{X-x})}{\log \frac{1}{t}} \frac{(\Lambda-1)^2}{2\sqrt{X}\sqrt{X-x}} \ind_{\{X > x + (\Lambda-1)^2 t\}} \right| \\
				&\leq \frac{\left( \mathcal{J}(t) + \left(1 + \log\left(\frac{X-x}{(\Lambda-1)^2 t}\right) \ind_{\{\frac{(\Lambda-1)^2}{X-x} < 1\}}\right)\right)}{\log(\frac{1}{t})} \frac{(\Lambda-1)^2}{2\sqrt{X}\sqrt{X-x}} \ind_{\{X > x + (\Lambda-1)^2 t\}} \\
				&\lesssim \frac{(\Lambda-1)^2}{\sqrt{X}\sqrt{X-x}} \ind_{\{X>x\}} + \sup_{u>1} \frac{\log u}{\sqrt{u}} \cdot \frac{|\Lambda-1|}{\sqrt{X}} \ind_{\{X>x\}}
			\end{align*}
			Conclude again by applying the Dominated Convergence Theorem.
		\end{proof}

        \begin{lemma}\label{supp:lemma: third convergence}
			Let $\xi_{n1}, \dots, \xi_{nn}$ be a triangular array of independent random variables that satisfy,
			\(
			\sum_{i=1}^{n} \ex|\xi_{ni}| = O(1) , \) and $\forall \, \varepsilon > 0 $, \( \sum_{i=1}^{n} \ex|\xi_{ni}| \ind_{|\xi_{ni}| > \varepsilon} \to 0.
			\) as $n \rightarrow \infty$.
			Then:
			\(
			\sum_{i=1}^{n} (\xi_{ni} - \ex \xi_{ni}) \overset{\pb}{\to} 0.
			\)
		\end{lemma}
		\begin{proof}
			Fix $\varepsilon >0$
			\[
			\text{Var} \left( \sum_{i=1}^{n} \xi_{ni} \ind_{|\xi_{ni}| \le \varepsilon} \right) = \sum_{i=1}^{n} \text{Var} \left( \xi_{ni} \ind_{|\xi_{ni}| \le \varepsilon} \right) \lesssim \sum_{i=1}^{n} \varepsilon \, \ex|\xi_{ni}|  \lesssim \varepsilon
			\]
			So, $\forall \, \varepsilon_n \to 0$:
			\[
			\sum_{i=1}^{n} (\xi_{ni} \ind_{|\xi_{ni}| \le \varepsilon_n} - \ex \xi_{ni} \ind_{|\xi_{ni}| \le \varepsilon_n}) \overset{\pb}{\to} 0.
			\]
			Now note that $\forall \, \varepsilon > 0$,
			\[
			\ex \left( \left|\sum_{i=1}^{n} (\xi_{ni} - \ex \xi_{ni}) - \sum_{i=1}^{n} (\xi_{ni} \ind_{|\xi_{ni}| \le \varepsilon} - \ex \xi_{ni} \ind_{|\xi_{ni}| \le \varepsilon}) \right| \right) \le 2 \sum_{i=1}^{n} \ex |\xi_{ni}| \ind_{|\xi_{ni}| > \varepsilon} \to 0,
			\]
			hence also for some $\varepsilon_n \to 0$. 
		\end{proof}

        \begin{lemma}\label{lemma: weak conv epsilons scaled}
            For \( \varepsilon_i \overset{\text{i.i.d.}}{\sim} \text{Exp}(1) \) and $Z_i \stackrel{i.i.d.}{\sim} G_0$:
            \begin{align}\label{supp:eq: weak conv for NBP}
                \frac{1}{\sqrt{n \log{n}}} \sum_{i=1}^n \frac{(\varepsilon_i -1)}{\sqrt{Z_i - x}} \ind_{\{ Z_i >x \}} \mid Z_1, \ldots, Z_n \rightsquigarrow N\left(0, g_0(x)\right) \, \, \text { under } \, G_0 \text {. }
            \end{align}
        \end{lemma}

        \begin{proof}
		We apply Theorem 6.16 in \cite{53} (with $\xi_{ni} := \frac{1}{\sqrt{n \log{n}}} \frac{\varepsilon_i -1}{\sqrt{Z_i - x}} \ind_{\{Z_i > x \}}$) conditionally on $Z_1,\ldots,Z_n$. Therefore we need to show: 
		\begin{align*}
			\sum_{i=1}^n \pb \left(\left. \left(\varepsilon_i-1\right)^2 \geq \delta \, n\log{n} \, |Z_i-x|, \, Z_i > x \, \,  \right\rvert\, Z_1, \ldots, Z_n\right) \stackrel{G_0}{\rightarrow} 0, \quad \forall \, \delta >0, \numberthis \label{supp:eq: pb convergence} \\
			\frac{1}{\sqrt{n \log{n}}} \sum_{i=1}^n \ex \left(\left. \frac{(\varepsilon_i-1)}{\sqrt{Z_i - x}} \ind_{\{\left(\varepsilon_i-1\right)^2 <  \, n\log{n} \, |Z_i-x|, \, Z_i > x \}} \,  \right\rvert\, Z_1, \ldots, Z_n\right) \stackrel{G_0}{\rightarrow} 0, \numberthis \label{supp:eq: first moment convergence}
		\end{align*}
		and:
		\begin{align}\label{supp:eq: var convergence}
			\frac{1}{n\log{n}}\sum_{i=1}^n \operatorname{Var}\left(\left.\frac{\left(\varepsilon_i-1\right)}{\sqrt{Z_i-x}} \ind_{\{\left(\varepsilon_i-1\right)^2 <  \, n\log{n} \, |Z_i-x|, \, Z_i > x \}}  \, \right\rvert\, Z_1, \ldots, Z_n\right) \stackrel{G_0}{\rightarrow} g_0(x).
		\end{align}
		First note that \eqref{supp:eq: pb convergence} and \eqref{supp:eq: first moment convergence} are implied by convergence in mean (which means unconditional probability and expectation). So \eqref{supp:eq: pb convergence} is implied by:
		\begin{align*}
			n \, \pb \left( \left(\varepsilon-1\right)^2 \geq \delta \, n\log{n} \, |Z-x|, \, Z > x \right) &= n \int_{0}^{\infty} \int_{x}^{x + \frac{(\lambda - 1)^2}{\delta n \log{n}}} e^{- \lambda} g_0(z) \, dz \, d \lambda \\
			&\leq n \, \|g_0 \|_{\infty} \frac{1}{\delta n \log{n} } \int_{0}^{\infty} (\lambda-1)^2 \, e^{-\lambda} \, d \lambda \rightarrow 0.
		\end{align*}
		Because $\ex (\varepsilon_i -1) = 0, \, \forall \, i$, the expected value of \eqref{supp:eq: first moment convergence} goes to zero because:
		\begin{align*}
			&\ex \, \bigg| \frac{1}{\sqrt{n \log{n}}} \sum_{i=1}^n \ex \left( \frac{(\varepsilon_i-1)}{\sqrt{Z_i - x}} \ind_{\{\left(\varepsilon_i-1\right)^2 \geq  \, n\log{n} \, |Z_i-x|, \, Z_i > x \}} \, \bigg| \, Z_1,\ldots,Z_n \right) \bigg| \\
			&\leq \frac{\sqrt{n}}{\sqrt{\log{n}}} \int_{0}^{\infty} \int_{x}^{x + \frac{(\lambda -1)^2}{n \log{n}}} \frac{|\lambda -1|}{\sqrt{z-x}} e^{- \lambda} g_0(z) \, dz \, d \lambda \\
			& \leq \frac{\sqrt{n}}{\sqrt{\log{n}}} \| g_0\|_{\infty} \int_{0}^{\infty} \sqrt{\frac{(\lambda -1)^2}{n \log{n}}} |\lambda - 1| e^{- \lambda} \, d \lambda \rightarrow 0.
		\end{align*}
 		In order to prove \eqref{supp:eq: var convergence} we proceed in steps. First note that \eqref{supp:eq: var convergence} equals:
		\begin{align*}
			\frac{1}{n \log{n}} \sum_{i=1}^n & \bigg\{  \ex \left(\left. \frac{(\varepsilon_i-1)^2}{Z_i - x} \ind_{\{\left(\varepsilon_i-1\right)^2 <  \, n\log{n} \, |Z_i-x|, \, Z_i > x \}} \,  \right\rvert\, Z_1, \ldots, Z_n\right)  \numberthis \label{supp:eq: first line var}\\
			& - \left( \ex \left(\left. \frac{(\varepsilon_i-1)}{\sqrt{Z_i - x}} \ind_{\{\left(\varepsilon_i-1\right)^2 <  \, n\log{n} \, |Z_i-x|, \, Z_i > x \}} \,  \right\rvert\, Z_1, \ldots, Z_n\right) \right)^2 \bigg\}. \numberthis \label{supp:eq: second line var}
 		\end{align*}
		First we show the term in \eqref{supp:eq: second line var} goes to zero. Define:
		\begin{align*}
		& \mathcal{J}(u)  := \ex \left( \varepsilon - 1 \right) \mathbf{1}_{\{ |\varepsilon - 1| < u \}}  \\
		&= 
		\begin{cases} 
			\ex \left( \varepsilon - 1 \right) \mathbf{1}_{\{ |\varepsilon - 1| > u \}} = \int_{1+u}^{\infty} (\lambda - 1) e^{-\lambda} \, d\lambda \lesssim e^{-u/2}, \, \,  u > 1 & \\[10pt]
		\int_{1-u}^{1+u} (\lambda - 1) e^{-\lambda} \, d\lambda \leq \int_{1-u}^{1+u} (\lambda - 1) \, d \lambda = u^2, \, \, \,  \quad  0 \leq u \leq 1 &
		\end{cases}\lesssim u^2 \wedge e^{-u/2}, \quad \text{all } u.
		\end{align*}
		Therefore for \eqref{supp:eq: second line var} we have:
		\begin{align*}
			&\ex \frac{1}{n \log n} \sum_{i=1}^n \left( \ex \left( \frac{\varepsilon_i -1}{\sqrt{Z_i - x}} \mathbf{1}_{\{ |\varepsilon_i -1|^2 < n \log n \mid Z_i - x|, \, Z_i \geq x \}} \, \bigg| \, Z_1,\ldots,Z_n \right) \right)^2 \\
			&\leq \frac{1}{\log n} \ex \left( \frac{\ind_{Z > x}}{Z-x} \mathcal{J}^2 \left( \sqrt{n \log n  \, (Z-x)_+}  \right) \right) \\
			&\lesssim \frac{1}{\log n} \ex \left( \frac{\ind_{Z > x}}{Z-x} \left( (Z-x) \, n \log n \wedge e^{-\sqrt{(Z-x) \, n \log n}/2} \right)^2 \right) \\
			&\leq \frac{1}{\log n} \ex \left( \mathbf{1}_{Z > x} \left( (Z-x) (n \log n)^2 \mathbf{1}_{\{ Z-x < \frac{\log \log n}{n \log n} \}} \right)  + \frac{1}{Z-x} e^{-\sqrt{\log \log n}} \mathbf{1}_{\{ Z-x > \frac{\log \log n}{n \log n} \}} \right) \\
			&\leq \frac{(n \log n)^2}{\log n} \int_{x}^{x + \frac{\log \log n}{n \log n}} (z-x) g_0(z) \, dz 
			+ \frac{e^{-\sqrt{\log \log n}}}{\log n} \ex \frac{1}{Z-x} \mathbf{1}_{\{ Z-x > \frac{\log \log n}{n \log n} \}} \\
			&\leq \| g_0 \|_{\infty} \frac{1}{2} \left( \frac{\log \log n}{\log n} \right)^2 
			+ \frac{e^{-\sqrt{\log \log n}}}{\log n} \log \left( \frac{n \log n}{\log \log n} \right) 
			\left( g_0(x) + O(1) \right).
		\end{align*}
		where in the last line we used Lemma \ref{supp:lemma: slow varying second moment}. Conclude that the last display goes to 0. Now we proceed with \eqref{supp:eq: first line var}. 
		\begin{align*}
			\eqref{supp:eq: first line var} &= \frac{1}{n \log n} \sum_{i=1}^n \ex \left( \frac{(\varepsilon_i - 1)^2}{Z_i - x} \ind_{\{|\varepsilon_i - 1|^2 < n \log n |Z_i - x|, \, Z_i > x\}} \, \bigg| \, Z_1,\ldots,Z_n \right) \\
			&= \frac{1}{n \log n} \sum_{i=1}^n \frac{\ind_{\{Z_i > x\}}}{Z_i - x} T(\sqrt{n \log n} \sqrt{Z_i - x})
			\end{align*}
			where $T(u) := \ex (\varepsilon - 1)^2 \ind_{|\varepsilon - 1| < u} \asymp u^3 \land 1$. Define
			\[
			\xi_{n,i} := \frac{\ind_{\{Z_i > x\}}}{Z_i - x} \cdot \frac{T(\sqrt{n \log n} \sqrt{Z_i - x})}{n \log n}
			\]
			note $0 \leq \xi_{n,i} \leq 1$, because $\frac{T(u)}{u^2} \leq 1$. Now we verify the conditions of Lemma \ref{supp:lemma: third convergence} and:
			\begin{align}\label{supp:eq: asymptotic variance}
				\sum_{i=1}^n \ex \xi_{ni} \rightarrow g_0(x).
			\end{align}
			First note that \eqref{supp:eq: asymptotic variance} and the fact that $\xi_{ni} \geq 0$, $\forall \, i,n$ imply the fact that $\sum_{i=1}^n \ex |\xi_{ni}| = O(1)$. Now:
			\begin{align*}
				\sum_{i=1}^{n} \ex|\xi_{ni}| \ind_{|\xi_{ni}| > \varepsilon} &\leq \sum_{i=1}^{n} \pb \left( |\xi_{ni}| > \varepsilon \right) \leq n \, \pb \left( \frac{1}{n\log{n}}\frac{\ind_{Z>x}}{Z-x} > \varepsilon \right) \\
				&\leq n \, \pb \left( 0 < Z-x < \frac{1}{\varepsilon n \log{n}} \right) \leq n \, \| g_0 \|_{\infty} \frac{1}{\varepsilon n \log{n}} \rightarrow 0.
			\end{align*}
			Finally note that the left hand side of \eqref{supp:eq: asymptotic variance} is equal to:
			\begin{align*}
				 \frac{1}{\log{n}} \ex \frac{(\varepsilon - 1)^2}{Z-x}  \ind_{\{|\varepsilon - 1|^2 < n \log n |Z - x|, \, Z > x\}} \sim \frac{\log{(n \log{n})}}{\log{n}} g_0(x).
			\end{align*}
			where in the last line we used Lemma \ref{supp:lemma: slow varying second moment}. 
	\end{proof}
    
    \begin{lemma}\label{lemma: result with cov structure}
        For \( \varepsilon_i \overset{\text{i.i.d.}}{\sim} \text{Exp}(1) \), $Z_i \stackrel{i.i.d.}{\sim} G_0$, the process $\mathbb{Z}^*_n$ is defined as:
        \begin{align}\label{supp:eq: mathbbZ process}
            \mathbb{Z}^*_n (t) := \frac{1}{\sqrt{n \log n}} \sum_{i=1}^n \left( \varepsilon_i - 1 \right) f_t^n(Z_i).
        \end{align}
        where $f_t^n(z) := 2(\delta_n^*)^{-1} \big\{ \sqrt{(z-x)_+} -\sqrt{(z-x-\delta^*_n t)_+}  \big\}$, $\delta_n = \sqrt{\log{n}}/\sqrt{n}$ and $\delta_n^{*}=\delta_n^{1 / \gamma}$ where $\gamma >1 / 2$ satisfies condition \eqref{eq: condition roughness}.
        
        Fix $t \in I_x$, where $I_x := [0, \infty) \text{ if } x = 0, \; \text{and} \; I_x := \mathbb{R} \text{ if } x > 0$. Then, in probability under $G_0$:
        \begin{align}\label{supp:eq: weak conv for IIP}
            \frac{1}{\sqrt{n \log n}} \sum_{i=1}^n \left( \varepsilon_i - 1 \right) f_t^n(Z_i) \mid Z_1, \ldots, Z_n \rightsquigarrow N\left(0, t^2  g_0(x)/ (2 \gamma) \right),
        \end{align}
        Furthermore:
        \begin{align}\label{supp:eq: cov structure}
            \text{Cov}\left(\mathbb{Z}_n^\ast(t), \mathbb{Z}_n^\ast(s) \,\middle|\, Z_1, \dots, Z_n\right) 
            = \frac{1}{2 \gamma} g_0(x) st \left( 1 - O_{\scriptscriptstyle{G_0}}\left(\frac{\log \log n}{\log n} \right) \right) + O_{\scriptscriptstyle{G_0}}\left(\frac{1}{\log n}\right).
        \end{align}
    \end{lemma}

    \begin{proof}
		Without loss of generality set $t=1$. We apply Theorem 6.16 in \cite{53} (with $\xi_{ni} = ((\varepsilon_i -1) f_1^n(Z_i)) / \sqrt{n \log{n}}  $) conditionally on $Z_1,\ldots,Z_n$. Therefore we need to show: 
		\begin{align*}
			\sum_{i=1}^n \pb \left(\left. \left|\varepsilon_i-1\right| f^n_1(Z_i) \geq  \sqrt{ \delta \, n\log{n}} \, \,\,  \right\rvert\, Z_1, \ldots, Z_n\right) \stackrel{G_0}{\rightarrow} 0, \quad \forall \, \delta >0, \numberthis \label{supp:eq: pb convergence 1} \\
			\frac{1}{\sqrt{n \log{n}}} \sum_{i=1}^n \ex \left(\left. (\varepsilon_i-1) f_1^n(Z_i) \ind_{\{\left|\varepsilon_i-1\right| f^n_1(Z_i) \geq  \sqrt{ \delta \, n\log{n}} \} } \,  \right\rvert\, Z_1, \ldots, Z_n\right) \stackrel{G_0}{\rightarrow} 0, \numberthis \label{supp:eq: first moment convergence 1}
		\end{align*}
		and:
		\begin{align}\label{supp:eq: var convergence 1}
			\frac{1}{n\log{n}}\sum_{i=1}^n \operatorname{Var}\left(\left. (\varepsilon_i-1) f_1^n(Z_i) \ind_{\{\left|\varepsilon_i-1\right| f^n_1(Z_i) \geq  \sqrt{ \delta \, n\log{n}} \} }  \, \right\rvert\, Z_1, \ldots, Z_n\right) \stackrel{G_0}{\rightarrow} \frac{g_0(x)}{2 \gamma}.
		\end{align}
		First note that \eqref{supp:eq: pb convergence 1} and \eqref{supp:eq: first moment convergence 1} are implied by convergence in mean (which means unconditional probability and expectation). In what follows, we will often use that $|f_n^1(Z)| \leq 2/ \sqrt{\delta^*_n}$. So \eqref{supp:eq: pb convergence 1} is implied by:
		\begin{align*}
			& n \, \pb\left( (f_n^1(Z))^2 > \frac{\delta \, n \log n}{(\varepsilon - 1)^2} \right) \lesssim n \, \pb\left( \frac{1}{\delta_n^*} > \frac{\delta \, n \log n}{(\varepsilon - 1)^2} \right) = n \, \pb\left( |\varepsilon - 1| > \sqrt{\delta \, \delta_n^* \, n \log n } \right) \\
			&  = n \int_{|\lambda - 1| > \sqrt{\delta_n^* \, n \log n \, \delta}} e^{-\lambda} d\lambda = n \int_{\lambda > 1 + \sqrt{\delta_n^* \, n \log n \, \delta}} e^{-\lambda} d\lambda = n \int_{1 + \sqrt{\delta_n^* \, n \log n \, \delta}}^{\infty} e^{-\lambda} d\lambda \\
			& = n \, e^{-1 - \sqrt{\delta_n^* \, n \log n \, \delta}} \to 0,
		\end{align*}
		because $\left( \frac{\log n}{n} \right)^{1/2\gamma} \cdot n \log n \gg n^{1 - 1/2\gamma} \log^2 n$, since $\delta_n^* = \left( \frac{\log n}{n} \right)^{1/2\gamma}, \quad \gamma > 1/2$. 
		
		Because $\ex(\varepsilon_i - 1) = 0$ and $\ex f_n^1(Z) = O(1)$ (see (A9) in \cite{27}), the expected value of \eqref{supp:eq: first moment convergence 1} goes to zero because:
		\begin{align*}
			&\ex \left| \frac{1}{\sqrt{n \log n}} \sum_{i=1}^n \ex \bigg( (\varepsilon_i - 1) f_n^1(Z_i) \ind_{\{|\varepsilon_i - 1| f^n_1(Z_i) > \sqrt{\delta \, n \log n}\}}  \bigg| Z_1, \ldots, Z_n \bigg) \right| \\
			& \lesssim \sqrt{\frac{n}{\log n}} \ex f_n^1(Z) \ex ( \varepsilon - 1 ) \ind_{\{|\varepsilon - 1| > \sqrt{\delta \, \delta_n^* \, n \log n}\}}  \lesssim \sqrt{\frac{n}{\log n}} \frac{( 1 + \sqrt{\delta_n^* n \log n})}{e^{\sqrt{\delta \, \delta_n^* \, n \log n} + 1}} ,
		\end{align*}
		which goes to zero as $\gamma > 1/2$ and thus $\left( \frac{\log n}{n} \right)^{1/2\gamma} \cdot n \log n \gg n^{1 - 1/2\gamma} \log^2 n \gtrsim n^{\eta}$ some $\eta>0$.

		In order to prove \eqref{supp:eq: var convergence 1} we proceed in steps. First note that \eqref{supp:eq: var convergence 1} equals:
		\begin{align*}
			\frac{1}{n \log{n}} \sum_{i=1}^n & \bigg\{  \ex \left(\left. (f^n_1(Z_i))^2 (\varepsilon_i-1)^2 \ind_{\{|\varepsilon_i - 1| f^n_1(Z_i) <  \sqrt{\delta \, n\log{n}}  \}} \,  \right\rvert\, Z_1, \ldots, Z_n\right)  \numberthis \label{supp:eq: first line var 1}\\
			& - \left( \ex \left(\left. f^n_1(Z_i) (\varepsilon_i-1) \ind_{\{|\varepsilon_i - 1| f^n_1(Z_i) <  \sqrt{\delta \, n\log{n}}  \}} \,  \right\rvert\, Z_1, \ldots, Z_n\right) \right)^2 \bigg\}. \numberthis \label{supp:eq: second line var 1}
 		\end{align*}
		First we show the term in \eqref{supp:eq: second line var 1} goes to zero. Because $\ex(\varepsilon -1) = 0$ and $\varepsilon \indep Z$:
		\begin{align*}
			&\ex \left( \frac{1}{n \log n} \sum_{i=1}^n  \ex \left( (\varepsilon_i - 1) f_n^1(Z_i) \ind_{\{|\varepsilon_i - 1| f_n^1(Z_i) > \sqrt{\delta \, n\log{n}}\}} \, \bigg| \, Z_1,\ldots,Z_n\right) \right)^2 \\
			& \leq \frac{1}{\log n} \ex (\varepsilon - 1)^2 (f_n^1(Z))^2 \ind_{\{|\varepsilon - 1| f_n^1(Z) > \sqrt{\delta \, n\log{n}}\}} \\
			& \leq \frac{1}{\log n} \ex (\varepsilon - 1)^2 (f_n^1(Z))^2 \ind_{\{|\varepsilon - 1| > \sqrt{\delta \, \delta_n^* \, n \log n}\}} \\
			& = \frac{1}{\log n} \ex (f_n^1(Z))^2 \ex (\varepsilon - 1)^2 \ind_{\{|\varepsilon - 1| > \sqrt{\delta \, \delta_n^* \, n \log n}\}} \to 0,
		\end{align*}
		because by (A7) in \cite{27} with $s=t=1$, $\frac{1}{\log n} \ex (f_n^1(Z))^2 = O(1)$ and for all $n$ big,
		\begin{align*}
		 &\ex (\varepsilon - 1)^2 \ind_{\{|\varepsilon - 1| > \sqrt{\delta \, \delta_n^* n \log n}\}} \lesssim \int_{1 + \sqrt{\delta \, \delta_n^* n \log n}}^{\infty} (\lambda - 1)^2 e^{-\lambda} d\lambda \numberthis \label{supp:eq: convergence usefull}\\
		 & = \left( \sqrt{\delta \, \delta_n^* n \log n}(\sqrt{\delta \, \delta_n^* n \log n} + 2) + 2 \right) e^{-\sqrt{\delta \, \delta_n^* n \log n} - 1} \to 0
		\end{align*}
		again because $\delta_n^* \, n \log n \sim n^{1 - 1/2\gamma} \log^{1 + \frac{1}{2\gamma}}(n)$, $\gamma > 1/2$.

		Now we show the behavior of \eqref{supp:eq: first line var 1}. Denote in what follows by $\ex_{\varepsilon}$ the expectation only under $\varepsilon$, thus conditionally on $Z_1,\ldots,Z_n$. Using that $\ex(\varepsilon -1)^2 =1$ and $Z \indep \varepsilon$:
		\begin{align*}
			\eqref{supp:eq: first line var 1} & = \frac{1}{n \log n} \sum_{i=1}^n \ex_{\varepsilon} (\varepsilon_i - 1)^2 (f_n^1(Z_i))^2 \ind_{\{|\varepsilon_i - 1| f_n^1(Z_i) < \sqrt{\delta \, n \log n}\}} \\
			& = \frac{1}{n \log n} \sum_{i=1}^n \ex_{\varepsilon} (\varepsilon_i - 1)^2 (f_n^1(Z_i))^2 \{1 - \ind_{\{|\varepsilon_i - 1| f_n^1(Z_i) \geq  \sqrt{\delta \, n \log n}\}} \} \\
			& = \frac{1}{n \log n} \sum_{i=1}^n (f_n^1(Z_i))^2 - \frac{1}{n \log n} \sum_{i=1}^n \ex_{\varepsilon} (\varepsilon_i - 1)^2 (f_n^1(Z_i))^2 \ind_{\{|\varepsilon_i - 1| f_n^1(Z_i) > \sqrt{\delta \, n \log n}\}}.
		\end{align*}
		First we show $\frac{1}{n \log n} \sum_{i=1}^n (f_n^1(Z_i))^2 \overset{G_0}{\to} g_0(x)/2\gamma$. We use again Lemma \ref{supp:lemma: third convergence} with $\xi_{ni} = (f_n^1(Z_i))^2/n\log n$. Now $\forall \, \varepsilon > 0$ and all $n$ big enough:
		\begin{align*}
		\sum_{i=1}^n \ex |\xi_{ni}| \ind_{\{|\xi_{ni}| > \varepsilon\}} &= \frac{1}{n\log n} \sum_{i=1}^n \ex (f_n^1(Z_i))^2 \ind_{\{(f_n^1(Z_i))^2 > \varepsilon n\log n\}} \\
		& \leq \frac{1}{n\log n} \sum_{i=1}^n \ex (f_n^1(Z_i))^2 \ind_{\{\frac{1}{\delta_n^*} > \varepsilon n\log n\}} = 0,
		\end{align*}
		because $\frac{1}{\delta_n^*} = (\frac{n}{\log n})^{1/2\gamma} \ll n \log n$ as $\gamma > 1/2$. Finally:
		\[
		\sum_{i=1}^n \ex \xi_{ni} = \frac{1}{\log n} \int (f_n^1(z))^2 g_0(z) dz \to \frac{1}{2\gamma} g_0(x)
		\]
		where we use (A7) of \cite{27} with $s=t=1$. The fact that $\xi_{ni}$ are nonnegative and $\sum_{i=1}^n \ex \xi_{ni} \to g_0(x)/2\gamma$ imply $\sum_{i=1}^n \ex |\xi_{ni}| = O(1)$. Using \eqref{supp:eq: convergence usefull} we conclude:
		\begin{align*}
			& \frac{1}{n \log n} \sum_{i=1}^n \ex_{\varepsilon} (\varepsilon_i - 1)^2 (f_n^1(Z_i))^2 \ind_{\{|\varepsilon_i - 1| f_n^1(Z_i) > \sqrt{\delta \, n \log n}\}} \\
			& \leq \frac{1}{n \log n} \sum_{i=1}^n (f_n^1(Z_i))^2 \ex_{\varepsilon} (\varepsilon - 1)^2 \ind_{\{|\varepsilon - 1| > \sqrt{\delta \, \delta_n^* n \log n}\}} \to 0
		\end{align*}
		since $\frac{1}{n \log n} \sum_{i=1}^n (f_n^1(Z_i))^2 = O_{\scriptscriptstyle{G_0}}(1)$. 

        We now prove \eqref{supp:eq: cov structure}. Fix \(s, t \in I_x\). Since:
		\begin{align}\label{supp:eq: cond expectation}
			\ex \left( \mathbb{Z}^*_n(t) \mid Z_1, \dots, Z_n \right)= \frac{1}{\sqrt{n\log{n}}} \sum_{i=1}^n \ex \left( \left(\varepsilon_i - 1\right) f_t^n(Z_i) \mid Z_1, \dots, Z_n \right) = 0.
		\end{align}
		Thus using \eqref{supp:eq: cond expectation}, the above relations, that $Z \indep \varepsilon$ and that $\varepsilon \overset{\text{i.i.d.}}{\sim} \mathrm{Exp}(1)$:
		\begin{align}\label{supp:eq: to determine behav}
		&\text{Cov}\left(\mathbb{Z}_n^\ast(t), \mathbb{Z}_n^\ast(s) \,\middle|\, Z_1, \dots, Z_n\right) = \frac{1}{n\log{n}} \sum_{i=1}^n f_t^n(Z_i) f_s^n(Z_i) \underbrace{\ex \left( \varepsilon_i - 1\right)^2}_{=1}.     
		\end{align}
        Now we apply Lemma \ref{supp:lemma: third convergence} with $\xi_{ni} := \frac{f_t^n(Z_i) f_s^n(Z_i)}{n \log{n}}$ to show:
		\begin{align*}
			\frac{1}{n \log{n}} \left\{ \sum_{i=1}^n f_t^n(Z_i) f_s^n(Z_i) - n \,  \ex \, f_t^n(Z) f_s^n(Z) \right\} \stackrel{G_0}{\rightarrow} 0.
		\end{align*}
		First, by using (A7) from \cite{27} and that $\delta^*_n = (\log{n}/n)^{1/2\gamma}$:
		\begin{align*}
			\frac{1}{\log{n}} \ex \, f_t^n(Z_i) f_s^n(Z_i) &= - st g_0(x) \frac{\log{\delta^*_n}}{\log{n}} +O\left(\frac{1}{\log{n}} \right) \\
                & = \frac{1}{2 \gamma} g_0(x) st \left( 1 - O\left(\frac{\log \log n}{\log n} \right) \right) + O\left(\frac{1}{\log n}\right). \numberthis \label{supp:eq: determines behav}
		\end{align*}
        This implies together with the nonnegativity of the functions $f_s^n$ that $\sum_{i=1}^n \ex |\xi_{ni}| = O(1)$. Moreover:
        \begin{align*}
            \sum_{i=1}^n \ex \bigg|\frac{f_t^n(Z_i) f_s^n(Z_i)}{n \log{n}} \bigg| \ind_{\big| \frac{f_t^n(Z_i) f_s^n(Z_i)}{n \log{n}} \big| > \varepsilon} \leq \sum_{i=1}^n \ex \bigg|\frac{f_t^n(Z_i) f_s^n(Z_i)}{n \log{n}} \bigg| \ind_{\big| \frac{\sqrt{|st|}}{\delta^*_n} \big| > \varepsilon n \log{n}} = 0,
        \end{align*}
        for all $n$ big enough, because $(n/ \log{n})^{1/2\gamma} \ll n \log{n}$ as $\gamma > 1/2$. Conclude that the behavior of \eqref{supp:eq: to determine behav} is given by \eqref{supp:eq: determines behav}.
	\end{proof}

	\begin{lemma}\label{lemma: cond asymp equicontinuity}
	The process $\mathbb{Z}_n^*$ is conditionally asymptotically equicontinuous, i.e.\ $\forall \, \eta >0$ as $n \rightarrow \infty$ followed by $\delta >0$:
	\begin{align*}
		\pb \bigg\{ \sup _{|s-t|< \delta} \big| \mathbb{Z}_n^*(s) - \mathbb{Z}_n^*(t) \big|>\eta \mid Z_1, \ldots, Z_n\bigg\} \stackrel{G_0}{\longrightarrow} 0.
	\end{align*}
	\end{lemma}
	\begin{proof}
		Take the class of functions for $\delta >0, \: M>0$:  
		\begin{align*}
		\mathscr{F}_{n, \delta}^M :=\big\{ (\log n)^{-1/2} (\varepsilon -1) (f_s^n - f_t^n)(z) : \: s, t \in I_x,|s-t|<\delta, \max (|s|,|t|)<M\big\}.  \numberthis \label{eq: class of functions chapter3} 
		\end{align*}
		and the class of functions $\mathscr{H}_{n, \delta}^M$, whose functions for $f \in \mathscr{F}_{n, \delta}^M$ are given by $f/(\varepsilon-1)$, thus without the $\varepsilon-1$. These classes have square integrable natural envelopes $F_{n,\delta}$ and $H_{n,\delta}$ that satisfy for all $\delta>0$ small (cf.\ Lemma 5 in \cite{27}):
		\begin{align*}
			\| F_{n,\delta} \|_2^2 = \ex (\varepsilon-1)^2  \| H_{n,\delta} \|_{2,\scriptscriptstyle{G_0}}^2 =\| H_{n,\delta} \|_{2,\scriptscriptstyle{G_0}}^2  \leq g_0(x + \delta^*_n s \wedge t) \, \delta^2 
		(\log{n})^{-1} ( \log{(\delta \delta^*_n)^{-2}}).
		\end{align*}
		For fixed $s, t \in I_x$, we prove  conditional asymptotic equicontinuity of $\mathbb{Z}_n^*$, i.e.\ $\forall \, \eta >0$ as $n \rightarrow \infty$ followed by $\delta >0$, the right-hand side of the display below converges to 0 under $G_0$:
		\begin{align*}
			\pb \bigg\{ \sup _{|s-t|< \delta} \big| \mathbb{Z}_n^*(s) - \mathbb{Z}_n^*(t) \big|>\eta \mid Z_1, \ldots, Z_n\bigg\} \leq \frac{1}{\eta} \ex \bigg[ \sup _{|s-t|< \delta} \big| \mathbb{Z}_n^*(s) - \mathbb{Z}_n^*(t) \big| \mid Z_1, \ldots, Z_n\bigg].
		\end{align*}
		By the maximal inequality in Theorem 2.14.1 in \cite{5}:
		\begin{align*}
			\ex \bigg[ \sup _{|s-t|< \delta} \big| \mathbb{Z}_n^*(s) - \mathbb{Z}_n^*(t) \big| \bigg] \lesssim J(1) \| F_{n,\delta} \|_2  \lesssim \sqrt{\|g_0\|_{\infty} \, \delta^2 
		(\log{n})^{-1} ( \log{(\delta \delta^*_n)^{-2}})} \to 0,
		\end{align*}
		where $J(1) := J(1,\mathscr{F}_{n, \delta}^M, \mathbb{L}_2 ) = \sup_{Q} \int_{0}^{1} \sqrt{\log N(\eta \| F_{n,\delta} \|_2,\mathscr{F}_{n, \delta}^M, \mathbb{L}_2(Q))} \, d \eta$ and where the square $\mathbb{L}_2(Q)$-distance is given for $s,t$ and $s^\prime, t^\prime \in I_x$ by:
		$
		\frac{1}{\log n} \int ((f^n_s - f^n_t)(z) - (f^n_{s^\prime} - f^n_{t^\prime})(z))^2 \, d Q(z).
		$
		(Note that we did not take $\varepsilon-1$ into account as $\ex (\varepsilon-1)^2=1$).
		Furthermore, we used that $J(1,\mathscr{F}_{n, \delta}^M, \mathbb{L}_2 ) < \infty$ as a consequence of Lemma 2.1 in \cite{27}, where it is shown that the function classes under consideration are VC of uniform bounded index.
	\end{proof}

	In the following Lemma, we make use of the concept of weak convergence of a process in $\ell^{\infty}(K)$ for every compact $K \subset \mathbb{R}$ and $\ell^{\infty}$ being the space of bounded functions on $K$ equipped with the supremum norm. For a review of the conditional weak convergence in $\ell^{\infty}$, see \cite{5} (c.f.\ Example 1.5.1 and Theorem 1.5.4).

	\begin{lemma}[Conditional argmax continuous mapping]\label{lemma: cond argmax}
		Let $\mathbb{Z}_n, \mathbb{Z}$ be stochastic processes indexed by $H \subseteq \mathbb{R}$ such that in probability given a sequence of random variables $Z_1,\ldots,Z_n$:
		\begin{align*}
			\mathbb{Z}_n \mid Z_1,\ldots, Z_n \rightsquigarrow \mathbb{Z} \quad \text{in} \quad \ell^{\infty}(K)
		\end{align*}
		for every compact $K \subset H$. Suppose that almost all sample paths $t \mapsto \mathbb{Z}(t)$ are upper semicontinuous and possess a unique maximum at a (random) point $\hat{t}$, which as a random map in $H$ is tight. If $\, \forall \, n \in \mathbb{N}$, $\mathbb{Z}_n (\hat{t}_n) = \sup_t \mathbb{Z}_n (t)$ a.s.\ and there exist a compact set $\tilde{K} \subset H$ such that, under $\pb$:
			\begin{align*}
				\pb \left( \hat{t}_n \notin \tilde{K} \mid Z_1,\ldots, Z_n  \right) \stackrel{\pb}{\rightarrow} 0,
			\end{align*}
		then $\hat{t}_n \mid Z_1,\ldots,Z_n \rightsquigarrow \hat{t}$ in probability in $H$.
	\end{lemma}
	\begin{proof}
		By the continuous mapping theorem, for every closed set \( F \) and \( A, B \) arbitrary sets of \(H\): $\sup_{t \in F \cap A} \mathbb{Z}_n(t) - \sup_{t \in B} \mathbb{Z}_n(t)$ converges conditionally on \(Z_1, \ldots, Z_n\) in distribution to $\sup_{t \in F \cap A} \mathbb{Z}(t) - \sup_{t \in B} \mathbb{Z}(t),$ where \( \mathbb{Z}(t)\) has a.s. continuous sample paths and a.s. unique maxima.
		Thus using the assumption on $\hat{t}_n$:
		\begin{align*}
		\limsup_{n \to \infty} \pb(\hat{t}_n \in F \cap A \mid Z_1, \ldots, Z_n)
		&\leq \limsup_{n \to \infty} \pb\left( \sup_{t \in F \cap A} \mathbb{Z}_n(t) > \sup_{t \in B} \mathbb{Z}_n(t) \mid Z_1, \ldots, Z_n \right) \\
		&= \pb\left( \sup_{t \in F \cap A} \mathbb{Z}(t) > \sup_{t \in B} \mathbb{Z}(t) \right)
		\leq \pb(\hat{t} \in F \cup B^c),
		\end{align*}
		where in the last inequality we use that \( \mathbb{Z}(\hat{t}) > \sup_{t \notin F^c, t \in A} \mathbb{Z}(t)\) and thus the event 
		\[
		\left\{\sup_{t \in F \cap A} \mathbb{Z}(t) > \sup_{t \in B} \mathbb{Z}(t)\right\} \subseteq \{ \hat{t} \in F \cup B^c\}.
		\]
		Because \(Z_n\) converges in distribution conditionally on \(Z_1, \ldots, Z_n\) to \(Z\) in \(\ell^\infty(K)\) for every compact \(K\), take now \(A = B = K\). Then for every open set \(G\) containing \(\hat{t}\), it holds
		\[
		\mathbb{Z}(\hat{t}) > \sup_{t \notin G, t \in K} \mathbb{Z}(t),
		\]
		because of the continuity of \(\mathbb{Z}\) and the uniqueness of \(\hat{t}\). Thus, the above yields:
		\begin{align*}
		&\limsup_{n \to \infty} \pb(\hat{t}_n \in F \mid Z_1, \ldots, Z_n) 
		\leq \limsup_{n \to \infty} \pb(\hat{t}_n \in F \cap K \mid Z_1, \ldots, Z_n) 
		 \\
		&+ \limsup_{n \to \infty} \pb(\hat{t}_n \in K^c \mid Z_1, \ldots, Z_n) = \pb(\hat{t} \in F \cup K^c).
		\end{align*}
		Note: $\pb(\hat{t} \in F \cup K^c) \leq \pb(\hat{t} \in F) + \pb(\hat{t} \in K^c),$
		where by the stochastic boundedness of \(\hat{t}_n \mid Z_1, \ldots, Z_n\) and the choice of \(K^c\), the terms \(\pb(\hat{t} \in K^c)\) can be made arbitrarily small and $\limsup_{n \to \infty} \pb(\hat{t}_n \in K^c \mid Z_1, \ldots, Z_n)$ converges to $0$.		
	\end{proof}

	\begin{lemma}[Consistency]\label{lemma: consistency chapter3}
		Assume \eqref{eq: condition roughness} holds true. The IIP is consistent, i.e. $\forall \, \varepsilon >0$ as $n \rightarrow \infty$:
		\begin{align}\label{eq: consistency}
			\Pi_n (| \hat{V}_{\scriptscriptstyle{G}}(x) - V_0(x)| > \varepsilon \mid Z_1,\ldots, Z_n) \stackrel{G_0}{\rightarrow} 0.
		\end{align}
		Moreover:
		\begin{align}\label{eq: consistency IIP}
			\Pi_n \left( \left| T_{\scriptscriptstyle{G}} \left( V_0(x) + \delta_n  (a + Z^{\scriptscriptstyle{G_0}}_n(1)) \right) - x \right| > \varepsilon \mid Z_1, \ldots, Z_n \right) \stackrel{G_0}{\to} 0.
		\end{align}
	\end{lemma}
	\begin{proof}
		We begin by showing \eqref{eq: consistency}. Let $\psi_x(z) := 2 \sqrt{z} - 2 \sqrt{(z-x)_+}$, using Proposition 4.3 in \cite{4} with $\psi_x$ we obtain:
		\begin{align*}
			&\ex \, (U_{\scriptscriptstyle{G}}(x) \mid Z_1,\ldots,Z_n) = \frac{1}{|\alpha|+n}\int \psi_x  \, d \, \bigg(\alpha + \sum_{i=1}^n \delta_{Z_i} \bigg), \\
			&\mathrm{Var} \, (U_{\scriptscriptstyle{G}}(x) \mid Z_1,\ldots,Z_n) \lesssim \frac{1}{1+|\alpha|+n} \frac{1}{|\alpha|+n} \int \psi_x^2 \,  d \, \bigg(\alpha + \sum_{i=1}^n \delta_{Z_i} \bigg).
		\end{align*}
		Thus, using the conditional Chebyshev inequality:
		\[
		\Pi_n \left( \lvert U_{\scriptscriptstyle{G}} (x) - U_0(x) \rvert > \varepsilon \mid Z_1, \ldots, Z_n \right) 
		\rightarrow 0, \quad \text{a.s.}
		\]
		Note now that the functions \(U_{\scriptscriptstyle{G}}\) are non-decreasing, \(U_0\) is bounded continuous non-decreasing, thus \(\forall \, \varepsilon > 0\) by Exercise 3.2 in \cite{2} we conclude: 
		$$\Pi_n \left( \| U_{\scriptscriptstyle{G}} - U_0 \|_{\infty} > \varepsilon \mid Z_1, \ldots, Z_n \right)  \to 0, \quad \text{a.s.}$$
		Thus, by Marshall's inequality (c.f.\ Exercise 3.1 in \cite{2}), with $ U_{\scriptscriptstyle{G}}^*$ least concave majorant of $ U_{\scriptscriptstyle{G}}$:
		\[
		\Pi_n \left(\lVert U_{\scriptscriptstyle{G}}^* - U_0 \rVert_\infty > \varepsilon \mid Z_1, \ldots, Z_n \right) \stackrel{G_0}{\rightarrow} 0.
		\]
		Now fix $\varepsilon >0$ and $h >0$ such that $| (U_0(x)- U_0(x-h))/h - V_0(x)| < \varepsilon$. By concavity of $U^*_{\scriptscriptstyle{G}}$ we have:
		\begin{align*}
			V_{\scriptscriptstyle{G}}(x) \leq \frac{U^*_{\scriptscriptstyle{G}}(x) - U^*_{\scriptscriptstyle{G}}(x-h)}{h}.
		\end{align*}
		Thus by Lemma 3.1 in \cite{2}:
		\begin{align*}
			&\Pi_n \, (V_{\scriptscriptstyle{G}}(x) \geq V_0(x) + 2 \varepsilon \mid Z_1,\ldots, Z_n) \\
			&\leq \Pi_n \, \left( \frac{U^*_{\scriptscriptstyle{G}} (x) - U^*_{\scriptscriptstyle{G}} (x-h)}{h} \geq \frac{U_{0} (x) - U_{0} (x-h)}{h} +\varepsilon \mid Z_1, \ldots, Z_n \right) \to 0, \quad \text{a.s.} 
		\end{align*}		
		Similarly $V_{\scriptscriptstyle{G}}(x) \geq (U^*_{\scriptscriptstyle{G}}(x+h) - U^*_{\scriptscriptstyle{G}}(x))/h$ gives $\Pi_n \, (V_{\scriptscriptstyle{G}}(x) \leq V_0(x) - 2 \varepsilon \mid Z_1,\ldots, Z_n) \to 0, \quad \text{a.s.}$
		This being true for every $\varepsilon >0$ gives \eqref{eq: consistency}. 
		
		From \eqref{eq: consistency} follows consistency of \(T_{\scriptscriptstyle{G}}\); by the switch relation, \ $\, \forall \, \varepsilon >0$:
		\begin{align*}
		&\Pi_n \left( \left| T_{\scriptscriptstyle{G}} \left( V_0(x) + \delta_n  (a + Z^{\scriptscriptstyle{G_0}}_n(1)) \right) - x \right| > \varepsilon \mid Z_1, \ldots, Z_n \right) \\
		&\leq \Pi_n \left( T_{\scriptscriptstyle{G}} \left( V_0(x) + \delta_n (a + Z^{\scriptscriptstyle{G_0}}_n(1)) \right)   > x+\varepsilon \mid Z_1, \ldots, Z_n \right)  \\
		&\quad + \Pi_n \left(  T_{\scriptscriptstyle{G}} \left( V_0(x) + \delta_n (a + Z^{\scriptscriptstyle{G_0}}_n(1)) \right)  < x - \varepsilon \mid Z_1, \ldots, Z_n \right) \\
		&\leq \Pi_n \left( \hat{V}_{\scriptscriptstyle{G}} (x+\varepsilon) > V_0(x) + \delta_n (a + Z^{\scriptscriptstyle{G_0}}_n(1))  \mid Z_1, \ldots, Z_n \right)  \\
		&\quad + \Pi_n \left( \hat{V}_{\scriptscriptstyle{G}} (x-\varepsilon) < V_0(x) + \delta_n  (a + Z^{\scriptscriptstyle{G_0}}_n(1))   \mid Z_1, \ldots, Z_n \right) \stackrel{G_0}{\to} 0, \numberthis \label{eq: consistency of T_G chapter3}
		\end{align*}
		where by \eqref{eq: consistency} above $\hat{V}_{\scriptscriptstyle{G}} (x \pm \varepsilon) \mid Z_1, \ldots, Z_n \rightarrow V_0(x \pm \varepsilon)$, $Z^{\scriptscriptstyle{G_0}}_n(1) \rightsquigarrow \mathbb{Z}(1)$ and $\delta_n \downarrow 0$, thus $\delta_n  (a + Z^{\scriptscriptstyle{G_0}}_n(1)) \stackrel{G_0}{\to} 0$ and $V_0$ is strictly decreasing at $x$ thus we have the strict inequalities $ V_0(x+ \varepsilon) < V_0(x) <  V_0(x- \varepsilon)$.  
	\end{proof}

	\begin{proof}[\hypertarget{proof of stoch boundedness}{\textbf{Proof of}} \eqref{eq: conditional stoch boundedness chapter3}]
		The claim is implied by the unconditional version of the convergence given in \eqref{eq: conditional stoch boundedness chapter3} and this entails proving that the rate for $\left| T_{\scriptscriptstyle{G}} \left( V_0(x) + \delta_n (a + Z^{\scriptscriptstyle{G_0}}_n(1)) \right) - x \right|$ is $\left( \delta_n^* \right)^{-1}$. Because of Lemma \ref{lemma: consistency chapter3}, we can proceed by verifying the rest of the conditions of Theorem 3.2.5 in \cite{5} (or see Theorem 3.4.1 for a more general version). 

		The argmax process we are interested in is $\tilde{T}_n :=  \argmax_t \big\{ \tilde{Z}^{\scriptscriptstyle{G}}_n(t)-a t \big\}$ where $\tilde{Z}^{\scriptscriptstyle{G}}_n$ is given in \eqref{eq: first components Z tilde}-\eqref{eq: second components Z tilde}. Thus note that we can rewrite it as: $\tilde{T}_n = (\delta^*_n)^{-1} \argmax_t \big\{ \mathbb{M}_n(t)\big\} $, where:
		\begin{align}\label{eq: mathbbM_n}
			\mathbb{M}_n(t) := \, \, &2 \int \bigg( \sqrt{(z-x)_+} - \sqrt{(z-x -t)_+} - \frac{1}{2} V_0(x)t \bigg) \, d (G - \mathbb{G}_n+G_0)(z)\\
			&- 2 \frac{t}{\delta^*_n}  \int \bigg( \sqrt{(z-x)_+} - \sqrt{(z-x -\delta^*_n)_+} \bigg) \, d (\mathbb{G}_n-G_0)(z) \\
			&+ 2  \int \bigg( \sqrt{(z-x)_+} - \sqrt{(z-x -t)_+} \bigg) \, d (\mathbb{G}_n-G_0)(z)  - at \delta_n .
		\end{align}
		Moreover define the deterministic function:
		\begin{align*}
			\mathbb{M}(t):= 2 \int \bigg( \sqrt{(z-x)_+} - \sqrt{(z-x -t)_+} - \frac{1}{2} V_0(x)t \bigg) \, d G_0(z).
		\end{align*}
		Note that $\mathbb{M}_n(0) = \mathbb{M}(0) = 0$. In \cite{27} (c.f.\ proof of (18)) it is shown that $\exists \, \delta_0>0$ such that $\forall \, |t|< \delta_0$:
		\begin{align}\label{eq: bound curvature}
			\mathbb{M}(t) - \mathbb{M}(0) \leq - K \, |t|^{\gamma +1}
		\end{align} 
		Following Theorem 3.2.5 in \cite{5}, define thus the metric $d(s,t) := |s-t|^{\frac{\gamma+1}{2}}.$

		Consider the class of functions $h_t :[0, \infty) \rightarrow \mathbb{R}$:
		\begin{align*}
		\mathscr{H} := \big\{h_t \: : \: h_t (z) =& \: 2 \sqrt{(z-x)_{+}} -2 \sqrt{(z-x-t)_{+}}  - V_0(x)t, \: t \geq 0 \big\}
		\end{align*}
		and, for $\eta>0$, its subclass $ \mathscr{H}_\eta :=\left\{h_t \in \mathscr{H}:|t| \leq \eta\right\}$. In \cite{1} it is shown that the envelope of $\mathscr{H}_\eta$, $H_{\eta} (z) := 2 \sqrt{(z-x)_{+}} -2 \sqrt{(z-x-\eta)_{+}} - V_0(x)\eta$, satisfies for all $\eta$ small enough (c.f.\ (11) in \cite{1}):
		\begin{align}\label{eq: bound envelope 1}
			\int H^2_{\eta}(z) \, g_0(z) \, dz \leq -2g_0(x) \eta^2 \log{(\eta)}.
		\end{align}

		Consider now the class of functions $\widetilde{\mathscr{H}}_n := \{ \tilde{h}^n_t \, : \, t \geq 0 \}$ where, for any fixed $t \geq 0$:
		\begin{align*}
		\tilde{h}^n_t (z) := \:  2 \bigg(\frac{t}{\delta^*_n} -1\bigg) \sqrt{(z-x)_{+}} -2\frac{t}{\delta^*_n} \sqrt{(z-x-\delta^*_n)_{+}}   +2 \sqrt{(z-x-t)_{+}} ,
		\end{align*}
		and, for $\eta>0$, its subclass $ \widetilde{\mathscr{H}}_{n,\eta}:=\big\{\tilde{h}^n_t \in \widetilde{\mathscr{H}}_n:|t| \leq \eta\big\}$. Using the same computations as in \cite{1}, it follows that the envelope of $\widetilde{\mathscr{H}}_{n,\eta}$, $\tilde{H}_{n,\eta} (z) := \sup_{|t|< \eta }|\tilde{h}^n_t(z)|$, satisfies for all $\eta$ small enough:
		\begin{align}\label{eq: bound envelope 2}
			\int \tilde{H}^2_{n,\eta} (z) \, g_0(z) \, dz \leq -2g_0(x) \eta^2 (\log{(\eta)} + \log{(\delta^*_n)}).
		\end{align}

		Using the representation \eqref{eq: representation process G} we have:
		\begin{align*}
			\mathbb{M}_n(t) - \mathbb{M}(t) &=  (1-V_n) \int h_t(z) \, d (\mathbb{B}_n - \mathbb{G}_n)(z)   + V_n \int h_t(z) \, d (Q- \mathbb{G}_n)(z) - at \delta_n \\
			& \quad -  \int \tilde{h}^n_t(z) \, d (\mathbb{G}_n-G_0)(z), 
		\end{align*}
		and, by applying the maximal inequality in Theorem 2.14.1 in \cite{5} for sufficiently small $\delta >0$,  we conclude ($d(t,0)< \delta \, \iff \, |t|< \delta^{\frac{2}{1+\gamma}}$):
		\begin{align*}
			&\ex \bigg[ \sup_{|t|< \delta^{\frac{2}{1+\gamma}}} \sqrt{n} |(\mathbb{M}_n - \mathbb{M})(t)| \bigg] \lesssim J(1, \mathscr{H}_{\delta^{\frac{2}{1+\gamma}}}, \mathbb{L}_2) \sqrt{\int H^2_{\delta^{\frac{2}{1+\gamma}}} \, dG_0} + \sqrt{n} \frac{1}{n} \delta^{\frac{2}{1+\gamma}} \\
			&\quad + |a| \delta^{\frac{2}{1+\gamma}} \delta_n \sqrt{n} + J(1, \widetilde{\mathscr{H}}_{n,\delta^{\frac{2}{1+\gamma}}}, \mathbb{L}_2) \sqrt{\int \tilde{H}^2_{n,\scriptscriptstyle{\delta^{\frac{2}{1+\gamma}}}}\, dG_0 \,}  \lesssim  \delta^{\frac{2}{1+\gamma}} \big(\sqrt{\log{\delta}} + \sqrt{\log{n}} \big).
		\end{align*}
		The first term after applying the maximal inequality comes from the $(1-V_n)(\mathbb{B}_n - \mathbb{G}_n)$-term in the representation above of $\mathbb{M}_n - \mathbb{M}$; the second one comes from the $V_n(Q - \mathbb{G}_n)$-term, the third one from the $at \delta_n $-term and the last one from the $\mathbb{G}_n-G_0$-term (having used also for this one the maximal inequality). The last inequality is just a consequence of \eqref{eq: bound envelope 1}, \eqref{eq: bound envelope 2} and the facts that $J(1, \mathscr{H}_{\delta^{\frac{2}{1+\gamma}}}, \mathbb{L}_2) < \infty$ by the proof of Lemma 2.1 in \cite{27} and by the same reasonings $J(1, \widetilde{\mathscr{H}}_{n,\delta^{\frac{2}{1+\gamma}}}, \mathbb{L}_2) $ is uniformly bounded.

		The above computations and Theorem 3.2.5 in \cite{5} imply that the rate for $d$ is determined as follows: $\delta^{\frac{2}{1+\gamma}} (\sqrt{\log{\delta}} + \sqrt{\log{n}}) \leq \sqrt{n} \delta^2 \ $, hence $ \delta \gtrsim \left( \frac{\log{n}}{n}\right)^{\frac{1+\gamma}{4 \gamma}}$.
		Because $d(s,t) = |s-t|^{\frac{\gamma+1}{2}}$, the rate for the distance $|\cdot|$ is given by: 
		$$\left( \frac{\log{n}}{n}\right)^{\frac{1+\gamma}{4 \gamma} \cdot \frac{2}{1+\gamma}} = \left( \frac{\log{n}}{n}\right)^{\frac{1}{2\gamma}} \equiv \delta^*_n.$$
		Altogether the above proves that the rate for $\argmax_t \big\{ \mathbb{M}_n(t)\big\}$ is $(\delta^*_n)^{-1}$.
	\end{proof}

\subsection{Proofs of complemental statements for Hierarchical Modelling}

\subsubsection{Proofs of the consistency results}\label{III.sec: proof consistency results}

\begin{proof}[\hypertarget{III.proof consistency in weak topology}{\textbf{Proof of Theorem} \ref{III.thm: weak topology consistency}}]
    Fix $\varepsilon > 0$, let $F_\varepsilon$ be $F_0$ restricted and renormalized to $[\varepsilon, L]$. Then $ g_{\scriptscriptstyle{F_\varepsilon}} > 0$ whenever $g_0 > 0$ and hence:
    \begin{equation}
    K(g_0, g_{\scriptscriptstyle{F}}) = K(g_0,  g_{\scriptscriptstyle{F_\varepsilon}}) + \int \left( \log \frac{ g_{\scriptscriptstyle{F_\varepsilon}}}{g_{\scriptscriptstyle{F}}} (z) \right) g_0(z) \, dz.
    \end{equation}
    We shall show: (i) that the first term on the right tends to 0 as $\varepsilon \downarrow 0$ and (ii) that the $\limsup$ of the second term is non-positive if $F \leadsto F_\varepsilon$ weakly, for any fixed $\varepsilon > 0$. Because the DP($\alpha$) prior puts positive mass on every weak neighborhood of $F_\varepsilon$ (see Theorem 4.15 in \cite{4}), together (i) and (ii) show $g_0 \in \operatorname{KL}(\Pi)$. Next the consistency is immediate from Example 6.20 in \cite{4}.
    
    (i) Because $F_0[\varepsilon, L] \rightarrow 1$ as $\varepsilon \downarrow 0$, we have $ g_{\scriptscriptstyle{F_\varepsilon}} \rightarrow g_0$ pointwise on $(0, \infty)$ and hence $\| g_{\scriptscriptstyle{F_\varepsilon}} - g_0\|_1 \rightarrow 0$ by Sheffe's Theorem. % possible typo: Sheffe's, Sheffé's or Scheffé's?
    We strengthen this convergence in KL-convergence by providing bounds on $\log \left( \frac{g_0}{ g_{\scriptscriptstyle{F_\varepsilon}}} \right)$.
    
    First $\log \frac{g_0}{ g_{\scriptscriptstyle{F_\varepsilon}}} \geq 1 - \frac{ g_{\scriptscriptstyle{F_\varepsilon}}}{g_0}$ since $-\log x \geq 1 - x$ for $x > 0$ where $\frac{ g_{\scriptscriptstyle{F_\varepsilon}}}{g_0}$ is integrable relative to $g_0$. Second: note we can bound $ g_{\scriptscriptstyle{F_\varepsilon}}$ from below as follows:
    \begin{equation*}
    g_{\scriptscriptstyle{F_\varepsilon}}(z) = \int_{z \vee \varepsilon}^{L} \frac{1}{2 \sqrt{x(x-z)}} \cdot \frac{dF_0(x)}{F_0[\varepsilon, L]} \geq \frac{F_0(\max(z,1), L)}{\sqrt{L^2 - Lz}}\geq \frac{F_0(\max(z,1), L)}{L}.
    \end{equation*}
    Now we check integrability:
    \begin{equation*}
        \ex_0 \left| \log L \frac{g_0(z)}{F_0(\max (Z,1), L)} \right| \lesssim  \ex_0 \log g_0(z) + \int_0^L g_0(z) \left| \log F_0(\max(z,1), L) \right| dz
    \end{equation*}
    which is integrable. This concludes the proof of (i).

    (ii) For fixed $x,x_1 \in \mathbb{R^+}$, $\psi(z;x) := \frac{\ind_{x > z}}{2\sqrt{x^2 - xz}}$:
    \begin{align}\label{III.eq: L1 distance kernels}
        \| \psi_x(\cdot) - \psi_{x_1} (\cdot) \|_1 = 2 \frac{\sqrt{x_1 -x}}{\sqrt{x_1}} \mathbbm{1}_{\{ x \leq x_1 \}} + 2 \frac{\sqrt{x -x_1}}{\sqrt{x}} \mathbbm{1}_{\{ x_1 \leq x \}}. 
    \end{align} 
    Now for any partition $[0, L] = \bigcup_{i=0}^N I_i$ and points $x_i \in I_i$, using \eqref{III.eq: L1 distance kernels}
    \begin{align*}
    \|g_{\scriptscriptstyle{F}} - g_{\scriptscriptstyle{F_\varepsilon}}\|_1 &= \left\| \sum_{i=1}^N \int_{I_i} (\psi_x - \psi_{x_i}) d(F - F_\varepsilon)(x) + \sum_{i=1}^N \psi_i (F - F_\varepsilon)(I_i) \right\| \\
    &\leq \sum_{i=1}^N \sup_{x \in I_i} \| \psi_x - \psi_{x_i} \|_1 (F + F_\varepsilon)(I_i) + \sum_{i=1}^N |(F - F_\varepsilon)(I_i)| \\
    &\leq 2 (F + F_\varepsilon)(I_0) + \max_{i \geq 1} \sup_{x \in I_i} 2 \sqrt{\frac{|x - x_i|}{x \vee x_i}} + \sum_{i=1}^N |(F - F_\varepsilon)(I_i)|.
    \end{align*}
    We can choose a partition in intervals with boundary points continuity points of $F_\varepsilon$ and $I_0 \subset [0, \varepsilon]$ so that $F_\varepsilon (I_0) = 0$. Then the first and last terms on the right tend to zero as $F \leadsto F_\varepsilon$ weakly. We can further choose the partition fine enough so that the second term is arbitrarily small. Thus $\|g_{\scriptscriptstyle{F}} - g_{\scriptscriptstyle{F_\varepsilon}}\|_1 \rightarrow 0$ as $F \leadsto F_\varepsilon$ weakly. It follows that $\log \frac{g_{\scriptscriptstyle{F_\varepsilon}}}{g_{\scriptscriptstyle{F}}} \rightarrow 0$ in measure.

    If $F \leadsto F_\varepsilon$ weakly, by Portmanteau lemma $\exists \, \delta > 0$ and $0 < c \leq 1$ such that $\forall \, z \in [L - \delta, L)$:
    \begin{equation*}
    \liminf F(z, L) \geq c F_\varepsilon (z, L).
    \end{equation*}
    Therefore we can choose neighborhoods of $F$s with positive prior probability such that $\forall \, z \in [L - \delta, L)$:
    \begin{equation*}
    g_{\scriptscriptstyle{F}}(z) = \int \frac{dF(x)}{2 \sqrt{x^2 - xz}} \geq \frac{F(z, L)}{2L} \geq \frac{c F_\varepsilon (z, L)}{2L}.
    \end{equation*}
    Conclude:
    \begin{align*}
    &\int_0^L g_0(z) \log \left( \frac{g_{\scriptscriptstyle{F_\varepsilon}}}{g_{\scriptscriptstyle{F}}} (z) \right) dz = \int_0^{L-\delta} g_0(z) \log \frac{g_{\scriptscriptstyle{F_\varepsilon}}}{g_{\scriptscriptstyle{F}}} (z) dz + \int_{L - \delta}^L g_0(z) \log \frac{g_{\scriptscriptstyle{F_\varepsilon}}}{g_{\scriptscriptstyle{F}}} (z) dz \\
    &\quad \leq \int_0^{L-\delta} g_0(z) \log \frac{2 g_0(z) L }{F_0(\varepsilon, L)  F(z, L)} dz + \int_{L-\delta}^L g_0(z) \log \frac{g_0(z) 2L}{c F_0(z, L)} dz \\
    &\quad \lesssim \int_0^{L-\delta} g_0(z) \log g_0(z) dz + \int_{L-\delta}^L g_0(z) \log \frac{1}{F_0(z, L)} dz < \infty,
    \end{align*}
    where we used that $F \leadsto F_\varepsilon$ and $F(L - \delta, L) > 0$ a.s.
\end{proof}

\begin{proof}[\hypertarget{III.proof consistency in L1}{\textbf{Proof of Theorem} \ref{III.thm: L1 consistency}}]
We verify the conditions of Theorem 6.23 from \cite{4}. We also use the result derived above for the Kullback-Leibler property. 
Now consider the following partitions of the space $\mathfrak{X}$ where the $X$'s take values: set $\mathfrak{X}_n = [\frac{1}{a_n},L]$, for $a_n = n^{\beta}$. Then $\text{diam}(\mathfrak{X}_n) = L - \frac{1}{a_n}$.
For given $\varepsilon>0$, for $N_n \sim$ $\eta n / \log n$ and small $\eta >0$ to be determined at the end of the proof, define $\mathcal{P}_n$ by:
$$
\mathcal{P}_n=\left\{\sum_{j=1}^{\infty} w_j \psi\left(\cdot ; x_j\right):\left(w_j\right) \in \mathbb{S}_{\infty}, \sum_{j>N_n} w_j<\frac{\varepsilon}{8}, x_1,\ldots,x_{N_n} \in \mathfrak{X}_n \right\} .
$$
By the series representation (Sethuraman) $F =\sum_j W_j \delta_{X_j}$ of the Dirichlet process, the prior density $g_{\scriptscriptstyle{F}}$ is contained in $\mathcal{P}_n$, unless the weights of the representation satisfy $\sum_{j>N_n} W_j \geq \varepsilon / 8$ or one of the points $X_1, \ldots, X_{N_n}$ is not in $\mathfrak{X}_n $. It follows that
$$
\Pi \left(g_{\scriptscriptstyle{F}} \notin \mathcal{P}_n\right) \leq \mathrm{P} \left(\sum_{j>N_n} W_j \geq \frac{\varepsilon}{8}\right) + N_n \overline{\alpha}(\mathfrak{X}_n^{c}).
$$
The last term is exponentially small by assumption. The stick-breaking weights are $W_j=V_j \prod_{l=1}^{j-1}\left(1-V_l\right)$, for $V_l \stackrel{\mathrm{iid}}{\sim} \operatorname{Be}(1,|\alpha|)$ and $\sum_{j>N} W_j=\prod_{j=1}^N\left(1-V_j\right)$. It can be checked that $-\log \left(1-V_j\right)$ is exponentially distributed with parameter $|\alpha|$, so that $R_n:=-\log \sum_{j>N_n} W_j$ possesses a gamma distribution with parameters $N_n$ and $|\alpha|$. Therefore the first term is equal to:
\begin{align*}
\mathrm{P} \left(R_n\leq \log \frac{8}{\varepsilon}\right)&=\int_0^{\log (8 / \varepsilon)} \frac{1}{\left(N_n-1\right) !} r^{N_n-1}|\alpha|^{N_n} e^{-|\alpha| r} d r \leq \\
&\leq \frac{(\log (8 / \varepsilon)|\alpha|)^{N_n}}{N_{n} !} \leq\left(\frac{e \log (8 / \varepsilon)|\alpha|}{N_n}\right)^{N_n}
\end{align*}
The last inequality follows, since $x^N / N ! \leq e^x$, for every $x>0$, including $x=N_n$. For large $n$ and fixed $\varepsilon >0$ the term within brackets is (easily, because $\log n \ll \sqrt{n}$ ) bounded above by $1 / \sqrt{n}$ and hence the right side is bounded by $e^{-\left(N_n / 2\right) \log N_n} \leq e^{n \eta /2}$.

Given any infinite probability vector $w=\left(w_1, w_2, \ldots\right)$, the probability vector $\tilde{w}$ obtained by setting the coordinates with $j>N_n$ to zero and renormalizing the remaining coordinates to sum to 1 (hence $\tilde{w}_j=0$ for $j>N_n$ and $\tilde{w}_j=w_j / \sum_{j \leq N_n} w_j$, otherwise) satisfies $\|w-\tilde{w}\|_1=$ $2 \sum_{j>N_n} w_j$. By the triangle inequality and the fact that $\|\psi(\cdot ; x)\|_1=1$, for $\sum_{j=1}^n w_j\psi\left(\cdot ; x_j\right) \in \mathcal{P}_n$ it follows:
$$
\left\|\sum_{j=1}^{\infty} w_j \psi\left(\cdot ; x_j\right)-\sum_{j \leq N_n} \tilde{w}_j \psi\left(\cdot ; x_j\right)\right\|_1 \leq\|w-\tilde{w}\|_1 = 2 \sum_{j>N_n} w_j \leq \frac{\varepsilon}{4}.
$$
Therefore, the finite mixtures $\sum_{j=1}^{N_n} w_j \psi\left(\cdot ; x_j\right)$, with $\left(w_1, \ldots, w_{N_n}\right)$ an arbitrary probability vector, form an $\varepsilon / 4$-net over $\mathcal{P}_n$ with respect to the $\mathbb{L}_1$-norm. Consequently, by the triangular inequality, a $3 \varepsilon / 4$-net over these finite mixtures is an $\varepsilon$-net over $\mathcal{P}_n$. To construct such a net we restrict $\left(w_1, \ldots, w_{N_n}\right)$ to an $\varepsilon / 4$-net over all probability vectors of length $N_n$, restrict $x_1,\ldots,x_{N_n}$ to an $\frac{\varepsilon^2}{16 a_n}$-net over $\mathfrak{X}_n$. By the triangle inequality:
\begin{align*}
\left\|\sum_{j \leq N_n} w_j \psi\left(\cdot ; x_j\right)- \hspace{-0.1cm} \sum_{j \leq N_n} \tilde{w}_j^{\prime} \psi\left(\cdot ; x_j^{\prime} \right)\right\|_1 &\hspace{-0.2cm} \leq\left\|w-w^{\prime}\right\|_1 \hspace{-0.1cm} +\max _{j \leq N_n}\left\|\psi\left(\cdot ; x_j\right)-\psi\left(\cdot ; x_j^{\prime}\right)\right\|_1 \\
& \hspace{-0.7cm} \leq \frac{\varepsilon}{4} + 2\sqrt{a_n} \max_{j \leq N_n} \sqrt{|x^{\prime}_j - x_j|} \leq \frac{\varepsilon}{4} + \frac{\varepsilon}{2} = \frac{3\varepsilon}{4}. 
\end{align*}
The second upper bound follows from \eqref{III.eq: L1 distance kernels}. It follows that for every $w, x$ there exist $w^{\prime}, x^{\prime}$ in the restricted set, so that the right side is smaller than $3 \varepsilon / 4$.  
By Proposition C.2 in \cite{4} the cardinality of this net is bounded above as follows (denote by $\mathcal{P}_n^{F}$ the set $\mathcal{P}_n$ but composed by finite mixtures, which constitute as shown above a $\frac{3}{4} \varepsilon$-net over $\mathcal{P}_n$):
\begin{align*}
N\left(  \frac{3\varepsilon }{4}, \mathcal{P}_n^{F}, \| \cdot \|_1 \right) &\leq N \left(\frac{\varepsilon}{4} , \{ w = (w_1,\ldots,w_{N_n}) \in \mathbb{R}^{N_n} \: : \: \| w \|_1 = 1 \}, \| \cdot  \|_1 \right) \cdot\\
& \hspace{-2cm} \cdot N \left(\frac{\varepsilon^2}{16a_n^2} , \{ x = (x_1,\ldots,x_{N_n}) \in \mathbb{R}^{N_n} \: : \: \| x \|_\infty \leq \text{diam}(\mathfrak{X}_n) \} , \| \cdot \|_\infty  \right)   \\
& \leq \left( \frac{12}{\varepsilon} \right)^{N_n} \times \left( \frac{3 \cdot 16  a^2_n \cdot \text{diam} (\mathfrak{X}_n)}{\varepsilon^2} \right)^{N_n}
 \end{align*}
Together with the definition $N_n \sim n \eta / \log n$ and the fact that $\log{a_n} \leq \beta \log{n}$ for $\beta >0$, this shows that the log-cardinality of $\mathcal{P}_n$ for sufficiently large $n$:  
$\log N\left(\varepsilon, \mathcal{P}_n,\|\cdot\|_1\right) \leq \frac{n \varepsilon^2}{2}$ , provided $\eta$ is chosen small enough.
\end{proof}

\begin{proof}[\hypertarget{III.proof weak consistency underlying}{\textbf{Proof of Theorem} \ref{III.thm: weak consistency underlying}}]
First note that by substituting $y = z/x$:
$$ \int_{0}^{L} z^k \int_{z}^{L} \frac{1}{\sqrt{x^2 -xz}} dF(x)dz = \int_{0}^{1}y^k \frac{1}{2} \frac{1}{\sqrt{1-y}}dy \int_{0}^{L} x^k dF(x)$$
And using the above derivation we conclude for $c_k = \int_{0}^{1}y^k \frac{1}{2} \frac{1}{\sqrt{1-y}}dy$:
\begin{align*}
    & \lim_{n \rightarrow \infty} \Pi_n \bigg( \bigg|c_k  \int_0^L x^k d(F - F_0) (x) \bigg| > \varepsilon \mid Z_1,\ldots,Z_n \bigg) \\
    &\leq \lim_{n \rightarrow \infty} \Pi_n \left(\bigg| \int_0^{L} z^k g_{\scriptscriptstyle{F}} (z) dz - \int_0^{L} z^k g_0 (z) dz \bigg| > \varepsilon \mid Z_1,\ldots,Z_n \right)  \\ 
    &\leq \lim_{n \rightarrow \infty} \Pi_n \left( \int_0^{L} |z|^k \big|g_{\scriptscriptstyle{F}} (z) - g_0 (z)\big| dz  > \varepsilon \mid Z_1,\ldots,Z_n \right)  \\
    &\leq \lim_{n \rightarrow \infty} \Pi_n (L^k \|g_{\scriptscriptstyle{F}} - g_0\|_1 > \varepsilon \mid Z_1,\ldots,Z_n ) = 0, \: \: \: \: \: \: \: G^n_0 \: \: \text{a.s.}
\end{align*}
By the above we conclude that for any polynomial $p_k$:
\begin{align}\label{III.eq: polynomial convergence}
\lim_{n \rightarrow \infty} \Pi_n \bigg( \bigg| \int_0^L p_k(x) d(F - F_0) (x) \bigg| > \varepsilon \mid Z_1,\ldots,Z_n \bigg) = 0, \: \: \: \: \: \: \:  G^n_0 \: \: \text{a.s.}
\end{align}
In order to prove the claim, by the Portmanteau Lemma, it is enough to show that $\forall \: f \in C_b([0,L])$:
$$ \lim_{n \rightarrow \infty} \Pi_n \bigg( \bigg| \int_0^L f(x) d(F - F_0) (x) \bigg| > \varepsilon \mid Z_1,\ldots,Z_n \bigg) = 0, \: \: \: \: \: \: \:  G^n_0 \: \: \text{a.s.}$$
Recall now that every continuous function on a compact interval can be approximated uniformly by polynomials using the Stone-Weierstrass theorem. This means that $\forall \: \varepsilon > 0$ there exists $k_\varepsilon$ such that $\forall \: x \in [0,L]$ we have $|p_{k_{\varepsilon}} (x) - f(x)| < \varepsilon/4$. Since:
\begin{align*}
    \varepsilon &< \bigg| \int_0^L f(x) d(F - F_0) (x) \bigg| \leq \bigg| \int_0^L f(x) - p_{k_{\varepsilon}}(x) d(F - F_0) (x) \bigg| +  \\ 
    &+ \bigg| \int_0^L p_{k_{\varepsilon}}(x) d(F - F_0) (x) \bigg| \leq \frac{\varepsilon}{2} + \bigg| \int_0^L p_{k_{\varepsilon}}(x) d(F - F_0) (x) \bigg| 
\end{align*}
which means:
\begin{align*}
    \bigg\{ F : \bigg| \int_0^L f(x) d(F - F_0) (x) \bigg| > \varepsilon \bigg\} \subseteq \bigg\{ F : \bigg| \int_0^L p_{k_{\varepsilon}}(x) d(F - F_0) (x) \bigg| > \frac{\varepsilon}{2} \bigg\} 
\end{align*}
Therefore we conclude by \eqref{III.eq: polynomial convergence}:
\begin{align*}
   & \lim_{n \rightarrow \infty} \Pi_n \left( \bigg| \int_0^L f(x)   d(F - F_0) (x) \bigg| > \varepsilon  \mid Z_1,\ldots,Z_n \right) \leq \\ 
    &\leq \lim_{n \rightarrow \infty} \Pi_n \left( \bigg| \int_0^L p_{k_{\varepsilon}}(x) d(F - F_0) (x) \bigg| > \frac{\varepsilon}{2} \mid Z_1,\ldots,Z_n \right)  = 0 \: \: \: \: \: \: \:  G^n_0 \: \: \text{a.s.}
\end{align*}
By Portmanteau Lemma (Theorem A.2 in \cite{4}) and the fact that $\mathbb{R}$ is separable, it follows that the above is equivalent to $\forall \: \varepsilon > 0$:
\begin{align*}
    \lim_{n \rightarrow \infty} \Pi_n \left( d_{\scriptscriptstyle{BL}} \left( F,F_0 \right) > \varepsilon \mid Z_1,\ldots,Z_n \right) = 0 \: \: \: \: \: \: \:  G^n_0 \: \: \text{a.s.}
\end{align*}
where $d_{\scriptscriptstyle{BL}}$ is the bounded Lipschitz metric on the space of probability measures. Again by Portmanteau Lemma (see proof of Proposition 6.2 from \cite{4}) the above is equivalent to: $d_{\scriptscriptstyle{BL}} \left( \Pi_n \left( \cdot \mid Z_1, \ldots,Z_n \right) , \delta_{F_0} (\cdot) \right) \rightarrow 0, \: \:  G^n_0 \: \: \text{a.s.}$.
\end{proof}

\begin{proof}[\hypertarget{III.proof weak consistency underlying}{\textbf{Proof of Theorem} \ref{III.thm: weak consistency underlying F^*}}] Because $g_{\scriptscriptstyle{F}} \stackrel{\mathbb{L}_1}{\to} g_0$, then $g_{\scriptscriptstyle{F}} \rightsquigarrow g_0$, which implies pointwise convergence of the Fourier transforms: $\hat{g}_{\scriptscriptstyle{F}} (\lambda) \to \hat{g}_0(\lambda)$. But $\hat{g}_{\scriptscriptstyle{F}} (\lambda) = \hat{k}(\lambda) \cdot \hat{V}_{\scriptscriptstyle{F}}(\lambda) $, where $\hat{k}$ is the Fourier transform of $k(x) =  -\ind_{\{x < 0 \}}/ \pi \sqrt{-x}$, since $g_{\scriptscriptstyle{F}} = k * V_{\scriptscriptstyle{F}} $. Similarly $\hat{g}_0 = \hat{k} \cdot \hat{V}_0$. Note that $\hat{k}$ is nonzero. Thus also: $\hat{V}_{\scriptscriptstyle{F}}(\lambda) \to \hat{V}_0(\lambda)$, meaning $V_{\scriptscriptstyle{F}} \rightsquigarrow V_0$. 

Because $V_0(0) = \int \delta_0 \, dV_0$ and it is possible to find a sequence of mollifiers $\{f_m\}_m$ (which are bounded continuous functions) such that $\int f_m \, dV_0 \to \int \delta_0 \, dV_0$ as $m \to \infty$, we conclude that also pointwise convergence $V_{\scriptscriptstyle{F}}(0) \to V_0(0)$ holds true. Because $F^*_0(x) = 1 - V_0(x)/V_0(0)$ conclude $\forall \: f \in C_b([0,L])$
$$  \bigg| \int_0^L f(x)   d(F^* - F^*_0) (x) \bigg| \lesssim |V_{\scriptscriptstyle{F}}(0) - V_0(0)| + V^{-1}_0(0) \int_0^L f(x)   d(V_{\scriptscriptstyle{F}} - V_0) (x),$$
which converges to zero as $V_{\scriptscriptstyle{F}} \rightsquigarrow V_0$.
\end{proof}


\subsubsection{Proofs of the contraction rates results}\label{III.sec: proof contraction rates}

\begin{proof}[\hypertarget{III.proof contraction hellinger}{\textbf{Proof of the contraction rate in Hellinger distance, Theorem} \ref{III.thm: contraction rate hellinger}}] We verify the conditions of Theorem 8.11 from \cite{4} with $d$ the Hellinger distance. Let $\bar{\varepsilon}_n = \varepsilon_n = (1 / n)^{1 / 4} \log^{1/4}n$. Correspondingly, we obtain the set $\mathcal{P}_n$ which is $\mathcal{P}_{\varepsilon}$ in \eqref{III.eq: partition sets} with $\varepsilon = \varepsilon_n$. Note that: $(1/ \varepsilon_n)^2 \log (1/ \varepsilon_n) \asymp n \varepsilon_n^2$ and that the conditions required by Proposition \ref{III.prop: bound of kullback neighborhood probability} are the most restrictive and imply the ones required by Lemma \ref{III.lemma: prior mass partition} and Proposition \ref{III.prop: complexity bound model partition}. The chosen $\varepsilon_n$, combined with Proposition \ref{III.prop: bound of kullback neighborhood probability} below implies condition (i) of Theorem 8.11 and Proposition \ref{III.prop: complexity bound model partition} implies condition (ii). Using Lemma \ref{III.lemma: prior mass partition} we obtain condition (iii):
\begin{align*}
    &\frac{\Pi(g_{\scriptscriptstyle{F}} \notin \mathcal{P}_n)}{\Pi \big(F: K\left(g_{\scriptscriptstyle{F}}, g_0\right) \leq \varepsilon^2_n, K_2\left(g_{\scriptscriptstyle{F}}, g_0\right) \leq \varepsilon^2_n \big)}= o \big(e^{-2n \varepsilon_n^2} \big).
\end{align*}
\end{proof}

For $\varepsilon \downarrow 0$ and $\mathfrak{X}_\varepsilon$ to be specified later, define:
\begin{align}\label{III.eq: partition sets}
    \mathcal{P}_{\varepsilon}=\left\{\sum_{j=1}^{\infty} w_j \psi\left(\cdot ; x_j\right):\left(w_j\right) \in \mathbb{S}_{\infty}, \sum_{j>N_{\varepsilon}} w_j<\frac{\varepsilon^2}{32}, x_1,\ldots,x_{N_\varepsilon} \in \mathfrak{X}_\varepsilon \right\}.  
\end{align}
\begin{lemma}[Sieves]\label{III.lemma: prior mass partition}
    Assume $\exists \, \beta >0$ such that $\overline{\alpha} [0,\varepsilon^{\beta}] \leq  e^{-C \frac{1}{\varepsilon^2}\log{\frac{1}{\varepsilon}}}$ for some constant $C>0$. Then for $\mathcal{P}_{\varepsilon}$ as in \eqref{III.eq: partition sets} with $N_\varepsilon \asymp \varepsilon^{-2}$ and all $\varepsilon$ sufficiently small,
    \begin{align*}
        \log{\Pi(g_{\scriptscriptstyle{F}} \notin \mathcal{P}_\varepsilon)} \lesssim -\frac{1}{\varepsilon^2} \log{\left( \frac{1}{\varepsilon} \right)}. 
    \end{align*}
\end{lemma}
\begin{proof}
    Let $\mathfrak{X}_\varepsilon = \operatorname{supp}( \alpha) \setminus [0,\varepsilon^{\beta}]$. By the series representation $F =\sum_j W_j \delta_{X_j}$ of the Dirichlet process, the prior density $g_{\scriptscriptstyle{F}}$ is contained in $\mathcal{P}_\varepsilon$, unless the weights of the representation satisfy $\sum_{j>N_\varepsilon} W_j \geq \varepsilon^2 / 32$ or $X_1,\ldots,X_{N_{\varepsilon}} \notin \mathfrak{X}_\varepsilon$. It follows that
    $$
    \Pi \left(g_{\scriptscriptstyle{F}} \notin \mathcal{P}_\varepsilon \right) \lesssim \mathrm{P} \left(\sum_{j>N_\varepsilon} W_j \geq \frac{\varepsilon^2}{32}\right) + N_\varepsilon \overline{\alpha}(\mathfrak{X}_\varepsilon^{c}).
    $$
    The last term is exponentially small with exponent $-C \frac{1}{\varepsilon^2}$ by assumption and the choice of $N_\varepsilon$.
    The stick-breaking weights in the first term satisfy $W_j=V_j \prod_{l=1}^{j-1}\left(1-V_l\right)$, for $V_l \stackrel{\mathrm{iid}}{\sim} \operatorname{Be}(1,|\alpha|)$ and $\sum_{j>N} W_j=\prod_{j=1}^N\left(1-V_j\right)$. It can be checked that $-\log \left(1-V_j\right)$ is exponentially distributed with parameter $|\alpha|$, so that $R_\varepsilon:=-\log \sum_{j>N_\varepsilon} W_j$ possesses a gamma distribution with parameters $N_\varepsilon$ and $|\alpha|$. Therefore this term is bounded above by:
    \begin{align*}
        &\mathrm{P} \left(R_\varepsilon <\log \frac{32}{\varepsilon^2}\right) =\int_0^{\log\frac{32}{\varepsilon^2}} \frac{1}{\left(N_\varepsilon -1\right) !} r^{N_\varepsilon -1}|\alpha|^{N_\varepsilon} e^{-|\alpha| r} d r  \\
        & \quad \quad \leq \frac{(\log (32 / \varepsilon^2)|\alpha|)^{N_\varepsilon}}{N_{\varepsilon} !}\leq\left(\frac{e \log (32 / \varepsilon^2)|\alpha|} {N_\varepsilon}\right)^{N_\varepsilon} \hspace{-0.2cm} \asymp \exp{\left\{ -\frac{1}{\varepsilon^2  }  \log{\left(\frac{1}{\varepsilon} \right)}   \right\}}.
    \end{align*}
    The second inequality follows, since $x^N / N ! \leq e^x$, for every $x>0$, including $x=N_{\varepsilon}$. 
\end{proof}
\begin{proposition}[Entropy]\label{III.prop: complexity bound model partition}
Let $\mathcal{P}_\varepsilon$ as in \eqref{III.eq: partition sets} with $\mathfrak{X}_\varepsilon = [\varepsilon^{\beta},L]$ for some $\beta >0, L>0$ as in Lemma \ref{III.lemma: prior mass partition} and $N_{\varepsilon} \asymp \varepsilon^{-2}$. Then:
\begin{align*}
    \log{N(\varepsilon,\mathcal{P}_\varepsilon,h)} \lesssim \left(\frac{1}{\varepsilon}\right)^2. 
\end{align*}
\end{proposition}
\begin{proof}
    Given any infinite probability vector $w=\left(w_1, w_2, \ldots\right)$, the probability vector $\tilde{w}$ obtained by setting the coordinates with $j>N_\varepsilon$ to zero and renormalizing the remaining coordinates to sum to 1 (hence $\tilde{w}_j=0$ for $j>N_\varepsilon$ and $\tilde{w}_j=w_j / \sum_{j \leq N_\varepsilon} w_j$, otherwise) satisfies $\|w-\tilde{w}\|_1=$ $2 \sum_{j>N_\varepsilon} w_j$. By the triangle inequality and the fact that $\|\psi_x\|_1=1$, it follows that for $F = \sum_{j=1}^{\infty} w_j \delta_{x_j}$ and $F_1 = \sum_{j\leq N_\varepsilon} \tilde{w}_j \delta_{x_j}$:
    $$
    h^2 (g_{\scriptscriptstyle{F}}, g_{\scriptscriptstyle{F_{\scaleto{1}{3.2pt}}}}) \leq \left\|\sum_{j=1}^{\infty} w_j \psi_{x_j}-\sum_{j \leq N_\varepsilon} \tilde{w}_j \psi_{x_j} \right\|_1 \leq\|w-\tilde{w}\|_1 = 2 \sum_{j>N_\varepsilon} w_j \leq \frac{\varepsilon^2}{16}.
    $$
    Therefore, the finite mixtures $\sum_{j\leq N_{\varepsilon}} w_j \psi_{x_j}$, with $\left(w_1, \ldots, w_{N_\varepsilon}\right)$ an arbitrary probability vector, form an $\varepsilon / 4$-net over $\mathcal{P}_\varepsilon$ with respect to the Hellinger distance. Consequently, by the triangle inequality, a $3 \varepsilon / 4$-net over these finite mixtures is an $\varepsilon$-net over $\mathcal{P}_\varepsilon$. To construct such a net we restrict $\left(w_1, \ldots, w_{N_\varepsilon}\right)$ to an $\varepsilon^2 / 16$-net over all probability vectors of length $N_\varepsilon$, and restrict $x_1,\ldots,x_{N_\varepsilon}$ to an $\frac{ \varepsilon^{4+\beta}}{16}$-net over $\mathfrak{X}_\varepsilon$. If $\bar{F} = \sum_{j \leq N_\varepsilon} w_j \delta_{x_j}$ and $F_2 = \sum_{j \leq N_\varepsilon} w^{\prime}_j \delta_{x^{\prime}_j}$ then by the triangle inequality:
    \begin{align*}
    h^2 (g_{\scriptscriptstyle{\bar{F}}}, g_{\scriptscriptstyle{F_{\scaleto{2}{3.2pt}}}}) & \leq  \left\|\sum_{j \leq N_\varepsilon} w_j \psi_{x_j}-\sum_{j \leq N_\varepsilon} w_j^{\prime} \psi_{x_j^{\prime}} \right\|_1 \leq\left\|w-w^{\prime}\right\|_1+\max _{j \leq N_\varepsilon}\left\|\psi_{x_j}-\psi_{x_j^{\prime}} \right\|_1  \\
    & \leq \frac{\varepsilon^2}{16} + 2 \max_{j \leq N_\varepsilon} \sqrt{\frac{|x^{\prime}_j - x_j|}{x^{\prime}_j \vee x_j}} \leq \frac{\varepsilon^2}{16} + \frac{ \varepsilon^2}{2}.
    \end{align*}
    The second upper bound follows from \eqref{III.eq: L1 distance kernels}. It follows that for every $w, x$ there exist $w^{\prime}, x^{\prime}$ in the restricted set, so that $h (g_{\scriptscriptstyle{\bar{F}}}, g_{\scriptscriptstyle{F_{\scaleto{2}{3.2pt}}}}) \leq 3 \varepsilon / 4$.  
    By Proposition C.2 in \cite{4} the cardinality of this net is bounded above by
    \begin{align*}
     N \left(\frac{\varepsilon^2}{16} , \mathbb{S}_{N_{\varepsilon}}, \| \cdot  \|_1 \right)  \cdot N \left(\frac{\varepsilon^{4+\beta}}{16} , \mathfrak{X}_\varepsilon^{N_{\varepsilon}} , \| \cdot \|_\infty  \right)   \leq \left( \frac{48}{\varepsilon^2} \right)^{N_\varepsilon} \hspace{-0.3cm} \cdot \left( \frac{3 \cdot 16 \cdot \text{diam} (\mathfrak{X}_\varepsilon)}{\varepsilon^{4+\beta}} \right)^{N_\varepsilon}.
     \end{align*}
    Since $N_\varepsilon \asymp \varepsilon^{-2} $, the logarithm of this is bounded by a multiple of
    $\left(\frac{1}{\varepsilon}\right)^2 $.
\end{proof}

\begin{lemma}[Finite approximation]\label{III.lemma: hellinger discrete approximation}
Let the probability measure $F_0$ be supported on $[\delta,L]$ for $0 < \delta < L$ with density $f_0$ bounded away from zero and infinity and $\| 1/ g_0 \|_1 < \infty$. Then there exists a discrete measure $F_1$ with at most $N \asymp \varepsilon^{-2} $ support points such that for $g_{\scriptscriptstyle{F_1}}$ as in \eqref{eq: g_F_0 expression}:
\begin{align*}
    h(g_0,g_{\scriptscriptstyle{F_{\scaleto{1}{3.2pt}}}}) \leq \varepsilon
\end{align*}
\end{lemma}
\begin{proof}
Using Lemma A.1 in \cite{23} and Lemma \ref{III.lemma: discrete supported measure}, let $F_1$ be a discrete measure on $[\delta,L]$ such that for $m \geq 1$:
\begin{align*}
&\int x^{-j} d\left(F_1-F_0\right)(x)=0, \quad \forall \: \: 0 \leq j \leq m, \\
&\int x^{j+1/4} d\left(F_1-F_0\right)(x)=0, \quad \forall \: \: 0 \leq j \leq 2m-1.
\end{align*}
then $F_1$ can be chosen to have at most $3m$ support points. Using the power series:
\begin{align*}
    &\frac{\ind_{\left\{x > z\right\}}}{2 \sqrt{x^2 - xz}} = \frac{1}{2x} \frac{\ind_{\left\{x>z\right\}}}{\sqrt{1 - \frac{z}{x}}} = \frac{1}{2x} \sum_{k=0}^{\infty} \binom{-1/2}{k} \left( - \frac{z}{x} \right)^k \ind_{x>z}, 
\end{align*}
and using the way $F_1$ is constructed, we obtain:
\begin{align*}
    &  h^2(g_0,g_{\scriptscriptstyle{F_{\scaleto{1}{3.2pt}}}}) = \frac{1}{2} \int_0^L  (g^{1/2}_0(z) - g^{1/2}_{\scriptscriptstyle{F_{\scaleto{1}{3.2pt}}}}(z))^2 \, dz \leq \frac{1}{4} \int_{[0, \infty)} \frac{\left(g_{\scriptscriptstyle{F_{\scaleto{1}{3.2pt}}}}(z)-g_0(z)\right)^2}{g_0(z)} \,  d z \\
     & = \int_{0}^L \frac{1}{g_0(z)} \bigg( \int_{z}^\infty \frac{1}{2x} \sum_{k=m+1}^{\infty} \binom{-1/2}{k} \left(- \frac{z}{x} \right)^k \, d\left(F_0 - F_1 \right)(x) \bigg)^2 \, dz \\
     &\lesssim \int_{0}^L \frac{1}{g_0(z)} \bigg( \int_{z}^\infty \frac{1}{2x} \sum_{k=m+1}^{\infty} \binom{-1/2}{k} \left(- \frac{z}{x} \right)^k \, dF_0 (x) \bigg)^2 \, dz \\
     &\quad \quad + \int_{0}^L \frac{1}{g_0(z)} \bigg( \int_{z}^\infty \frac{1}{2x} \sum_{k=m+1}^{\infty} \binom{-1/2}{k} \left(- \frac{z}{x} \right)^k \, dF_1 (x) \bigg)^2 \, dz \\
     &\lesssim \max \{1, \|f_0 \|_{\infty} \} \int_{0}^L \frac{1}{g_0(z)} \bigg( \sum_{k=m+1}^{\infty } \binom{-1/2}{k} (-z)^k \int_{z}^\infty \frac{1}{x^{k+1}} \, dx \bigg)^2 \, dz \numberthis \label{III.eq: to quantify asymp} \\
     & \lesssim \bigg( \sum_{k=m+1}^{\infty}  \frac{1}{\sqrt{k}} \frac{L }{k+1} \bigg)^2 \lesssim \bigg( \sum_{k=m+1}^{\infty}  \frac{1}{k^{3/2}} \bigg)^2 \leq \bigg( \int_{m}^{\infty} \frac{1}{x^{3/2}} \, dx \bigg)^2 \asymp \frac{1}{m}.
\end{align*}
In \eqref{III.eq: to quantify asymp} we used the fact that for all $k$ big enough:
\begin{align*}
    \bigg| \binom{-1/2}{k} \bigg| = \frac{\sqrt{\pi}}{ | \Gamma(\frac{1}{2}-k) \Gamma (k+1)| } \asymp \frac{1}{\sqrt{k}}, 
\end{align*}
 Moreover, it is possible to see that for $k>0$, $\binom{-1/2}{k} (-1)^k >0$, so for fixed $z$ the functions $\binom{-1/2}{k} (-1)^k \frac{z^k}{x^{k+1}}$ are decreasing in $x$ and nonnegative. Therefore using the integral test for series and the fact that $F_1 \{ x_i \}  \in [0,1]$:
\begin{align*}
    \sum_{x_i>z} \frac{1}{2x_i} \sum_{k=m+1}^{\infty} \binom{-1/2}{k} \left(- \frac{z}{x_i} \right)^k \, F_1 \{ x_i \} \leq \int_{z}^\infty \frac{1}{2x} \sum_{k=m+1}^{\infty} \binom{-1/2}{k} \left(- \frac{z}{x} \right)^k \, dx.
\end{align*}
Similarly, we can upper bound the other integral by using $\| f_0 \|_{\infty} < \infty$. 

By choosing $m \asymp \varepsilon^{-2}$ we obtain the claim. 
\end{proof}

\begin{lemma}\label{III.lemma: discrete supported measure}
    Let $F_0$ as in Lemma \ref{III.lemma: hellinger discrete approximation}. Then there exists a discrete probability measure $F$ on $[0,L]$ with at most $N+1$ support points such that
    $$
    \int \frac{1}{x^j} \, d F(x)=\int \frac{1}{x^j} \, d F_0(x) \quad \text { for } j=1, \ldots, N .
    $$
    \end{lemma}
    \begin{proof}
    This lemma is an adaptation of Lemma A.1 in \cite{23}. 
    
    Consider the set
    $
    C=\left\{\left(1/x, \ldots, 1/x^N \right): x \in (0,L]\right\}
    $
    and its convex hull $\operatorname{conv}(C)$. Then for every probability measure $F_0$ on $[0,L]$ with finite inverse moments $\int_{0}^L \frac{1}{x^k} \, dF_0(x) < \infty$ for $k=0,1,2,...$ as in Lemma \ref{III.lemma: hellinger discrete approximation}, 
    $$
    v_0=\left(\int \frac{1}{x} \, d F_0(x), \ldots, \int \frac{1}{x^N} \, d F_0(x)\right)^T \in \operatorname{conv}(C)
    $$
    (see, e.g., Lemma 3 p.\ 74 of \cite{28}). Since $\operatorname{conv}(C) \subset \mathbb{R}^N$, any element of $\operatorname{conv}(C)$ may be written as a convex combination of at most $N+1$ elements of $C$ (see, e.g., \cite{29}). Thus there exist $\lambda_1, \ldots, \lambda_{N+1} \geq$ $0, \sum_{j=1}^{N+1} \lambda_j=1, x_1 \ldots, x_{N+1} \in K$ such that
    $$
    \begin{aligned}
    v_0 & =\sum_{j=1}^{N+1} \lambda_j\left(\frac{1}{x_j}, \ldots, \frac{1}{(x_j)^N}\right)^T =\left(\int \frac{1}{x} \, d F(x), \ldots, \int \frac{1}{x^N} \, d F(x)\right)^T,
    \end{aligned}
    $$
    where $F=\sum_{j=1}^{N+1} \lambda_j \delta_{x_j}$.
\end{proof}


\begin{proposition}[Prior mass]\label{III.prop: bound of kullback neighborhood probability}
Let the density $g_0$ be bounded and let the conditions of Lemma \ref{III.lemma: hellinger discrete approximation} hold true. Let $\Pi$ be the Dirichlet process $\operatorname{DP}(\alpha)$ with precision $|\alpha|$ and center measure $\alpha$ that has a Lebesgue density bounded away from 0 and $\infty$ and bounded support, that satisfies $\mathrm{supp}(F_0) \subseteq [\delta,L] \subset \mathrm{supp}(\alpha)$. For all sufficiently small $\varepsilon>0$,
\begin{align}\label{III.eq: lower bound kullback balls}    
\log \Pi \bigg(F: K\left(g_{\scriptscriptstyle{F}}, g_0\right) \leq \varepsilon^2, K_2\left(g_{\scriptscriptstyle{F}}, g_0\right) \leq \varepsilon^2\bigg) \gtrsim - \left(\frac{1}{\varepsilon}\right)^{2}  \log \left(\frac{ 1}{\varepsilon}\right).
\end{align}
\end{proposition}
\begin{proof}
By Lemma \ref{III.lemma: hellinger discrete approximation} there exists a discrete measure $F_1$ with $N \lesssim \varepsilon^{-2} $ support points such that $h(g_{\scriptscriptstyle{F_{\scaleto{1}{3.2pt}}}}, g_0) \leq \varepsilon$. We can take the support points $x_j$ of $F_1$ such that $\forall \, j$ there exists an interval $I_j$ with $\sup_{x\in I_j} \sqrt{|x_j-x|/ (x_j \vee x)} \lesssim \varepsilon^2$ and $\bigcup_{j=1}^N I_j \subset [0,L]$, $F_1 (x \, : \, x \notin I_j, \, \forall \, j) \leq \varepsilon^2$. If not, we take a maximal set of points in the support points of $F_1$ which respect this property, and replace $F_1$ by the discrete measure obtained by relocating the masses of $F_1$ to the nearest points in the net which fulfill the above requirement. Call the obtained measure $F_2$. Then $h(g_{\scriptscriptstyle{F_{\scaleto{2}{3.2pt}}}}, g_{\scriptscriptstyle{F_{\scaleto{1}{3.2pt}}}}) \lesssim \varepsilon$, as seen in the proof of Proposition \ref{III.prop: complexity bound model partition}. 

By Lemma B.2 in \cite{4} we have: 
\begin{align}\label{III.eq: K_2 neighborhoods subset hellinger ones}
    \max \left(K\left(g_0,g_{\scriptscriptstyle{F}} \right), K_2\left(g_0,g_{\scriptscriptstyle{F}}\right)\right) \lesssim h^2\left(g_0,g_{\scriptscriptstyle{F}}\right) \log^2_{-}{ h\left(g_0,g_{\scriptscriptstyle{F}}\right) }
\end{align}
where $\log_{-}x = \max(-\log x,0)$. In order to conclude we need a further step. We need to upper bound the $\mathbb{L}_1$ distance between $g_0$ and $g_{\scriptscriptstyle{F}}$. Let $F_1=\sum_{i=1}^N w_j \delta_{x_j}$, with the support points $x_j \in I_j$ for every $j$ as set in the beginning of the proof. Next note:
\begin{align*}
     &|g_{\scriptscriptstyle{F}}(z) -g_{\scriptscriptstyle{F_{\scaleto{1}{3.2pt}}}} (z)| = \bigg|\int \psi_x^z dF(x) - \int \psi_x^z dF_1(x) \bigg| \\
    &\leq \int_{x\notin I_j, \: \forall j} \psi_x^z d(F+ F_1)(x) + \sum_{j=1}^N \int_{I_j} |\psi_x^z -\psi_{x_j}^z| dF(x) + \sum_{j=1}^N \psi_{x_j}^z |F(I_j) - w_j|
\end{align*}
Now note that: $\norm{\psi_x - \psi_{x_j}}_1 = 2 \sqrt{\frac{|x-x_j|}{x_j \vee x}}$, that $\norm{\psi_x^z}_1 = 1$, that by construction:
$$\sup_{x\in I_j} \sqrt{|x_j-x|/ (x_j \vee x)} \lesssim \varepsilon^2$$ 
and that:
\begin{align*}
    F(x \: : \: x \notin I_j, \: \forall j) = 1- \sum_{j=1}^{N} F(I_j) \leq \sum_{j=1}^{N} |F (I_j) - w_j|
\end{align*}
Altogether, this gives:
\begin{align*}
    &\norm{g_{\scriptscriptstyle{F}}- g_{\scriptscriptstyle{F_{\scaleto{1}{3.2pt}}}}}_1 \\
    &\leq F (x \: : \: x \notin I_j, \: \forall \: j) + 2 \sum_{j=1}^{N} \int_{I_j} \sqrt{\frac{|x_j-x|}{ x_j \vee x}} dF(x) + \sum_{j=1}^{N} |F(I_j)-w_j| + \varepsilon^2 \\ 
    &\leq 2\sum_{j=1}^{N} |F(I_j) - w_j| + 2 \sum_{j=1}^{N} \sup_{x \in I_j} \sqrt{\frac{|x_j-x|}{ x_j \vee x}} F(I_j) \lesssim 2\sum_{j=1}^{N} |F(I_j) - w_j| + 3 \varepsilon^2. \numberthis \label{III.eq: bound L1 distance with discrete F}
\end{align*}
Thus there exist constants $d_1,d_2$ such that for $F_1=\sum_{i=1}^N w_j \delta_{x_j}$, with the support points $x_j \in I_j, \: \forall j$, 
\begin{align*}
\left\{F: \max \left(K, K_2\right)\left(g_0, g_{\scriptscriptstyle{F}}\right)<d_1 \varepsilon^2 \right\} & \stackrel{\eqref{III.eq: K_2 neighborhoods subset hellinger ones}}{\supset} \left\{F: h\left(g_0, g_{\scriptscriptstyle{F}}\right) \leq 2 \varepsilon / \log_{-}\varepsilon \right\}  \\
& \: \stackrel{\Delta \: \text{in.}}{\supset}\left\{F: h\left(g_{\scriptscriptstyle{F_{\scaleto{1}{3.2pt}}}}, g_{\scriptscriptstyle{F}}\right) \leq \varepsilon / \log_{-}\varepsilon \right\} \\
& \: \stackrel{\eqref{III.eq: bound hellinger with L1}}{\supset}\left\{F:\left\|g_{\scriptscriptstyle{F}}-g_{\scriptscriptstyle{F_{\scaleto{1}{3.2pt}}}}\right\|_1 \leq d_2 \varepsilon^2 / \log^2_{-}\varepsilon \right\} \\
& \stackrel{\eqref{III.eq: bound L1 distance with discrete F}}{\supset}\left\{F: \sum_{j=1}^N\left|F(I_j)-w_j\right| \leq \varepsilon^2/ \log^2_{-}\varepsilon\right\}.
\end{align*}
By Lemma A.2 in \cite{23}, since the base measure $\alpha$ has density bounded away from zero and infinity on its support by assumption, we have
$$
\log \Pi\bigg(F: \sum_{j=1}^N\left|F(I_j)-w_j\right| \leq \varepsilon^2 / \log^2_{-}\varepsilon\bigg) \gtrsim-N \log \left(\frac{ \log_{-}\varepsilon}{\varepsilon}\right).
$$
The proposition follows upon combining the preceding.
\end{proof}

\begin{proof}[\hypertarget{III.proof contraction pointwise V}{\textbf{Proof of pointwise contraction rate for V, Theorem} \ref{III.thm: contraction rate V}}] 
    Let $\mathcal{F}_n$ be such that $\Pi_n (F \in \mathcal{F}_n \mid Z_1, \dots, Z_n) \xrightarrow{G_0} 1$. Without loss of generality, let the $F \in \mathcal{F}_n$ satisfy Theorem \ref{III.thm: contraction rate hellinger}. Denote by $\| \cdot \|_{\mathcal{F}_n} = \sup_{F \in \mathcal{F}_n} | \cdot |$. Now note for $\eta_n = \frac{1}{\sqrt{\log n}}$:
    \begin{align*}
    |V_{\scriptscriptstyle{F}}(x) - V_{0}(x)| &= \left| \int_x^\infty \frac{1}{\sqrt{z-x}} (g_{\scriptscriptstyle{F}} - g_{0})(z) dz \right| \\
    &= \left| \left( \int_x^{x+\eta_n} + \int_{x+\eta_n}^\infty \right) \frac{1}{\sqrt{z-x}} (g_{\scriptscriptstyle{F}} - g_{0})(z) dz \right| \\
    &\leq \left| \int_x^{x+\eta_n} \frac{1}{\sqrt{z-x}} (g_{\scriptscriptstyle{F}} - g_{0})(z) dz \right| + \frac{2}{\sqrt{\eta_n}} h(g_{\scriptscriptstyle{F}}, g_{0})
    \end{align*}
    Then define $V_{\scriptscriptstyle{F}}^{\eta_n}(x) = \int_x^{x+\eta_n} \frac{1}{\sqrt{z-x}} g_F(z) dz$ and note:
    \begin{align*}
        &\Pi_n \left( \frac{\sqrt[4]{n}}{(\log n)^{3/4}} |V_{\scriptscriptstyle{F}}(x) - V_{0}(x)| > M \mid Z_1, \dots, Z_n \right) \\
        &\quad \leq \Pi_n \left( \frac{\sqrt[4]{n}}{(\log n)^{3/4}} |V_{\scriptscriptstyle{F}}^{\eta_n}(x) - V_{0}^{\eta_n}(x)| > M \mid Z_1, \dots, Z_n \right) \\
        &\quad \quad + \Pi_n \left( \frac{\sqrt[4]{n}}{(\log n)^{3/4}} \cdot \sqrt{\log n} \cdot h(g_{\scriptscriptstyle{F}}, g_{0}) > M \mid Z_1, \dots, Z_n \right)
    \end{align*}
    where the second term goes to 0 under $G_0$. By using Lemma 14 in \cite{26} the claim is proven if we show that $\forall \, t$ in a neighborhood of 0, $r_n = \frac{\sqrt[4]{n}}{(\log n)^{3/4}}$:
    \begin{align}\label{III.eq: contraction statement 1}
        \ex \left[e^{t r_n (V_{\scriptscriptstyle{F}}^{\eta_n}(x) - V_{0}^{\eta_n}(x))} \mid Z_1, \dots, Z_n, F \in \mathcal{F}_n \right] \xrightarrow{G_0} 1 = \int e^{ts} d \delta_0(s).
    \end{align}
    The first step to prove the claim consists of the definition of the appropriate perturbations of $F$. Let $\eta_n^* = \frac{1}{\sqrt{n}}$ and define the function 
    $$h_n^x(u) := \frac{2 \sqrt{u}}{\pi f_0(u)}  \frac{1}{u-x} \ind_{ \{\eta_n^* \leq u - x \leq \eta_n \} }$$
    then $\| h_n^x \|_\infty = O(\frac{1}{\eta_n^*})$. Define for $t_n = t r_n/n$.
    \begin{align}\label{III.eq: biased perturbation 1}
        dF_{t_n}(u) := D^{-1}_{t_n} \left( 1 + t_n h^x_n (u) \right) dF(u), \quad D_{t_n} := \int \left( 1 + t_n h^x_n (u) \right) dF(u)
    \end{align}
    It follows that: 
    \begin{align*}
        D_{t_n} g_{\scriptscriptstyle{F_{\scaleto{t_n}{3.5pt}}}}(z) = \int_{z}^{\infty} \frac{D_{t_n} dF_{t_n}(v)}{2\sqrt{v^2 - vz}} =  g_{\scriptscriptstyle{F}}(z) + t_n  \int_{z}^{L} \frac{1}{\pi \sqrt{u-z}} \frac{\ind_{ \{\eta_n^* \leq u - x \leq \eta_n \} }}{u-x}  \, \frac{dF(u)}{f_0(u)}
    \end{align*}
    The proof's second step involves rewriting the left-hand side of \eqref{III.eq: contraction statement 1} as:
    \begin{align}\label{III.eq: rewrite ex bayes}
        \ddfrac{\int_{\mathcal{F}_n} e^{t r_n (V_{\scaleto{F}{3pt}} (x) - V_0(x)) + \sum_{i=1}^n \left(\log{g_{\scaleto{F}{3pt}}(Z_i)} - \log{g_{\scaleto{F_{-t_n}}{5.2pt}}(Z_i)} \right)} \prod_{i=1}^n g_{\scriptscriptstyle{F_{\scaleto{-t_n}{3.5pt}}}}(Z_i) \: d \Pi(F)}{\int_{\mathcal{F}_n} \prod_{i=1}^n g_{\scriptscriptstyle{F}}(Z_i) \: d \Pi(F)}
    \end{align}
    where $F_{-t_n}$ is as in \eqref{III.eq: biased perturbation 1} (note that the inserted perturbation has negative sign).  The proof's third step consists of studying the sum of the log-likelihood ratios:
    \begin{align}\label{III.eq: ll ratio}
        &\sum_{i=1}^n \log{\frac{g_{\scriptscriptstyle{F_{\scaleto{-t_n}{3.5pt}}}}}{g_{\scriptscriptstyle{F}}} (Z_i)} = \sum_{i=1}^n \log \left\{ 1 - \frac{t_n }{D_{t_n}} Y_{ni}^{\scriptscriptstyle{F}} \right\}
    \end{align}
    where $Y_{ni}^{\scriptscriptstyle{F}} = V_{ni}^{\scriptscriptstyle{F}} - \ex V_{ni}^{\scriptscriptstyle{F}}$, where $\ex$ denotes the expectation under $g_{\scriptscriptstyle{F}}$ and:
    \begin{align*}
        V_{ni}^{\scriptscriptstyle{F}} = \frac{1}{g_{\scriptscriptstyle{F}}(Z_i)} \int_{Z_i}^L \frac{1}{\sqrt{u-Z_i}} \cdot \frac{\ind_{\eta_n^* \leq u - x \leq \eta_n}}{\pi (u-x)}  \frac{dF(u)}{f_0(u)}.
    \end{align*}
    Now, using the Taylor expansion of the logarithm $\log \{ 1 - \frac{tY}{1-tK} \}$ at $t=0$ in \eqref{III.eq: ll ratio} at $t=0$, we obtain for some $\tilde{t}$ between $0$ and $t$:
    \begin{align*}
        &\eqref{III.eq: ll ratio} = \sum_{i=1}^n \left\{ \frac{-t_n Y_{ni}^{\scriptscriptstyle{F}}}{D_{\tilde{t}_n} (D_{\tilde{t}_n} - \tilde{t}_n Y_{ni}^{\scriptscriptstyle{F}})} + t_n Y_{ni}^{\scriptscriptstyle{F}} - t_n Y_{ni}^{\scriptscriptstyle{F}} \right\}. 
    \end{align*}
    The supremum over $F \in \mathcal{F}_1$ of the above expression can be upper bounded as follows:
    \begin{align*}
        &n \max_{i \leq i \leq n} \left\| \frac{t_n Y_{ni}^{\scriptscriptstyle{F}}}{D_{\tilde{t}_n} (D_{\tilde{t}_n} - \tilde{t}_n Y_{ni}^{\scriptscriptstyle{F}})} \right\|_{\mathcal{F}_1} \left\| D_{\tilde{t}_n}^2 - \tilde{D}_{\tilde{t}_n} \tilde{t}_n Y_{ni}^{\scriptscriptstyle{F}} - 1 \right\|_{\mathcal{F}_1} \leq \\
        &  n \max_{i \leq i \leq n} \underbrace{\| t_n Y_{ni}^{\scriptscriptstyle{F}} \|_{\mathcal{F}_1}}_{(I)} \underbrace{\max_{i \leq i \leq n} \left\| \frac{1}{D_{\tilde{t}_n} (D_{\tilde{t}_n} - \tilde{t}_n Y_{ni}^{\scriptscriptstyle{F}})} \right\|_{\mathcal{F}_1}}_{(II)} \underbrace{\max_{i \leq i \leq n} \| D_{\tilde{t}_n}^2 - D_{\tilde{t}_n} \tilde{t}_n Y_{ni}^{\scriptscriptstyle{F}} - 1 \|_{\mathcal{F}_1}}_{(III)}
    \end{align*}
    which we prove goes to zero under $G_0$ in probability.

    $(I)$ Since $t_n = t r_n / n$ we show: $\max_{i \leq i \leq n} \| t_n Y_{ni}^{\scriptscriptstyle{F}} \|_{\mathcal{F}_1} = O_{\scriptscriptstyle{G_0}} \left( \frac{r_n}{n} \frac{1}{\sqrt{\eta^*_n}} \right)$. First note:
    \begin{align*}
    &|\ex V_{ni}^{\scriptscriptstyle{F}}| = \left| \int_0^L \int_0^\infty \frac{1}{\sqrt{u-z}} \cdot \frac{\ind_{\eta_n^* \leq u - x \leq \eta_n}}{\pi(u-x)}  \frac{dF(u)}{f_0(u)} dz \right| \\
    &\leq \left\| \frac{1}{f_0} \right\|_\infty \left| \int_0^L \int_0^\infty \frac{1}{\sqrt{u-z}} \cdot \frac{\ind_{\eta_n^* \leq u - x \leq \eta_n}}{\pi(u-x)}  \, dF(u) \, dz \right|  \\
    &\leq \left\| \frac{1}{f_0} \right\|_\infty \int_0^L \int_0^\infty \frac{1}{\sqrt{u-z}} \cdot \frac{1}{\pi(u-x)} \, du \, dz \lesssim \left\| \frac{1}{f_0} \right\|_\infty \int_x^L \frac{1}{\sqrt{z-x}} dz = O(1).
    \end{align*}
    Now, using Lemma 8 from \cite{27} and $ \| \sup_{z \in [x, x + \eta_n]} \frac{1}{g_{\scriptscriptstyle{F}}(z)} \|_{\mathcal{F}_1}< \infty$, for all $1 \leq i \leq n$ uniformly over $\mathcal{F}_1$:
    \begin{align*}
        |V_{ni}^{\scriptscriptstyle{F}}| &\leq \left\| \frac{1}{f_0} \right\|_\infty \frac{1}{g_{\scriptscriptstyle{F}}(Z_i)} \int_{Z_i}^L \frac{1}{\sqrt{u-Z_i}} \cdot \frac{\ind_{\eta_n^* \leq u - x \leq \eta_n}}{\pi(u-x)}  \, dF(u) \\
        &\lesssim \frac{1}{g_{\scriptscriptstyle{F}}(Z_i)} \int_{Z_i}^L \frac{1}{\sqrt{u-Z_i}} \cdot \frac{\ind_{\eta_n^* \leq u - x \leq \eta_n}}{\pi(u-x)} \, du = O_{\scriptscriptstyle{G_0}} \left( \frac{1}{\sqrt{\eta^*_n}} \right).
    \end{align*}
    

    $(II)$ Uniformly for $F \in \mathcal{F}_1$, we have:
    \begin{align*}
    |D_{\tilde{t}_n}| \leq 1 + |\tilde{t}| \frac{r_n}{n} \sup_{F \in \mathcal{F}_1} \int h_n^x(u) dF(u) \leq 1 + O \left( \frac{r_n}{n} \log \left( \frac{1}{\eta_n^*} \right) \right),
    \end{align*}
    because we can use the integral test for series as $(u-x)^{-1}$ is decreasing and nonnegative on $[x+\eta_n^*, x+\eta_n]$. We have:
    \begin{align*}
    \int h_n^x(u) dF(u) &= \int \frac{2 \sqrt{u}}{\pi f_0(u)} \cdot \frac{\ind_{\eta_n^* \leq u - x \leq \eta_n}}{u-x} \, dF(u) \\
    &\leq \left\| \frac{1}{f_0} \right\|_\infty \frac{2 \sqrt{x + \eta_n}}{\pi} \int_{x + \eta_n^*}^{x + \eta_n} \frac{du}{u-x} \sim \log \left(\frac{1}{\eta_n^*}\right).
    \end{align*}
    Using also the results derived for $(I)$ we conclude $(II) = 1 + O_{G_0}(1)$ because:
    \begin{align*}
    &\| D_{\tilde{t}_n} (D_{\tilde{t}_n} - \tilde{t}_n Y_{ni}^{\scriptscriptstyle{F}}) \|_{\mathcal{F}_1} \leq \| D_{\tilde{t}_n} \|_{\mathcal{F}_1}^2 + \| D_{\tilde{t}_n} \|_{\mathcal{F}_1} \| \tilde{t}_n Y_{ni}^{\scriptscriptstyle{F}} \|_{\mathcal{F}_1} \\
    &\quad \quad = 1 + O \left( \frac{r_n}{n} \log \frac{1}{\eta_n^*} \right) + O \left( \frac{r_n^2}{n^2} \log^2 \frac{1}{\eta_n^*} \right) + O_{\scriptscriptstyle{G_0}} \left( \frac{r_n}{n} \frac{1}{\sqrt{\eta_n^*}} \right).
    \end{align*}

    $(III)$ First note that $D_{\tilde{t}_n}^2 - 1 = 2 \int \tilde{t}_n h_n^x(u) dF(u) + \left( \int \tilde{t}_n h_n^x(u) dF(u) \right)^2$. Thus:
    \begin{align*}
    \max_{i \leq i \leq n} \| D_{\tilde{t}_n}^2 - 1 - D_{\tilde{t}_n} \tilde{t}_n Y_{ni}^{\scriptscriptstyle{F}} \|_{\mathcal{F}_1} &\leq \| D_{\tilde{t}_n}^2 - 1 \|_{\mathcal{F}_1} + \| D_{\tilde{t}_n} \|_{\mathcal{F}_1} \| \tilde{t}_n Y_{ni}^{\scriptscriptstyle{F}} \|_{\mathcal{F}_1} \\
    &= O_{\scriptscriptstyle{G_0}} \left( \frac{r_n}{n} \left( \frac{1}{\sqrt{\eta_n^*}} + \log \left( \frac{1}{\eta_n^*} \right) \right) \right).
    \end{align*}

    We conclude that we can reduce the study of the numerator of \eqref{III.eq: rewrite ex bayes} to:
    \begin{align*}
        \int_{\mathcal{F}_n} e^{  t r_n (V_{\scriptscriptstyle{F}}^{\eta_n}(x) - V_0^{\eta_n}(x)) + \sum_{i=1}^n \left\{ \frac{t r_n}{n} Y_{ni}^{\scriptscriptstyle{F}} \right\} \pm t r_n \ex_0 Y_{ni}^{\scriptscriptstyle{F}} } \, \prod_{i=1}^n g_{\scriptscriptstyle{F_{\scaleto{-t_n}{3.5pt}}}}(Z_i) \: d \Pi(F)
    \end{align*}
    where $\ex_0$ denotes the expectation under $g_0$. The proof's fourth step consists in studying the exponent given in the integral in the above display. We now prove:
    \begin{align*}
        &r_n \left\| V_{\scriptscriptstyle{F}}^{\eta_n}(x) - V_0^{\eta_n}(x) + \ex_0 Y_{ni}^{\scriptscriptstyle{F}} \right\|_{\mathcal{F}_n} = o(1), \numberthis \label{III.eq: first to show 0} \\
        &r_n \left\| \frac{1}{n} \sum_{i=1}^n Y_{ni}^{\scriptscriptstyle{F}} - \ex_0 Y_{ni}^{\scriptscriptstyle{F}} \right\|_{\mathcal{F}_n} = o_{\scriptscriptstyle{G_0}}(1). \numberthis \label{III.eq: second to show 0}
    \end{align*}
    Denote for $k_n^x(u) = \frac{\ind_{\eta_n^* \leq u - x \leq \eta_n}}{\pi(u-x)}$, its fractional $1/2$ integral by $(I^{1/2} k_n^x)(z) = \int_z^\infty \frac{k_n^x(u)}{\sqrt{u-z}} \, du$ and analogously $(I_{\scriptscriptstyle{F/f_0}}^{1/2} k_n^x)(z) = \int_z^\infty \frac{k_n^x(u)}{\sqrt{u-z}} \frac{dF(u)}{f_0(u)}$. Then:
    \begin{align*}
        &\eqref{III.eq: first to show 0} = r_n \left\| \int_0^L (g_{\scriptscriptstyle{F}} - g_0)(z) \left( \frac{\ind_{x \leq z \leq x + \eta_n}}{\sqrt{z-x}} - \frac{1}{g_{\scriptscriptstyle{F}}(z)} (I_{\scriptscriptstyle{F/f_0}}^{1/2} k_n^x)(z) \right) \right\|_{\mathcal{F}_n} \\
        &\, \,  \leq r_n \left\| \int_0^L (g_{\scriptscriptstyle{F}} - g_0)(z) \frac{1}{\sqrt{z-x}} \ind_{x \leq z \leq x + \eta^*_n} dz \right\|_{\mathcal{F}_n}  \numberthis \label{III.eq: A} \\
        &\, \, \, + r_n \left\| \int_0^L \frac{\ind_{\eta^*_n \leq z - x \leq \eta_n}}{\sqrt{z-x}}  \left( g_0 - g_{\scriptscriptstyle{F}} \right)(z) \left( 1 - \frac{g_0}{g_{\scriptscriptstyle{F}}} (z) \right) \, dz \right\|_{\mathcal{F}_n} \numberthis \label{III.eq: B}\\
        &\, \, \, + r_n \left\| \int_0^L \left( 1- \frac{g_0}{g_{\scriptscriptstyle{F}}}(z)  \right) \left( \frac{g_0(z)}{\sqrt{z-x}} \ind_{\eta^*_n \leq z - x \leq \eta_n} - (I_{\scriptscriptstyle{F/f_0}}^{1/2} k_n^x)(z) \right) \, dz \right\|_{\mathcal{F}_n} \numberthis \label{III.eq: C}
    \end{align*}
    
    We consider first \eqref{III.eq: A}. Using the substitution $vs - x = y(s-x)$ thus $(1-v) = (1-y) \frac{s-x}{s}$:
    \begin{align*}
        \int_x^{x+\eta^*_n} &\frac{1}{\sqrt{z-x}} g_{0}(z) dz = \int_0^L \int_{0 < vs- x < \eta^*_n} \frac{\ind_{0 < v < 1}}{\sqrt{vs-x} \sqrt{1-v}} \, dv \, dF_0(s)\\ 
        &= \int_0^{x+\eta^*_n} \int_0^{\frac{\eta^*_n}{s-x} \wedge 1}  \frac{s-x}{s\sqrt{(y(s-x))}  } \frac{\sqrt{s}}{2\sqrt{(1-y)(s-x)}}  \, dy \, dF_0(s) \\
        &= \int_x^{x+\eta^*_n} \frac{1}{\sqrt{s}} \int_0^{\frac{\eta^*_n}{s-x} \wedge 1} \frac{1}{2 \sqrt{y(1-y)}} \, dy \,  dF_0(s) \leq \frac{\pi}{2} \int_x^{x+\eta^*_n} \frac{f_0(s)}{\sqrt{s}}  \, ds \\
        & \leq \pi \left\| f_0 \right\|_\infty (\sqrt{x+\eta^*_n} - \sqrt{x}) \leq \pi \sqrt{\eta^*_n} \left\| f_0 \right\|_\infty 
    \end{align*}
    where we used that: $\mathrm{Be}(\frac{\eta^*_n}{s-x} \wedge 1 ; \frac{1}{2}, \frac{1}{2}) = -\mathrm{arcsin} \sqrt{1 - \frac{\eta^*_n}{s-x} \wedge 1} + \frac{\pi}{2} \leq \frac{\pi}{2}$ and that since $x \mapsto \frac{1}{\sqrt{x}}$ is decreasing, by the integral test for series:
    \begin{align*}
        \int_x^{x+\eta^*_n} \frac{dF(s)}{2 \sqrt{s}}  &= \hspace{-0.2 cm} \sum_{i \, : \, 0 < x_i - x < \eta^*_n} \frac{F\{x_i\}}{2 \sqrt{x_i}}  \leq \hspace{-0.2 cm}\sum_{i \, : \, 0 < x_i - x < \eta^*_n} \frac{1}{2 \sqrt{x_i}} \leq \int_x^{x+\eta^*_n} \frac{ds}{2 \sqrt{s}} \leq \sqrt{\eta^*_n}
    \end{align*}
    thus $\eqref{III.eq: A} = O(r_n \sqrt{\eta^*_n}) =o(1)$.

    Now note that because $\sup_{z \in [x, x+\eta^*_n]} \big| \big( 1 + \frac{g_0}{g_{\scriptscriptstyle{F}}} (z) \big)^2  \big| < \infty$ uniformly for $F \in \mathcal{F}_1$ and using the contraction rate given in Theorem \ref{III.thm: contraction rate hellinger} for the Hellinger distance: 
    \begin{align*}
        \eqref{III.eq: B} &\leq \frac{r_n}{\eta^*_n} \sup_{F \in \mathcal{F}_1} \sup_{z \in [x, x+\eta^*_n]} \left| \left( 1 + \frac{g_0}{g_{\scriptscriptstyle{F}}}(z) \right)^2  \right| \sup_{F \in \mathcal{F}_n} h^2(g_{\scriptscriptstyle{F}}, g_0) = o(1).
    \end{align*}
    
    The quantity in \eqref{III.eq: C} can be upper bounded by the following quantity:
    \begin{align*}
        &r_n \left( \int_0^{x+\eta_n} \left| \frac{g_0}{g_{\scriptscriptstyle{F}}}(z) - 1 \right|^2 \, dz \right)^{\frac{1}{2}} \left( \int_0^L \left( \frac{\ind_{\eta^*_n \leq z - x \leq \eta_n}}{\sqrt{z-x}} g_0(z) - I_{\scriptscriptstyle{F/f_0}}^{1/2} k_n^x(z) \right)^2 dz \right)^{\frac{1}{2}} 
    \end{align*}
    and by using Lemma 8 from \cite{27}:
    \begin{align*}
    \left( \int \left( I^{1/2} k_n^x(z) \right)^2 dz \right)^{1/2} &\leq \left( \int \left( \frac{1}{\sqrt{z-x}} \ind_{\eta^*_n \leq z - x \leq \eta_n} \right)^2 dz \right)^{1/2} \hspace{-0.2cm} + O(\eta_n^{-1}) \\
    &= O \left( \sqrt{\log \frac{\eta_n}{\eta^*_n}} + \eta_n^{-1} \right).
    \end{align*}
    Moreover since the function $u \rightarrow \frac{1}{\sqrt{u-z}} \frac{1}{u-x} \ind_{\eta^*_n \leq u - x \leq \eta_n}$ is decreasing and nonnegative in $u$, we can thus use the integral test for series thus $\forall \, z \in [0,L]$: $\big| I_{\scriptscriptstyle{F/f_0}}^{1/2} k_n^x(z) \big| \leq \left| I^{1/2} k_n^x(z) \right|$ and conclude:
    \begin{align*}
        \left( \int \left( I_{\scriptscriptstyle{F/f_0}}^{1/2} k_n^x(z) \right)^2 dz \right)^{1/2} = O \left( \sqrt{\log \frac{\eta_n}{\eta^*_n}} + \eta_n^{-1} \right).
    \end{align*}    
    Finally since uniformly for $F \in \mathcal{F}_1$, $\sup_{z \in [0, x+\eta_n]} g^{-1/2}_{\scriptscriptstyle{F}}(z) < \infty$ and similarly $\sup_{z \in [0, x+\eta_n]} \big( \big( 1 + \frac{g_0}{g_{\scriptscriptstyle{F}}} (z)  \big)^2 \big)^{1/2} < \infty$, we conclude:
    \begin{align*}
        \eqref{III.eq: C} \lesssim r_n  h(g_{\scriptscriptstyle{F}}, g_0) \, O \left( \sqrt{\log \frac{\eta_n}{\eta^*_n}} +\eta_n^{-1} \right). 
    \end{align*}

    We now prove \eqref{III.eq: second to show 0}. We consider the class of functions:
    \begin{align*}
        \overline{\mathcal{V}}_n = \left\{ z \mapsto \frac{\ind_{0 \leq z \leq x + \eta_n}}{g_{\scriptscriptstyle{F}}(z)} \int_z^L \frac{1}{\sqrt{u-z}}  \frac{\ind_{\eta_n^* \leq u - x \leq \eta_n}}{\pi (u-x)}  \frac{dF(u)}{f_0(u)} \, \, : \, \, F \in \mathcal{F}_n \right\}
    \end{align*}
    where $D^* = \sup_{F \in \mathcal{F}_1} \int \frac{\ind_{\eta_n^* \leq u - x \leq \eta_n}}{(u-x) \pi} \frac{dF(u)}{f_0(u)}  = O \big( \log \frac{\eta_n}{\eta_n^*} \big)$, then $\mathcal{J} (1, \overline{\mathcal{V}}_n, V_n^*) < \infty$, and $V_n^{*}$ is the maximizer over support of the DP of the functions $V_n^{\scriptscriptstyle{F}} \in \overline{\mathcal{V}}_n$.
    
    Here comes in the importance of the localization argument because now the function: $z \mapsto \int_z^L \frac{\ind_{\eta_n^* \leq u - x \leq \eta_n}}{\pi(u-x) \sqrt{u-z}} \frac{dF(u)}{f_0(u)}$ is $0$ for all $z \geq x + \eta_n$ by definition; this means that we can use: $\sup_{F \in \mathcal{F}_1} \sup_{z \in [0, x+\eta_n]} g^{-1}_{\scriptscriptstyle{F}}(z) < \infty$ if $x < L$.
    
    The class $\overline{\mathcal{V}}_n$ is the product of two classes:
    \begin{align*}
        \mathcal{U}_n = \left\{ z \mapsto \frac{1}{D^*} \int_{z}^L \frac{1}{\sqrt{u-z}} \frac{\ind_{\eta_n^* \leq u - x \leq \eta_n}}{\pi (u-x)}  \frac{dF(u)}{f_0(u)} \, : \, F \in \mathcal{F}_n \right\}
    \end{align*}
    and $\phi(\mathcal{G}_n)$ where $g \mapsto \phi(g)$ is the mapping $g(z) \mapsto g^{-1}(z) \ind_{[0, x+\eta_n]}(z)$ and
    \begin{align*}
        \mathcal{G}_n = \left\{ z \mapsto g_{\scriptscriptstyle{F}}(z) = \sum_{i=1}^{\infty} \frac{F\{x_i\}}{ \sqrt{x_i^2 - x_i z}} \, : \,  F \in \mathcal{F}_n \right\}.
    \end{align*}
    
    We study for $C_n = [0, x + \eta_n]$ and $U_n^*$ the maximizer over the support of the DP of the functions $U_n^{\scriptscriptstyle{F}} \in \mathcal{U}_n$, analogously $g_*$ means the maximizer over the support of the DP of the functions $g_{\scriptscriptstyle{F}}$:
    \begin{align*}
        &\int_0^1 \sup_Q \sqrt{ \log N(\varepsilon \| \overline{V}_n^* \|_{Q, 2}, \overline{\mathcal{V}}_n, \mathbb{L}_2(Q)) } \, d\varepsilon \\
        &\quad \quad \leq \int_0^1 \sup_Q \sqrt{ \log N \left( \frac{\varepsilon}{2} \left\| \frac{\ind_{C_n}}{g_*} \right\|_{Q, 2}, \phi(\mathcal{G}_n), \mathbb{L}_2(Q) \right) } d\varepsilon \\
        &\quad \quad \quad \quad  \quad \quad \quad \quad + \int_0^1 \sup_Q \sqrt{ \log N \left( \frac{\varepsilon}{2} \| U_n^{\scriptscriptstyle{F}\star} \|_{Q, 2}, \mathcal{U}_n, \mathbb{L}_2(Q) \right) } d\varepsilon.
    \end{align*}

    Take now the class of functions: $\mathcal{S} := \left\{ z \mapsto \frac{\ind_{x > z}}{2 \sqrt{x^2 - xz}} \: : \: x \in \mathbb{R}^{+} \right\}$. We prove that this class of functions is VC. It is immediate using Lemma 2.6.16 and 2.6.18 from \cite{5} that the class of functions $\left\{ z \mapsto x(x-z) \: : \: x \in \mathbb{R}^+ \right\}$ is VC. Next by Lemma 2.6.18 (ii) in \cite{5} we see that the class $\left\{ z \mapsto x(x-z)_{+} \: : \: x \in \mathbb{R}^+ \right\}$ is VC. Again by Lemma 2.6.18 (viii) in \cite{5} we see that the class of functions $\left\{ z \mapsto \left( x(x-z)_{+} \right)^{-1/2} \: : \: x \in \mathbb{R}^+ \right\}$ is VC, as the square root of a function is a monotone function. Note that taking the reciprocal of a function is monotone except for one point, which is the point where the original function assumes value 0, and in that case we assign to the reciprocal also value 0. This makes the reciprocal function monotone except at one point. We can view the resulting class as the union of two VC classes (one obtained by taking the inverse except at zero and the other assigning value zero when the original function gets value zero), thus still VC. This proves $\mathcal{S}$ is VC.
    
    Since in $\mathcal{F}_n$ are contained realizations of the Dirichlet process with probability one, and these are almost surely discrete, we can look at the infinite mixtures of functions belonging to the set $\mathcal{S}$. We can therefore define $\mathcal{G}_n := \{z \mapsto g_{\scriptscriptstyle{F}} (z) = \sum_{i=1}^{\infty} \frac{\ind_{x_i > z}}{2 \sqrt{x_i^2 - x_i z}} F(x_i) \: : \: F \in \mathcal{F}_n \} $. Then $\mathcal{G}_n$ is a VC-hull of the class $\mathcal{S}$, as we showed previously that $\mathcal{S}$ is VC. Using Corollary 2.6.12 from \cite{5} we have that for a constant $K$ that depends only on the VC-index $V_m(\mathcal{G}_n)$:
    \begin{align}\label{III.eq: finiteness entropy}
        \log{N\left(\varepsilon \|  g_* \|_{Q,2}, \mathcal{G}_n, \mathbb{L}_2(Q) \right)} \leq K \left(\frac{1}{\varepsilon} \right)^{2 - 2 V^{-1}_m(\mathcal{G}_n)}
    \end{align}
    Now consider the functions $f \mapsto \phi(f) := \frac{1}{f}$, then using \eqref{III.eq: finiteness entropy} we have:
    \begin{align*}
        &\int_{0}^{1} \sup_{Q} \sqrt{\log{N\left(\varepsilon \| \ind_{C_n} g^{-1}_{*} \|_{Q,2}, \phi(\mathcal{G}_n), \mathbb{L}_2(Q)  \right)}} \, d\varepsilon \\
        &\leq \int_{0}^{1} \sup_{Q} \sqrt{\log{N\left(\varepsilon \| g_* \|_{Q,2}, \mathcal{G}_n, \mathbb{L}_2(Q)  \right)}} \, d\varepsilon < \infty \numberthis \label{III.eq: bounded entropy densities}
    \end{align*}
    where in the last inequality we used the fact that we can use Theorem 2.10.20 from \cite{5} as $\sup_{z \in [0, x+\eta_n]} | g_{\scriptscriptstyle{F}}^{-4}(z) | < \infty$ uniformly over $\mathcal{F}_1$, and thus:    
    \begin{align*}
        |\phi \cdot g_{\scriptscriptstyle{F}_1}(z) - \phi \cdot g_{\scriptscriptstyle{F}_2}(z)|^2 \leq \big\| \sup_{z \in [0, x+\eta_n]} | g_{\scriptscriptstyle{F}}^{-4}(z) | \big\|_{\mathcal{F}_1} |g_{\scriptscriptstyle{F}_1}(z) - g_{\scriptscriptstyle{F}_2}(z)|
    \end{align*}

    We show:
    \begin{align}\label{III.eq: eq to verify}
    J(1, \mathcal{U}_n, U_n^{*}) = \int_{0}^{1} \sup_{Q} \sqrt{\log{N\left(\varepsilon \| U_n^{*}\|_{Q,2},\mathcal{U}_n, \mathbb{L}_2(Q)  \right)}} \, d\varepsilon < \infty.
    \end{align}
    Note that we can write any function $U^{\scriptscriptstyle{F}}_n \in \mathcal{U}_n$ as: $z \mapsto   \frac{1}{D_{*}}\sum_{u_i >z} \frac{w(u_i)}{\sqrt{u_i-z}}  \, F\{u_i\}$ for weights $w(u_i)$:
    $$ w(u_i) := \frac{1}{f_0(u_i)}\frac{\ind_{\eta_n^* \leq u_i - x \leq \eta_n}}{\pi (u_i-x)}, \quad \text{and} \quad D_{*} = \sup_{F \in \mathcal{F}_1}\sum_i |w(u_i)| F(u_i)$$
    where $\mathcal{F}_1$ denotes the weak support of the Dirichlet process prior. Note $D_{*}$ is an appropriate renormalizing constant and that $D_{*} < \infty$. Because the $\mathbb{L}_1$ norm of the weights is bounded for every $n$, we conclude that the weights are absolutely summable in the sense: $(D^*)^{-1} \sum |w(u_i)| F\{u_i\} \leq 1 $ for every $n$. \eqref{III.eq: eq to verify} now follows from the same reasoning used for $\mathcal{G}_n$. Indeed the class of functions $\left\{ z \mapsto ((u-z)_+)^{-1/2} \: : \: u \in \mathbb{R}^+ \right\}$ is VC, thus taking its convex hull, using Corollary 2.6.12 from \cite{5} and the fact that the weights are absolutely summable, gives \eqref{III.eq: eq to verify}.
    
    We can now conclude by using the maximal inequality:
    \begin{align*}
        &\pb_0 \left( r_n \left\| \frac{1}{n} \sum_{i=1}^n Y_{n,i}^{\scriptscriptstyle{F}} - E_0 Y_{n,i}^{\scriptscriptstyle{F}} \right\|_{\mathcal{F}_n} > \varepsilon \right) \\
        &= \pb_0 \left( \sqrt{n} \left\| \frac{1}{n} \sum_{i=1}^n \overline{V}_{n,i}^{\scriptscriptstyle{F}} - \ex_{0} \overline{V}_{n,i}^{\scriptscriptstyle{F}} \right\|_{\mathcal{F}_n} \hspace{-0.2cm} > \frac{\varepsilon \sqrt{n}}{r_n D^*} \right) \leq \frac{ r_n}{\varepsilon \sqrt{n}} O  \left( \sqrt{\log \frac{1}{\eta_n^*}} + \frac{1}{\eta_n} \right), 
    \end{align*}
    because
    \begin{align*}
        \| V_n^* \|_{G_0, 2}^2 &= \frac{1}{(D^*)^2} \int_0^L \frac{g_{0}(z)}{g_{\scriptscriptstyle{F}}^2(z)} \left( \int_z^L \frac{1}{\sqrt{u-z}} \frac{\ind_{\eta_n^* \leq u - x \leq \eta_n}}{(u-x)} \frac{dF(u)}{f_0(u)} \right)^2 dz \\
        &= \frac{1}{(D^*)^2} \left( O \left( \log \frac{1}{\eta_n^*} \right) + O(\eta_n^{-2}) \right).
    \end{align*}
    The last display converges to zero under $G_0$ proving \eqref{III.eq: second to show 0}.

    The last step of the proof consists in showing $\forall \, t$ in a neighborhood of $0$:
    \begin{align}\label{III.eq: invariance measure condition}
        \ddfrac{\int_{\mathcal{F}_n} \prod_{i=1}^n g_{\scriptscriptstyle{F_{\scaleto{t_n}{3.5pt}}}}(Z_i) \: d \Pi(F)}{\int_{\mathcal{F}_n} \prod_{i=1}^n g_{\scriptscriptstyle{F}}(Z_i) \: d \Pi(F)} \stackrel{G_0}{\longrightarrow} 1.
    \end{align}
    This condition entails the so-called "invariance" under changing the prior from $F \sim \Pi$ (in the denominator) to $F \circ \phi^{-1}$ (in the numerator) for $\phi \, : \, g_{\scriptscriptstyle{F}} \mapsto g_{\scriptscriptstyle{F_{\scaleto{t_n}{3.5pt}}}}$.

    If $F \sim \operatorname{DP}(\alpha)$, $F$ is the normalized Gamma process, a \textit{completely random measure}. So if $F = \frac{\Gamma}{|\Gamma|}$ the:
    \begin{align}\label{III.eq: perturbation as crm}
        dF_t = \frac{(1+th^x) \, dF}{\int (1+th^x) \, dF} = \frac{(1+th^x) \, d \Gamma}{\int (1+th^x) \, d \Gamma} 
    \end{align}
    Hence $F_t$ is the normalized process $\overline{\Gamma}$ defined, for any measurable $A$, by:
    \begin{align*}
        \overline{\Gamma}(A) = \int_A (1+th^x) \, d \Gamma
    \end{align*}
    Since $F$ has base measure $\alpha$, we prove that the process $\overline{\Gamma}$ has intensity measure:
    \begin{align*}
        \nu (d s, dv)=s^{-1} \exp \left\{-\frac{s}{(1+t h^x(v))} \right\} d s \, d \alpha(v)
    \end{align*}
    A CRM is characterized by its intensity measure. For a CRM $\overline{\Gamma}$ without fixed atoms, this is the measure on $(0, \infty) \times \mathfrak{X}$ so that:
    $$
    -\log \mathrm{E} \, e^{-\int f d \overline{\Gamma}}=\int_{\mathfrak{X}} \int_0^{\infty}\left(1-e^{-s f(v)}\right) \nu (d s, d v)
    $$
    for any measurable $f$ which is integrable with respect to $\overline{\Gamma}$ almost surely.
    For the Gamma Process $v(d s, d v)=s^{-1} e^{-s} \, d s \, d \alpha(v)$. Because $\int f d \bar{\Gamma}=\int f(1+s h) d \Gamma$ it follows, by substituting $s(1+t h)=u$:
    \begin{align*}
    &-\log \mathrm{E} \, e^{-\int f d \bar{\Gamma}}=\iint\left(1-e^{-s f(v)(1+t h^x(v))}\right) s^{-1} e^{-s} \,  d s \, d \alpha(v) \\
    & =\iint\left(1-e^{-u f(v)}\right) \frac{(1+t h^x(v))}{u} \exp{\left\{-\frac{u}{(1+th^x(v))} \right\} } \frac{d u}{1+t h^x(v)} \,  d \alpha(v) \\
    & =\iint\left(1-e^{-s f(v)}\right) s^{-1} \exp{\left\{-\frac{s}{(1+th^x(v))} \right\} } \, d s \, d \alpha(v)
    \end{align*}
    which proves the claim.

    Now for $X_1, \ldots, X_n \mid F \sim F_t$, for any measurable functions $\varphi_1, \ldots, \varphi_n$ (taken to be $\varphi\left(Z_i, \cdot\right)$ in our case):
    \begin{align*}
    \mathrm{E} \bigg( \prod_{i=1}^n \varphi_i \left( X_i \right) \bigg)=\mathrm{E} \, \mathrm{E}\bigg(\prod_{i=1}^n \varphi_i\left(X_i\right) \mid F\bigg)= \mathrm{E} \prod_{i=1}^n \int \varphi_i \, d F_t  
    \end{align*}
    therefore in our case, the functions $\varphi\left(Z_i, \cdot\right)$ are given as follows:
    \begin{align*}
        &\ddfrac{\int_{\mathcal{F}_n} \prod_{i=1}^n g_{\scriptscriptstyle{F_{\scaleto{t_n}{3.5pt}}}}(Z_i) \: d \Pi(F)}{\int_{\mathcal{F}_n} \prod_{i=1}^n g_{\scriptscriptstyle{F}}(Z_i) \: d \Pi(F)} =\ddfrac{\int_{\mathcal{F}_n} \prod_{i=1}^n \int_{Z_i}^{\infty} \frac{dF_{t_n}(u)}{2\sqrt{(u)^2 - u Z_i }} \: d \Pi(F)}{\int_{\mathcal{F}_n} \prod_{i=1}^n \int_{Z_i}^{\infty} \frac{dF(u)}{2\sqrt{(u)^2 - u Z_i }} \: d \Pi(F)} \\
        &=\ddfrac{\ex \left[ \prod_{i=1}^n  \frac{\ind_{X_{i,n} > Z_i}}{2\sqrt{(X_{i,n})^2 - X_{i,n} Z_i }} \right]}{\ex \left[ \prod_{i=1}^n \frac{\ind_{X_i > Z_i}}{2\sqrt{(X_i)^2 - X_i Z_i }} \right]} \numberthis \label{III.eq: to rewrite using lemmas 1}
    \end{align*}
    Let $\mathscr{S}_n$ stand for the collection of all partitions of $\{1, \ldots, n\}$ that emerge from the ties obtained by sampling from a Dirichlet process, and write a partition as $\mathcal{S}=\left\{S_1, \ldots, S_N\right\}$, for $N=\# \mathcal{S}$ the number of partitioning sets. Then using Proposition 4.7 from \cite{4} the denominator of \eqref{III.eq: to rewrite using lemmas 1} can be written as:
    \begin{align*}
        \frac{1}{|\alpha|^{[n]}} \sum_{\mathcal{S} \in \mathscr{S}_n} \prod_{j=1}^{\# \mathcal{S}}\left[\left(\# S_j-1\right) ! \int \prod_{l \in S_j} \varphi_l(v) d \alpha(v)\right]. \numberthis \label{III.eq: denominator 1}
    \end{align*}
    Now we are going to give both an upper bound and lower bound for the numerator in \eqref{III.eq: to rewrite using lemmas 1} and we are going to show that both the upper bound and the lower bound, divided by \eqref{III.eq: denominator 1}, go to 1.

    For $X_1, \ldots, X_n \mid F \sim F_t$, for any measurable functions $\varphi_1, \ldots, \varphi_n$: 
    \begin{align*}
    &\mathrm{E} \prod_{i=1}^n \varphi_i\left(x_i\right) = \numberthis \label{III.eq: relation to prove}\\
    &\sum_{\mathcal{S} \in \mathscr{S}_n} \int \cdots \int \prod_{j=1}^k \prod_{l \in S_j} \varphi_l\left(v_j\right) \int_0^{\infty} \hspace{-0.05cm} \frac{\lambda^{n-1}}{\Gamma(n)} e^{-\psi(\lambda)} \prod_{j=1}^k \tau_{\scriptscriptstyle{\left|S_j\right|}}\left(\lambda \mid v_j\right) d \lambda d \alpha\left(v_1\right) \cdots d \alpha\left(v_k\right)
    \end{align*}
    where:
    \begin{align*}
        &\psi(\lambda) = \int \log{(1+\lambda(1+th^x(v)))} \, d\alpha(v) \\
        &\tau_{\scriptscriptstyle{\left|N\right|}}\left(\lambda \mid v\right) = \Gamma(N) \left(\frac{1+t h^x(v)}{1+ \lambda(1+th^x(v))} \right)^N
    \end{align*}
    The idea is that the realizations of $F_t$ are discrete, whence $X_1, \ldots, X_n$ can be grouped in sets of equal values. The joint distribution of the random partitions and values is given by Lemma 14.62 in \cite{4}. In order to prove \eqref{III.eq: relation to prove} by using  Lemma 14.62 in \cite{4} we need to compute:
    \begin{align*}
    \psi(\lambda) & =\iint\left(1-e^{-\lambda s}\right) \nu(d v, d s)= \\
    & =\iint\left(1-e^{-\lambda s}\right) s^{-1} e^{-\frac{s}{1+th^x(v)}} \, d s \, d \alpha(v) \\
    & =\iint\left(1-e^{-\lambda(1+th^x(v)) s}\right) s^{-1} e^{-s} \, d s \, d \alpha(v)= \\
    & =\int \log (1+\lambda(1+th^x(v))) \, d \alpha(v)
    \end{align*}
    where by example J.7 in \cite{4}, for $\nu^c (dv,ds) = s^{-1} e^{-s} \, ds \, d \alpha(v)$ we have for $\theta >0$:
    \begin{align*}
        \iint (1- e^{-\theta s f(v)}) \, \nu^{c}(dv,ds) = \int \log{(1+ \theta f(v))} \, d\alpha(v).
    \end{align*}
    Moreover using $\int_{0}^{\infty} s^{n-1} e^{-(c+k)s} \, ds = \Gamma(n)(c+k)^{-n}$:
    \begin{align*}
        \tau_{\scriptscriptstyle{N}}(\lambda \mid v) &= \int s^N e^{-\lambda s} \, \rho(ds \mid v) = \int s^N e^{-\lambda s} s^{-1} \exp{\left\{- \frac{s}{(1+th^x(v))} \right\} }\, ds\\
        &= \Gamma(N)\frac{1}{(\lambda + \frac{1}{(1+th^x(v))})^N}.
    \end{align*}
    We conclude that \eqref{III.eq: relation to prove} holds true and as a special case for $t=0$ we retrieve $F_t= F \sim \operatorname{DP}(\alpha)$, which gives $\psi(\lambda) = (\log{(1+\lambda)})|\alpha|$ and $\tau_{\scriptscriptstyle{N}} (\lambda \mid v)= \Gamma(N)/(1+\lambda)^{N}$ and it is possible to prove that \eqref{III.eq: relation to prove} reduces to \eqref{III.eq: denominator 1}.

    Now for $t>0$ small, using that $\frac{1}{1 + \lambda(1-r)} \leq \frac{1+r}{(1-r)(1+\lambda(1+r))}$we have: 
    \begin{align*}
    & \int_0^{\infty} \lambda^{n-1} e^{-\psi(\lambda)} \prod_{j=1}^k \frac{\tau_{\scriptscriptstyle{\left|S_j\right|}}\left(\lambda \mid v_j\right)}{\Gamma\left(\left|S_j\right|\right)} d \lambda= \\
    & =\int_0^{\infty} \lambda^{n-1} e^{-\int \log (1+\lambda(1+t h^x(v))) \, d \alpha (v) } \prod_{j=1}^k\left(\frac{1+t h^x\left(v_j\right)}{1+\lambda\left(1+t h^x\left(v_j\right)\right)}\right)^{\left|S_j\right|} d \lambda \\
    & \leq \int_0^{\infty} \lambda^{n-1} \frac{1}{\left(1+\lambda\left(1-t\|h\|_{\infty}\right)\right)^{|\alpha|}}\left(\frac{1+t\|h\|_{\infty}}{1+\lambda\left(1+t\|h\|_{\infty}\right)}\right)^n d \lambda \\
    &\leq \left(\frac{1+t\|h\|_{\infty}}{1-t\|h\|_{\infty}}\right)^{|\alpha|} \int_0^{\infty} \frac{\lambda^{n-1}\left(1+t\|h\|_{\infty}\right)^n}{\left(1+\lambda\left(1+t\|h\|_{\infty}\right)\right)^{|\alpha|+n}} d \lambda \\
    &= \left(\frac{1+t\|h\|_{\infty}}{1-t\|h\|_{\infty}}\right)^{|\alpha|} \frac{\Gamma(|\alpha|) \Gamma(n)}{\Gamma(|\alpha| + n)}.
    \end{align*}
    The lower bound is given similarly by:
    \begin{align*}
        \int_0^{\infty} \lambda^{n-1} e^{-\psi(\lambda)} \prod_{j=1}^k \frac{\tau_{\scriptscriptstyle{\left|S_j\right|}}\left(\lambda \mid v_j\right)}{\Gamma\left(\left|S_j\right|\right)} d \lambda \geq \left(\frac{1-t\|h\|_{\infty}}{1+t\|h\|_{\infty}}\right)^{|\alpha|} \frac{\Gamma(|\alpha|) \Gamma(n)}{\Gamma(|\alpha| + n)}.
    \end{align*}
    Therefore we can upper-bound (and similarly lower-bound) \eqref{III.eq: relation to prove} by: 
    \begin{align*}
        \left(\frac{1+t\|h\|_{\infty}}{1-t\|h\|_{\infty}}\right)^{|\alpha|} \frac{1}{|\alpha|^{[n]}} \sum_{\mathcal{S} \in \mathscr{S}_n} \prod_{j=1}^{\# \mathcal{S}}\bigg[ \int \prod_{l \in S_j} \varphi_l(v) \underbrace{\prod_{i}^{N} \Gamma(|S_j|)}_{=\left(\# S_j-1\right) !} d \alpha(v)\bigg] . 
    \end{align*}
    because $t$ was chosen to depend on $n$ in the way: $t_n = t r_n / n$ and $\| h^x_n\|_{\infty} = O(\sqrt{n})$, we conclude that: $ (1 -O(1/\sqrt[4]{n}))^{|\alpha|}\leq\eqref{III.eq: to rewrite using lemmas 1} \leq (1 +O(1/\sqrt[4]{n}))^{|\alpha|}$ which proves \eqref{III.eq: invariance measure condition}. \begin{flushright} \qed \end{flushright}    
\end{proof}




\begin{proof}[\hypertarget{III.proof contraction wasserstein}{\textbf{Proof of contraction rate in Wasserstein metric, Theorem} \ref{III.thm: contraction rate F}}] 
    We follow the structure of an analogous result given in \cite{25}. The first thing is to upper bound the Wasserstein distance in terms of the Hellinger distance and this is done in Lemma \ref{III.lemma: wasserstein bound} below. 

    We establish the contraction rate for the "discrepancy" $d(g_{\scriptscriptstyle{F}},g_0)=\gamma(h(g_{\scriptscriptstyle{F}},g_0))$ for $\varepsilon \mapsto \gamma(\varepsilon) = D_k \varepsilon^{1/(k+2)}[\log{(C_k/\varepsilon)}]^{\frac{k+1/2}{k+2}}$. Note that $\gamma^{-1} (d) \leq h(g_{\scriptscriptstyle{F}},g_0)$ and it possesses the properties (a)-(d) in Section 9 of \cite{25} and the balls of $d$ are convex. We can therefore apply Theorem 3 from \cite{25} to obtain a rate of posterior contraction relative to $d$ and thus relative to $ h^{1/(k+2)} (\log(1/h))^{\frac{k+1/2}{k+2}}$. We apply the theorem with the same set $\mathcal{P}_n$  as in the proof of Theorem \ref{III.thm: contraction rate hellinger}. Thus by Proposition \ref{III.prop: complexity bound model partition}:
   \begin{align*}
       \log{N(\varepsilon, \mathcal{P}_n, d)} = \log{N \bigg(\varepsilon^{\frac{1}{k+2}} \left( \log{\frac{1}{\varepsilon}} \right)^{\frac{k+1/2}{k+2}}, \mathcal{P}_n, h \bigg)} \lesssim \left( \frac{1}{\varepsilon}\right)^{\frac{2}{k+2}} \hspace{-0.1cm} \left( \log{\frac{1}{\varepsilon}} \right)^{\frac{2k+1}{k+2}}.
   \end{align*}
   Therefore condition (9.1) from Theorem 3 in \cite{25} holds for the rate $\varepsilon_n \gtrsim n^{-\gamma}$ and any $\gamma < (k+2)/(2k+4)$. The prior mass condition proved in Proposition \ref{III.prop: bound of kullback neighborhood probability} is satisfied with the rate $\varepsilon_n \asymp (\log n/n)^{1/4}$. The slowest of the two rates is $ \varepsilon_n \asymp (\log n/n)^{1/4}$ which, by substitution, yields a rate:
   \begin{align*}
       \varepsilon_n^{2/(k+2)} (\log{(1/\varepsilon_n)} )^{\frac{2k+1}{k+2}} \asymp n^{-1 /(2 k + 4)}(\log n)^{(4k +3)/(2k+4)}.
   \end{align*}
Theorem \ref{III.thm: contraction rate V} gives the contraction rate of $|V_{\scriptscriptstyle{F}}(0) - V_0(0)|$ and it is faster than the contraction rate for $d$. The latter is thus the dominating contraction rate. 
\end{proof}

\begin{lemma}\label{III.lemma: wasserstein bound}
For probability measures $F$ and $F_0$ supported on $[0, L]$, and $g_{\scriptscriptstyle{F}}= \int \psi \, dF$ for $\psi$ density as in \eqref{eq: g_F_0 expression} and for $x \geq 0$, $F^*(x) = 1 - V_{\scriptscriptstyle{F}}(x) / V_{\scriptscriptstyle{F}}(0)$ and similarly $F^*_0(x)$ as in \eqref{eq: solution range equation} with $V_0(0)>0$, then for any $k \geq 1$,
$$
W_k\left(F^*, F^*_0 \right) \lesssim h\left(g_{\scriptscriptstyle{F}}, g_0\right)^{1 /(k+2)}\bigg(\log \frac{C_k}{h\left(g_{\scriptscriptstyle{F}}, g_0\right)}\bigg)^{\frac{(k+1/2)}{(k+2)}} + |V_{\scriptscriptstyle{F}}(0) - V_0(0)|.
$$
where $C_k$ is a constant such that the map $\varepsilon \mapsto \varepsilon\left[\log \left(C_k / \varepsilon\right)\right]^{k+1 / 2}$ is monotone on $(0,2]$.
\end{lemma}
\begin{proof}
    The proof of this lemma follows very closely the proof of Lemma 7 in \cite{25}, with some adaptations for our case. The strategy consists of approximating $F^*$ and $F^*_0$ by convolving them with a mollifier. Using Theorem 6.15 from \cite{31}, $W_k\left(F^*, F^*_0 \right) \lesssim \left( \int |x|^k \, d |F^* - F^*_0| (x)\right)^{1/k}$, where $|F^* - F^*_0|$ is the total variation measure of $F^* - F^*_0$. For a measure $\mu = \mu^+ - \mu^-$, $|\mu|=\mu^+ + \mu^-$, where respectively $\mu^+$ and $\mu^-$ are the positive and negative variation of $\mu$. We apply the above upper bound to the convolutions of $F^*$ and $F^*_0$ with the Gaussian mollifier $\phi_{\delta}(x)=\frac{1}{\delta \sqrt{2\pi}} e^{-\frac{x^2}{2 \delta^2}}$ (denote by $\Phi_{\delta}$ the normal distribution). For any $M>0$ and $Z$ standard normal variable:
    \begin{align*}
        &W_k\big(F * \Phi_\delta,\,  F_0 * \Phi_\delta\big)^k \lesssim \int|x|^k \,  d \left|\left(F^*-F^*_0\right) * \Phi_\delta  \right| (x) \\
        &\quad \leq \frac{|V_{\scriptscriptstyle{F}}(0) - V_0(0)|}{|V_{\scriptscriptstyle{F}}(0) V_0(0)|} \int |x|^k \, d|V_{\scriptscriptstyle{F}} * \Phi_{\delta}| (x) + \frac{1}{V_0(0)} \int |x|^k \, d \left|\left(V_{\scriptscriptstyle{F}}-V_0\right) * \Phi_\delta  \right| (x)\\
        &\quad \lesssim |V_{\scriptscriptstyle{F}}(0) - V_0(0)| + \int |x|^k \, d \left|\left(V_{\scriptscriptstyle{F}}-V_0\right) * \Phi_\delta  \right| (x)
    \end{align*}
    where we used that $V_{\scriptscriptstyle{F}}(0), V_0(0) >0$ and that $\int |x|^k \, d|V_{\scriptscriptstyle{F}} * \Phi_{\delta}| (x) < \infty, \, \forall \, \delta >0$. Now we analyze separately the second expression in the inequality above:
    \begin{align*}
        & \int |x|^k \, d \left|\left(V_{\scriptscriptstyle{F}}-V_0\right) * \Phi_\delta  \right| (x) =\int |x|^{k} \, \left| \left( V_0  - V_{\scriptscriptstyle{F}}  \right)  * \phi_\delta (x)  \right| \, dx\\
        &\quad \leq\left(\int_{-M}^M x^{2 k} d x \int_{-M}^M \left| \left( V_0  - V_{\scriptscriptstyle{F}}  \right)  * \phi_\delta (x)  \right|^2 d x\right)^{1 / 2} \\
        &\quad \quad+e^{-M} \int_{|x|>M}|x|^k e^{|x|} \left| \left( V_0  - V_{\scriptscriptstyle{F}}  \right)  * \phi_\delta (x)  \right| d x \\
        &\quad \lesssim M^{k+\frac{1}{2} }\left\| \left(V_0  - V_{\scriptscriptstyle{F}} \right)  * \phi_\delta  \right\|_2+e^{-M} \max\{1,e^{2|L|}\} \ex e^{2|\delta Z|},
    \end{align*}
    where we used that the measures $V_{\scriptscriptstyle{F}}$ and $V_0$ are supported on $[0,L]$, the fact for differentiable functions $f$ and $g$ in $\mathbb{L}_1$, $\frac{d}{dx} (f*g)(x) =  (\frac{d}{dx}f(x))*g(x)$ and the fact for a measure that is absolutely continuous with respect to the Lebesgue measure we have $\frac{d}{d\lambda}|\mu| = |\frac{d\mu}{d\lambda}|$, where $\lambda$ is the Lebesgue measure and $|\mu|$ the total variation measure. Note also that the number $\max\{1,e^{2|L|}\} \ex e^{2|\delta Z|}$ is uniformly bounded. The reason why we upper bounded the distance $ W_k\big(F^* * \Phi_\delta, F^*_0 * \Phi_\delta\big)$ with an expression in terms of $V_0 $ and $ V_{\scriptscriptstyle{F}}$ is that we can express the density of the observations $g_{\scriptscriptstyle{F}}$ as a convolution with respect to the measure $ V_{\scriptscriptstyle{F}}$: $g_{\scriptscriptstyle{F}} = k * V_{\scriptscriptstyle{F}} $ for $k(x) = -\ind_{\{x < 0 \}}/ \pi \sqrt{-x}$. This will allow us to use tools like the convolution theorem for Fourier transforms and Plancherel's theorem. We denote the Fourier transform $\mathcal{F}$ as follows $f \mapsto \mathcal{F}(f) = \hat{f}$. Then:
    $$
    \begin{aligned}
        &\left\|\left(V_0  - V_{\scriptscriptstyle{F}}  \right)  * \phi_\delta \right\|_2^2 = \int\left|\hat{V}_0-\hat{V}_{\scriptscriptstyle{F}}\right|^2(\lambda) \hat{\phi}_\delta^2(\lambda) d \lambda\\
        &=\int\left|\hat{f}(\hat{V}_0-\hat{V}_{\scriptscriptstyle{F}})\right|^2 (\lambda) \frac{\hat{\phi}_\delta^2}{|\hat{f}|^2}(\lambda) \, d \lambda = \left\|(g_{\scriptscriptstyle{F}}-g_0) * \phi_{\delta} * \mathcal{F}(\hat{f}^{-1}) \right\|_2^2 \\
        & \leq 4 d_{\scriptscriptstyle{TV}}^2 (g_{\scriptscriptstyle{F}},g_0) \,  \|\phi_{\delta} * \mathcal{F}(\hat{f}^{-1}) \|_2^2  \lesssim h^2\left(g_{\scriptscriptstyle{F}}, g_0\right)\delta^{-2},
    \end{aligned}
    $$
    where $d_{\scriptscriptstyle{TV}}$ is the total variation distance. Above we used a known fact (cf.\ \cite{32} pp.\ 391) that for an arbitrary (signed) measure $\mu$ on $\mathbb{R}$ and a function $g \in \mathbb{L}_2 (\mathbb{R})$, there holds $\| \mu * g \|_2 \leq |\mu| \| g\|_2$ ($|\mu|$ total variation measure). By Plancherel's theorem and that $|\hat{f}(\lambda)|$ is proportional to $1/ \sqrt{|\lambda|}$, 
    $$\int \frac{\hat{\phi}_{\delta}^2 (\lambda)} { |\hat{f} (\lambda)|^2}  d \lambda \lesssim \int \lambda e^{- \delta^2 \lambda^2} \, d\lambda \lesssim \delta^{-2}.$$
    If $U \sim F^*$ is independent of $Z \sim N(0,1)$, then $(U, U+\delta Z)$ gives a coupling of $F^*$ and $F^* * \phi_\delta$. Under this coupling $\ex |U + \delta Z - U|^k  = O(\delta^k)$ which implies $W_k\left(F^*, F^* * \phi_\delta\right)^k \leq \mathrm{E}|\delta Z|^k \lesssim \delta^k$. By combining the preceding inequalities we see that, for any $M >0$:
    \begin{align*}
        W_k\left(F^*, F^*_0 \right)^k  &\leq W_k\big(F^* , F^*_0 * \phi_\delta\big)^k +  W_k\big( F^*_0 * \phi_\delta, F^* * \phi_\delta\big)^k + W_k\big(F^* * \phi_\delta, F^*_0 \big)^k\\
        &\lesssim M^{k+1 / 2} h\left(g_{\scriptscriptstyle{F}}, g_0\right) \delta^{-2}+e^{-M}+\delta^k
    \end{align*}
    As done in \cite{25}, the lemma follows by choosing for $\varepsilon = h\left(g_{\scriptscriptstyle{F}}, g_0\right)$ the optima $M=k /(k+1/2) \log \left(C_k / \varepsilon\right) \text { and } \delta=\left(M^{k+ 2} \varepsilon\right)^{1 /(k+1/2)}$, which are eligible by assumption.
\end{proof}



 








\section{Bibliographic remarks}

There is a vast literature on Wicksell's problem, see for instance \cite{18, 20, 27, 17, 10,1,  2,16,3,19, 9, 15}. Wicksell's problem has present-day applications, where it is common to see the use of the so-called Saltykov method \cite{49}, and numerical discretization methods \cite{34}, often lacking theoretical guarantees. These methods are far from being efficient and in some cases not even consistent. The estimation of $V_0$ given data from $g_0$ is of interest not only for stereological procedures, but also from the mathematical point of view, both because of the unusual rates of convergence and because of the non-standard efficiency theory.

In the nonparametric frequentist approach to Wicksell's problem, the standard rate of convergence for pointwise asymptotic normality in estimating \( V_0 \) is \( (\log n)^{-1/2} \sqrt{n} \) (cf.\ \cite{27, 1, 15}). The first time the naive plug-in estimator was studied is in \cite{15}. Later in \cite{1} the proof of the asymptotic normality of this estimator was simplified considerably. However, viewed as a function, the naive plug-in estimator is severely ill-behaved due to infinite discontinuities and non-monotonicity, and this led to the definition of the IIE in \cite{27,1}. 
The Nonparametric Maximum Likelihood Estimator is studied in \cite{18,3}. For Wicksell's problem, the usual definition of the NPMLE is also ill-posed, as the log-likelihood can attain value infinity. However, work has been done, in \cite{18} and \cite{3}, to obtain likelihood based estimators. Confidence interval results can be found in \cite{20,19}. The results of \cite{20} provide confidence bands by developing a methodology for isonization procedures, which often lead to nuisance parameters in the limiting distributions, that tackles this issue. However, they do not cover the setting considered in the present chapter. Extending their results to our case appears non-trivial, as their proof techniques rely on at least the differentiability of the estimated object, an assumption we do not make. The bootstrap results in \cite{19} are more directly relevant and will be discussed in detail in the simulation section, including a comprehensive numerical comparison. The main advantage of the methodology proposed in this chapter, if compared to the bootstrap methods of \cite{19}, is that it relies on significantly less stringent assumptions. In particular, our results do not depend on the assumption that the underlying distribution function $F_0$ is differentiable at $x$. Instead, our methodology is provably consistent for all local smoothness levels $\gamma > 1/2$ of $V_0$, where $\gamma$ is defined in condition \eqref{eq: condition roughness} below.


In \cite{10}, the authors investigate efficient estimation in the classical Wicksell problem, focusing on the scenario where the distribution function $F^*_0$ of interest belongs to a H\"{o}lder class of smoothness with degree $\gamma > \frac{1}{2}$. They propose a kernel density-based estimator that, under appropriate bandwidth selection, achieves asymptotic efficiency. However, a notable limitation is that the estimator requires prior knowledge or guesstimation of the smoothness parameter $\gamma$, making it non-adaptive. Moreover, efficiency is only guaranteed under global smoothness assumptions, i.e., the entire distribution function must belong to a global H\"{o}lder class. Despite these limitations, their work remains the first foundational contribution to semiparametric efficiency in Wicksell's problem. \cite{27} builds upon earlier work of \cite{10} and \cite{1} by addressing both the adaptivity and locality of efficient estimation in Wicksell's problem. The key innovation consists in showing that the Isotonized Inverse Estimator (IIE) achieves asymptotic efficiency under the same local smoothness assumptions on the target function $F^*_0$ of the present chapter. Specifically, instead of assuming that $F^*_0 \in \Lambda_\gamma$ globally as in \cite{10}, the analysis in \cite{27} only requires that $F^*_0$ satisfies local H\"{o}lder-type smoothness at the evaluation point $x$ and at the origin. The estimator adapts automatically to the local regularity of $F^*_0$, offering a more flexible framework for practitioners. Additionally, \cite{27} provides a local asymptotic minimax lower bound, showing that the IIE attains the asymptotically optimal risk in this refined setting. The estimator thus achieves the same efficiency as the KDE-based approaches, but without requiring tuning of smoothness parameters, and without imposing strong global assumptions. Finally, other results concerning efficient estimation in Wicksell's problem are contained in \cite{50} and \cite{17}. \cite{50} shows that in the case the underlying cdf is constant on a \textit{known} interval of strictly positive length the IIE evaluated at $x$ is not efficient, for $x$ in the interval of constancy. However, it behaves in practice very closely to informed estimators that make use of the fact that the cdf is constant on that interval. Finally, one of the authors of \cite{10} did some further work in \cite{17} where it was shown how to attain efficient smoothing for the plug-in estimator under squared losses. 

While Bayesian methods for inverse problems are widely adopted, their theoretical guarantees remain scarce, particularly in infinite-dimensional settings. Although recent years have seen significant advances in this area \cite{62,35,37,41,36,44,38,39,40,42,26,43}, the technical frameworks developed there do not readily apply to the specific inverse problem and Dirichlet process prior setting of this chapter. We provide a detailed discussion of this related work in Supplement \cite{64}. A similar methodology to the one proposed in this chapter --- using an unconstrained conjugate prior for the distribution function of the observables and projecting the posterior samples into the space of monotone functions --- has been very recently adopted in isotonic regression settings (c.f.\ \cite{54,55,56}), isotonic density estimation (c.f.\ \cite{60}) and ODEs settings (c.f.\ \cite{57,58}). \cite{57,58} develop Bayesian two-step methods for parameter inference in ODE models by embedding them in a nonparametric regression framework with B-spline priors. The parameter posterior is obtained by minimizing a discrepancy between the estimated regression function and the ODE solution, either via derivative matching or Runge-Kutta approximations. \cite{58} further establishes a Bernstein-von Mises theorem for the Runge-Kutta-based method, showing that the resulting Bayes estimator is asymptotically efficient. However, their framework relies on a very different type of prior modelling and the projection method does not entail isotonization. In \cite{54}, the authors study a nonparametric Bayesian regression model for a response variable \( Y \) with respect to a predictor variable \( X \in [0,1] \), given by \( Y = f(X) + \varepsilon \), where \( f \) is a monotone increasing function and \( \varepsilon \) is a mean-zero random error. They propose a "projection-posterior" estimator \( f^* \), constructed using a finite random series of step functions with normal basis coefficients as a prior for \( f \). Their main finding is that when \( f^* \) is centered at the maximum likelihood estimator (MLE) and rescaled appropriately, the Bernstein-von Mises (BvM) theorem does not hold. However, when \( f^* \) is instead centered at the true function \( f_0 \), a BvM result is obtained. This phenomenon resonates with the known inconsistency of the bootstrap for the Grenander estimator in isotonic regression (c.f.\ \cite{63}) and represents a key difference with the present work. In our case, the "projection-posterior" IIP \( \hat{V}_{\scriptscriptstyle{G}} \), when centered at the IIE \( \hat{V}_n \), is shown to converge in distribution.  This may also not be surprising given the existing results in \cite{19}, where bootstrap methods are shown to be consistent for Wicksell's problem. 

Recent research contributions \cite{62,35,37,41,36,44,38,39,40,42,26,43} have introduced proof techniques aimed at establishing theoretical guarantees for inference based on posterior distributions in a variety of models. These techniques rely on appropriately chosen priors and are applicable across some contexts. \cite{62} and \cite{35} made significant advances in understanding the Bernstein--von Mises (BvM) phenomenon in nonparametric Bayesian analysis. Specifically, \cite{62} provided foundational results in the Gaussian white noise model, while \cite{35} extended these findings to more general settings, introducing innovative tools such as multiscale spaces, which naturally accommodate certain nonparametric priors and posteriors. However, their results do not cover the setting of the current chapter, as the theoretical tools there developed do not naturally accommodate DP priors in inverse settings. \cite{37} extended the BvM theorem to semiparametric models, focusing on the asymptotic behavior of posterior distributions for smooth functionals of the parameter of interest. However, their methodology is not applicable to our setting for two key reasons. First, the primary functional of interest in our work, $V_0$, which is essentially the $1/2$-integral of $g_0$, is not a smooth functional of $g_0$. This is reflected in the fact that, since their framework applies to smooth functionals, they obtain $\sqrt{n}$ convergence rates, whereas we obtain the slower $\sqrt{n} / \sqrt{\log{n}}$ rate. Second, the study of the so-called \textit{least favorable direction} in our context, as conducted in \cite{27}, does not lead to the classical Laplace expansion of the posterior (cf.\ Eq.\ (2.13) in \cite{37}). This discrepancy arises because the LAN expansion for Wicksell's model involves a local functional, which does not yield the global $\mathbb{L}_2$ contractions needed. \cite{41} provides a comprehensive treatment of quantile estimation under the DP prior; however, their results do not involve any inverse setting. \cite{36} and \cite{44} establish semiparametric and nonparametric BvM results in the proportional hazards and right-censoring models, respectively. While they develop highly useful techniques for handling DP process priors in survival analysis, these methods are difficult to generalize beyond that framework. In addition, their approaches do not employ projection-based posterior methodologies, as used in the present chapter. \cite{38} investigates the statistical inverse problem of recovering a function $f$ from general X-ray transform measurements corrupted by Gaussian noise. The proof techniques of this chapter rely on the fact that the noise is Gaussian and that the used prior is a Gaussian process, making it difficult to extend them to our setting (with beta noise and DP process prior). \cite{39} considers the inverse problem of estimating the potential in a Schr\"{o}dinger equation while the solution of the equation is corrupted by Gaussian noise, while \cite{40} considers the inverse problem of recovering the parameters governing compound Poisson processes. In both settings, the BvM results are based on ad-hoc priors constructed with wavelet basis functions, which is far from the DP modelling we use in the current chapter. \cite{42} establishes a BvM result in the Gaussian white model with priors constructed with wavelet basis functions and \cite{26} for the Average Causal Effect in a causal setting, using Gaussian process priors. The latter uses techniques closely related to \cite{37}, which cannot be applied to our setting. Finally \cite{43} also uses wavelet-based priors to derive semiparametric BvM results for linear functionals of densities. While the approach of this chapter constitutes a relevant advancement in the semiparametric BvM theory, the inverse problem treated in the current work cannot be framed in the context of \cite{43}.

One of the key advantages our approach shares with the works in \cite{57,58,54,55,56} is the computational efficiency, driven by posterior conjugacy, as well as the relative simplicity of the asymptotic analysis. However, unlike in \cite{54}, our method ensures that the limiting distribution in the Bernstein-von Mises (BvM) theorem coincides with that of the asymptotically efficient estimator --- the IIE of \cite{27}.  
